\documentclass{book}
\usepackage{amsmath, amsthm, amssymb}
\usepackage{amscd}
\usepackage[backref=section]{hyperref}
\usepackage{geometry}
\usepackage{chngcntr}
\usepackage{tikz-cd}
\usepackage{pifont}
\usepackage{mathtools}
\usepackage{mathrsfs}
\usepackage{url}
\usepackage{xcolor} % 用于还原红色的 (***)
\usepackage{stmaryrd} % 用于 \mapsfrom 符号
\usepackage{docmute}
\hypersetup{colorlinks=true, linkcolor=blue}
% Customizing backref to show only section numbers
\renewcommand*{\backref}[1]{}
\renewcommand*{\backrefalt}[4]{%
    \ifcase #1 %
    \or
        #2%
    \else
        #2%
    \fi
}
\geometry{
  a4paper,
  left=20mm,
  right=20mm,
  top=25mm,
  bottom=25mm,
  bindingoffset=0mm
}

% 定义常用算子
\DeclareMathOperator{\Spec}{Spec}
\DeclareMathOperator{\rk}{rank}
\DeclareMathOperator{\Img}{Im} % 像 Image
\DeclareMathOperator{\Sym}{Sym}     % Sym 算子
\providecommand{\C}{\mathbb{C}}
\providecommand{\T}{\mathcal{T}}
\providecommand{\A}{\mathbb{A}}
\providecommand{\End}{\operatorname{End}}
% \providecommand{\O}{\mathcal{O}}
\providecommand{\D}{\mathcal{D}}
\providecommand{\Hom}{\operatorname{Hom}}
\providecommand{\Ext}{\operatorname{Ext}}
\providecommand{\Der}{\operatorname{Der}}
\providecommand{\gr}{\operatorname{gr}}
\providecommand{\ord}{\operatorname{ord}}
\providecommand{\Mod}{\mathrm{Mod}}
\providecommand{\Z}{\mathbb{Z}}
\providecommand{\Supp}{\operatorname{Supp}}
\providecommand{\Ch}{\operatorname{Ch}}
\providecommand{\Ann}{\operatorname{Ann}}
\providecommand{\rad}{\operatorname{rad}}
\providecommand{\cO}{\mathcal{O}}
\providecommand{\Id}{\operatorname{Id}}
\providecommand{\Poisson}{\mathrm{P}}
\providecommand{\Rees}{\operatorname{R}}
\providecommand{\Diff}{\operatorname{Diff}}
\providecommand{\Ker}{\operatorname{Ker}}
\providecommand{\Tor}{\operatorname{Tor}}
\providecommand{\Lie}{\operatorname{Lie}}


\newcommand{\uM}{\mathcal{M}}

% 为了方便，定义一些快捷命令
\newcommand{\Dx}{\mathcal{D}_X}
\newcommand{\ox}{\omega_X}
\newcommand{\op}{\mathrm{op}}
\newcommand{\oxinv}{\omega_X^{-1}}
\newcommand{\oxstar}{\omega_X^*}

% 定义环境和命令
\newcommand{\FilM}[1]{\Omega_{#1} M}
\newcommand{\FilF}[1]{F_{#1} M}
\newcommand{\FilA}[1]{\Gamma_{#1} A}

% 定义切空间符号，手写体通常是 \mathcal{T}
\newcommand{\Tv}{\mathcal{T}}
% 为了美观，定义 smooth 的上标为正体
\newcommand{\sm}{\mathrm{sm}}

\newcommand{\G}{\mathbb{G}}
\newcommand{\DR}{\mathrm{DR}}
\newcommand{\Ox}{\mathcal{O}_X}
\newcommand{\Oy}{\mathcal{O}_Y}
\newcommand{\m}{\mathfrak{m}}

\counterwithout{section}{chapter}
\renewcommand{\thesection}{\arabic{section}}
\numberwithin{equation}{section}

\newtheoremstyle{thmdot}{}{}{\itshape}{}{\bfseries}{.}{ }{}
\newtheoremstyle{defdot}{}{}{\upshape}{}{\bfseries}{.}{ }{}

\theoremstyle{thmdot}
\newtheorem{theorem}{Theorem}[section]
\newtheorem{lemma}[theorem]{Lemma}
\newtheorem{proposition}[theorem]{Proposition}
\newtheorem{corollary}[theorem]{Corollary}

\theoremstyle{defdot}
\newtheorem{definition}[theorem]{Definition}
\newtheorem{example}[theorem]{Example}
\newtheorem{exercise}[theorem]{Exercise}
\newtheorem{remark}[theorem]{Remark}
\newtheorem{note}[theorem]{Note}
\newtheorem{notation}[theorem]{Notation}

\title{D-modules}
\begin{document}
\maketitle
\tableofcontents

\chapter{Basics of D-modules on smooth affine algebraic varieties}
\section{Recall of smooth affine algebraic geometry over $\C$}
Let $R$ be an integral commutative $\C$-algebra of finite type. Set
\[
  X := \Spec R,
\]
 the affine algebraic variety associated to $R$, endowed with the Zariski topology.
For an ideal $I\subseteq R$, write
\[
  V(I) := \{\, \mathfrak p\in \Spec R \mid \mathfrak p\supseteq I \,\},
\]
so that a subset $V\subseteq X$ is closed iff $V=V(I)$ for some $I$.

We denote the module of K"ahler differentials of $R$ over $\C$ by
\[
  \Omega_R = \Omega_X = \Omega_{R/\C}.
\]
\begin{definition}[K\"ahler differentials]\label{def:kahler}
The $R$-module $\Omega_{R/\C}$ is generated by formal symbols $\mathrm df$ with $f\in R$,
subject to the relations, for all $f,g\in R$ and $a,b\in\C$,
\begin{align*}
  \mathrm d(fg) &= f\,\mathrm dg + g\,\mathrm df \qquad\text{(Leibniz rule)},\\
  \mathrm d(af+bg) &= a\,\mathrm df + b\,\mathrm dg.
\end{align*}
\end{definition}
\begin{exercise}
Show that $\mathrm dc=0$ for every constant $c\in\C$.
\end{exercise}
\begin{proof}[Solution to the exercise]
Using linearity and the Leibniz rule,
\(\mathrm d(1)=\mathrm d(1\cdot 1)=1\,\mathrm d1+1\,\mathrm d1=2\,\mathrm d1\), hence
\(\mathrm d1=0\). For $c\in\C$ we have
\(\mathrm dc=\mathrm d(c\cdot 1)=c\,\mathrm d1+1\,\mathrm dc=\mathrm dc\), and the only possibility is
\(\mathrm dc=0\).
\end{proof}

Define the $R$-module
\[
  \T_X := \Hom_R\!\big(\Omega_{R/\C},\, R\big) \cong \Der_{\C}(R,R),
\]
called the module of $\C$-derivations on $X$ (the vector fields on $X$).

\begin{definition}[Smoothness]
We say that $R$ (or $X=\Spec R$) is \emph{smooth} if $\Omega_{R/\C}$ is a free $R$-module of
finite rank. In this case, the rank
\[
  n := \operatorname{rank}_R(\Omega_{R/\C})
\]
which is equal to the topological dimension of $X$.

\begin{proof}[Prove that $\mathrm{rank}_R(\Omega_{R/\mathbb{C}}) = \dim X$]
Let $K = \mathrm{Frac}(R)$ denote the fraction field of $R$. Since $R$ is integral.

First, recall a fundamental result from field theory \cite[Chapter 2, Proposition 8.6A.]{hartshorne1977algebraic}: for a field extension $K/\mathbb{C}$, the module of Kähler differentials $\Omega_{K/\mathbb{C}}$ is a vector space over $K$, and its dimension is equal to the transcendence degree of the extension:
\[
    \dim_K (\Omega_{K/\mathbb{C}}) = \mathrm{tr.deg}_{\mathbb{C}}(K).
\]

One should note that $\mathbb{C}$ is of characteristic 0, which ensures all extensions are separable, even when extensions are not algebraic.

By the definition of the dimension of an algebraic variety \cite[Chapter 1, Proposition 1.8A.(a)]{hartshorne1977algebraic}, we have:
\[
    \dim X = \mathrm{tr.deg}_{\mathbb{C}}(K).
\]

Next, we relate the differentials on $R$ to those on $K$ via localization. We have the canonical isomorphism Chapter 2, Proposition 8.2A. \cite[Chapter 2, Proposition 8.2A.]{hartshorne1977algebraic}:
\[
    \Omega_{K/\mathbb{C}} \cong \Omega_{R/\mathbb{C}} \otimes_R K.
\]
The hypothesis states that $\Omega_{R/\mathbb{C}}$ is a \textbf{free $R$-module} of rank $n$. A key property of free modules is that their rank is preserved under base change to the fraction field. Therefore:
\[
    \dim_K (\Omega_{R/\mathbb{C}} \otimes_R K) = \mathrm{rank}_R (\Omega_{R/\mathbb{C}}) = n.
\]

Combining these equalities, we conclude:
\[
    n = \dim_K (\Omega_{R/\mathbb{C}} \otimes_R K) = \dim_K (\Omega_{K/\mathbb{C}}) = \mathrm{tr.deg}_{\mathbb{C}}(K) = \dim X.
\]  

\end{proof}

\end{definition}

If $X$ is smooth, a tuple $(x_1,\dots,x_n)$ with $x_i\in R$ is called an \emph{algebraic
coordinate system} if
\[
  \Omega_{R/\C} \cong \bigoplus_{i=1}^n R\,\mathrm d x_i.
\]
Equivalently, the dual basis gives
\[
  \T_X \cong \bigoplus_{i=1}^n R\,\partial_{x_i},
\]
where $\partial_{x_i}$ are the vector fields dual to $\mathrm d x_i$.

\begin{example}[Polynomial ring]
Let $R=\C[x_1,\dots,x_n]$. Then
\[
  \Omega_{R/\C} \cong \bigoplus_{i=1}^n R\,\mathrm d x_i,
\]
so $X=\Spec R\simeq \A^n_{\C}$ is smooth and $(x_1,\dots,x_n)$ is an algebraic coordinate system.
\end{example}

From now on assume $X=\Spec R$ is smooth affine and $(x_1,\dots,x_n)$ are algebraic coordinates.
The symmetric algebra of vector fields is
\[
  \Sym^{\bullet}\T_X = \bigoplus_{m\ge 0}\Sym^{m}\T_X \simeq R[\xi_1,\dots,\xi_n],
\]
canonically an $R$-algebra (identify $\xi_i$ with the class of $\partial_{x_i}$). Define the
cotangent bundle
\[
  T^{\ast}X := \Spec\, \Sym^{\bullet}\T_X \xrightarrow{\quad\pi\quad} X.
\]
and $\pi$ is induced by the $R$-algebra map $R = \Sym^{0}\T_X \subset \bigoplus_{m\ge 0}\Sym^{m}\T_X$.


\section{D-modules on smooth affine varieties}
Let $\End_{\C}(R)$ denote the $\C$-algebra of all $\C$-linear maps $R\to R$. There are natural
embeddings
\[
  R\hookrightarrow \End_{\C}(R),\quad f\longmapsto (r\mapsto fr),
\]
and
\[
  \T_X\hookrightarrow \End_{\C}(R),\quad v\longmapsto v':=[f\mapsto v(\mathrm d f)],
  \begin{tikzcd}[row sep=3em, column sep=1.5em]
      R \arrow[rr, "d"] \arrow[dr, "v'"]  % v' 在箭头上方（默认）
      & & \Omega_R \arrow[dl, "v"]       % v 在箭头下方（不使用 ' 换边）
      \\
      & R &
  \end{tikzcd}.
\]

\begin{definition}[Ring of differential operators]
The ring $\D_X$ of algebraic differential operators on $X$ is the $\C$-subalgebra of
$\End_{\C}(R)$ generated by $R$ and $\T_X$.
\end{definition}

\begin{example}[Weyl algebra]
If $X=\A^n_{\C}=\Spec \C[x_1,\dots,x_n]$, then
\[
  \D_X \cong \C[x_1,\dots,x_n]\langle \partial_{x_1},\dots,\partial_{x_n}\rangle,
\]
with relations
\[
  [\partial_{x_i},x_j]=\delta_{ij},\quad [x_i,x_j]=0,\quad [\partial_{x_i},\partial_{x_j}]=0.
\]
\end{example}

Thus $\D_X$ is non-commutative. Every operator $P\in\D_X$ can be written in multi-index notation as
\[
  P=\sum_{\underline{\alpha}\in\mathbb N^n} f_{\underline{\alpha}}\,\partial^{\underline{\alpha}},\qquad
  \partial^{\underline{\alpha}}:=\partial_{x_1}^{\alpha_1}\cdots\partial_{x_n}^{\alpha_n},\quad f_{\underline{\alpha}}\in R.
\]

\paragraph{Basic modules.}\mbox{}\par

\noindent 1.  The coordinate ring $R$ is a left $\D_X$-module where$\partial_{x_i}\cdot r = \partial_{x_i}(\mathrm d r)\in R$.

\noindent Equivalently, writing $\mathrm dr=\sum_i (\partial_{x_i}r)\,\mathrm d x_i$, we have $\partial_{x_i}r$ the
$i$-th coefficient.\\

\par\noindent 2.  The top exterior power
\[
  \omega_X := \bigwedge{}^{n}\Omega_{R/\C} \cong R\,\mathrm d x_1\wedge\cdots\wedge\mathrm d x_n
\]
is the \textbf{dualizing module} of $R$. It carries a natural right $\D_X$-module structure via the transpose
anti-involution $t\colon \D_X\to\D_X$ determined by
\[
  \prescript{t}{}{x_i}=x_i,\qquad \prescript{t}{}{\partial_{x_i}}=-\partial_{x_i}.
\]
If $P=\sum_{\underline{\alpha}} f_{\underline{\alpha}}\,\partial^{\underline{\alpha}}$, then
\[
  \prescript{t}{}{P}=\sum_{\underline{\alpha}} (-\partial)^{\underline{\alpha}} f_{\underline{\alpha}},
\]
and the action is
\[
  (r\,\mathrm d x_1\wedge\cdots\wedge\mathrm d x_n)\cdot P := \left(\prescript{t}{}{P}\cdot r\right)\,\mathrm d x_1\wedge\cdots\wedge\mathrm d x_n.
\]

\subsection*{Side change and opposite ring}
Let $A$ be a ring and let $A^{\mathrm{op}}$ denote the opposite ring. i.e. $A^{\op} = \{ a^{\op} \mid a \in A \}$, $a^{\op} \cdot b^{\op} = (ba)^{\op}$. Then right $A$-modules are the
same as left $A^{\mathrm{op}}$-modules.

\begin{lemma}\label{lem:DXop}
There is a natural isomorphism of rings
\[
  \D_X^{\mathrm{op}} \simeq \omega_X \otimes_R \D_X \otimes_R \omega_X^{\!*},\qquad
\text{where } \omega_X^{\!*}:=\Hom_R(\omega_X,R) \cong R\,(\mathrm d x_1\wedge\cdots\wedge\mathrm d x_n)^{-1}.
\]
\end{lemma}
\begin{proof}
Define
\[
  \Phi(P^{\mathrm{op}}):=\mathrm d x_1\wedge\cdots\wedge\mathrm d x_n\;\otimes\; \prescript{t}{}{P}\;\otimes\;
  (\mathrm d x_1\wedge\cdots\wedge\mathrm d x_n)^{-1}.
\]
Then for $P,Q\in\D_X$ we have
\begin{align*}
  \Phi(P^{\mathrm{op}}Q^{\mathrm{op}})
  &= \Phi((QP)^{\mathrm{op}})\\
  &= \mathrm d x\;\otimes \prescript{t}{}{(QP)}\otimes (\mathrm d x)^{-1}\\
  &= \mathrm d x\;\otimes \prescript{t}{}{P}\prescript{t}{}{Q}\otimes (\mathrm d x)^{-1}\\
  &=: \left(\mathrm d x\;\otimes \prescript{t}{}{P}\otimes (\mathrm d x)^{-1}\right)
    \cdot \left(\mathrm d x\;\otimes \prescript{t}{}{Q}\otimes (\mathrm d x)^{-1}\right)\\
  &= \Phi(P^{\mathrm{op}})\,\Phi(Q^{\mathrm{op}}),
\end{align*}

so $\Phi$ is a ring homomorphism. 
\end{proof}

\begin{remark}
  $\D_X$ is a $(R,R)$-bimodule in the natural way. One should note that the left and the right R-module actions on $\D_X$ are NOT the same.
\end{remark}

Write $\Mod(\D_X^{\mathrm{op}})$ for the category of right $\D_X$-modules and $\Mod(\D_X)$ for the
category of left $\D_X$-modules.

\begin{theorem}\label{thm:left-right-equivalence}
Tensoring with the canonical bundle yields an equivalence of categories
\[
  \Mod(\D_X) \xrightarrow{\quad\sim\quad} \Mod(\D_X^{\mathrm{op}}),\qquad M\longmapsto \omega_X\otimes_R M,
\]
with quasi-inverse $N\longmapsto \omega_X^{\!*}\otimes_R N$.
\end{theorem}
\begin{proof}
    % 第一部分
    \noindent 1. Let $M \in \Mod(\Dx)$.
    \begin{align*}
        \ox \otimes_R \Dx \otimes_R \oxinv \cdot (\ox \otimes_R M) 
        &= \ox \otimes_R (\Dx \cdot M) \\
        &\subseteq \ox \otimes_R M
    \end{align*}
    \[
        \implies \ox \otimes_R M \in \Mod(\Dx^{\mathrm{op}})
    \]

    % 第二部分
    \noindent 2. Let $N \in \Mod(\Dx^{\mathrm{op}}) = \Mod(\ox \otimes_R \Dx \otimes_R \oxstar)$.
    \begin{align*}
        \Dx \cdot (\oxstar \otimes_R N) 
        &= (\oxstar \otimes_R \ox \otimes_R \Dx \otimes_R \oxstar) \cdot N \\
        &= \oxstar \otimes_R (\ox \otimes_R \Dx \otimes_R \oxstar \cdot N) \\
        &\subseteq \oxstar \otimes_R N
    \end{align*}
    \[
        \implies \oxstar \otimes_R N \in \Mod(\Dx) 
    \]
\end{proof}

In particular, we will henceforth only consider left $\D_X$-modules, written simply as $\D_X$-modules.

\begin{remark}
  For a comprehensive treatment of side change, see \cite[Lemma 1.2.7, Proposition 1.2.9, Proposition 1.2.12]{hotta2008d}.
\end{remark}
\chapter{Filtered rings and modules}
\section{Filtered rings I}
\begin{definition}[$\mathbb Z$-filtration]\label{def:z_filtration}
    Let $A$ be a unital ring. \(\Gamma_\bullet A\) is a family \(\{\Gamma_i A\}_{i\in\mathbb Z}\) called a \textbf{$\mathbb Z$-filtration} of $A$ if
    \begin{itemize}
        \item Each \(\Gamma_i A\) is an additive subgroup of \(A\);
        \item $\Gamma_i A \subseteq A$;
        \item $\Gamma_i A \subseteq \Gamma_{i+1} A$ and $1\in \Gamma_0 A$;
        \item $\bigcup_{i\in\mathbb Z} \Gamma_i A = A$;
        \item $\Gamma_i A\cdot \Gamma_j A \subseteq \Gamma_{i+j} A$ for all $i,j$; and often we assume the
              filtration is separated: $\bigcap_{i\in\mathbb Z}\Gamma_i A=0$.
    \end{itemize}
\end{definition}
We say $A$ is \textbf{$\mathbb Z$-filtered} if $\Gamma_\bullet A$ is$\mathbb Z_{\ge 0}$-filtered if additionally $\Gamma_i A=0$ for all $i<0$.

The \textbf{associated graded ring} is
\[
    \operatorname{gr}_\bullet A = \bigoplus_{i\in\mathbb Z} \Gamma_i A/\Gamma_{i-1} A.
\]
The Rees ring (with respect to $\Gamma$) is
\[
    R^{\Gamma}_\bullet(A) = \bigoplus_{i\in\mathbb Z} \Gamma_i A\,u^i \subseteq A[u,u^{-1}],
\]
which records the filtration multiplicatively.

\subsection*{Basic consequences}
\begin{lemma}\label{lem:rees-quotients}
    For the Rees ring $R^{\Gamma}_\bullet (A)$ one has natural isomorphisms
    \[
        \gr_\bullet (A) \simeq R^{\Gamma}_\bullet (A)/u\,R^{\Gamma}_\bullet (A),\qquad
        A \simeq R^{\Gamma}_\bullet (A)/(u-1)\,R^{\Gamma}_\bullet (A).
    \]
\end{lemma}

\begin{proposition}\label{prop:noeth}
    Let $A$ be a filtered ring with filtration $\Gamma$.
    \begin{enumerate}
        \item If $R^{\Gamma}_\bullet (A)$ is Noetherian, then both $A$ and $\gr_\bullet (A)$ are Noetherian.
        \item If the filtration is nonnegative (i.e. $\Gamma_iA=0$ for $i<0$) and
              $\gr_\bullet (A)$ is Noetherian, then $A$ is Noetherian.
    \end{enumerate}
\end{proposition}
\begin{proof}
    (1) By Lemma~\ref{lem:rees-quotients}, both $\operatorname{gr}_\bullet (A)$ and $A$ are surjective homomorphic images of the Rees ring $R^{\Gamma}_\bullet (A)$. Since the homomorphic image of a Noetherian ring is Noetherian, if $R^{\Gamma}_\bullet (A)$ is Noetherian, then so are its quotients $\operatorname{gr}_\bullet (A)$ and $A$.

    (2) Assume the filtration is nonnegative, i.e., $\Gamma_i A = 0$ for $i < 0$, which implies $A = \bigcup_{p \ge 0} \Gamma_p A$.
    Let $I_1 \subseteq I_2 \subseteq \cdots$ be an ascending chain of left ideals of $A$.
    Consider the associated graded ideals in $\operatorname{gr}_\bullet (A)$ defined by
    \[
        \mathfrak{a}_j := \operatorname{gr}_\bullet(I_j) = \bigoplus_{p \ge 0} \frac{I_j \cap \Gamma_p A + \Gamma_{p-1} A}{\Gamma_{p-1} A}.
    \]
    Since $I_j \subseteq I_{j+1}$, we have an ascending chain $\mathfrak{a}_1 \subseteq \mathfrak{a}_2 \subseteq \cdots$ in $\operatorname{gr}_\bullet (A)$.
    Because $\operatorname{gr}_\bullet (A)$ is Noetherian, this chain stabilizes. Choose $N$ such that $\operatorname{gr}_\bullet(I_j) = \operatorname{gr}_\bullet(I_N)$ for all $j \ge N$.

    We claim that $I_j = I_N$ for all $j \ge N$, which implies $A$ is Noetherian.
    Fix $j \ge N$. Since $I_N \subseteq I_j$, it suffices to show equality. We proceed by induction on the filtration degree $p$ that
    \[
        I_j \cap \Gamma_p A = I_N \cap \Gamma_p A.
    \]
    For $p < 0$, both sides are $0$. Assume the equality holds for $p-1$.
    Since $\operatorname{gr}_\bullet(I_j) = \operatorname{gr}_\bullet(I_N)$, their degree-$p$ components are identical:
    \[
        \frac{I_j \cap \Gamma_p A + \Gamma_{p-1} A}{\Gamma_{p-1} A} = \frac{I_N \cap \Gamma_p A + \Gamma_{p-1} A}{\Gamma_{p-1} A}.
    \]
    This equality implies
    \[
        I_j \cap \Gamma_p A + \Gamma_{p-1} A = I_N \cap \Gamma_p A + \Gamma_{p-1} A.
    \]
    Intersecting both sides with $I_j \cap \Gamma_p A$ and using the modular law ($I_N \subseteq I_j$, so the modular law applies), we get:
    \[
        I_j \cap \Gamma_p A = I_N \cap \Gamma_p A + (I_j \cap \Gamma_{p-1} A).
    \]
    By the induction hypothesis, $I_j \cap \Gamma_{p-1} A = I_N \cap \Gamma_{p-1} A \subseteq I_N \cap \Gamma_p A$. Thus,
    \[
        I_j \cap \Gamma_p A = I_N \cap \Gamma_p A.
    \]
    Since $A = \bigcup_{p} \Gamma_p A$, taking the union over all $p$ yields $I_j = I_N$.
    Thus, every ascending chain stabilizes, and $A$ is Noetherian.
\end{proof}

\begin{remark}
    It is not recommended to use ``s.t.'' to replace ``such that'' in the formal writing,
    since ``s.t.'' may represent ``subject to'' in some contexts.
\end{remark}

\subsection*{Filtered modules and good filtrations}
Let $A$ be a $\mathbb{Z}$-filtered ring with filtration $\Gamma_\bullet A$ and let $M$ be a left
$A$-module.

\begin{definition}[Compatible and Good Filtration]\label{def:filtered-module}
    Let $\Omega_\bullet M$ be a $\mathbb{Z}$-filtration on $M$ (also, each $\Omega_i M$ is an additive subgroup of $M$) such that $\bigcup_k \Omega_k M = M$.

    \begin{enumerate}
        \item $\Omega_\bullet M$ is called a \textbf{compatible filtration} over $\Gamma_\bullet A$ if:
              \[
                  \Gamma_k A \cdot \Omega_i M \subseteq \Omega_{k+i} M.
              \]

        \item $\Omega_\bullet M$ is called a \textbf{good filtration} over $\Gamma_\bullet A$ if it is compatible and the Rees module:
              \[
                  R_\bullet(M) = R_\bullet^\Omega(M) = \bigoplus_i \Omega_i M \cdot u^i
              \]
              is finitely generated (f.g.) over $R(A)$.

              This implies that the associated graded module:
              \[
                  \gr_\bullet M = \gr_\bullet^\Omega M = \bigoplus_i \frac{\Omega_i M}{\Omega_{i-1} M}
              \]
              is finitely generated (f.g.) over $\gr_\bullet A$.
    \end{enumerate}
\end{definition}

\begin{lemma}\label{lem:good-filtration-exists}
    \mbox{}
    \begin{enumerate}
        \item If $M$ is finitely generated over $A$, then good filtrations on $M$ exist.
        \item If $\Omega_\bullet M$ is good, then there exist integers $k_1,\dots,k_\ell$ and
              elements $m_1,\dots,m_\ell\in M$ with $m_i\in\Omega_{k_i}M$ such that, for every $p$,
              \[
                  \Omega_pM \,=\, \sum_{i=1}^{\ell} \Gamma_{p-k_i}A\cdot m_i.
              \]
    \end{enumerate}
\end{lemma}
\begin{proof}
    \textbf{(1)} Assume $M$ is generated by $m_1, \dots, m_\ell$ over $A$.
    Then define the filtration by:
    \[
        \FilM{p} = \sum_{i=1}^{\ell} \FilA{p} \cdot m_i.
    \]
    This defines a good filtration on $M$.

    \vspace{1em}

    \textbf{(2)} By definition, if $\FilM{\bullet}$ is a good filtration, then the Rees module $\Rees^\Omega(M)$ is finitely generated (f.g.) over the Rees ring $\Rees(A)$.

    This implies we can find homogeneous generators:
    \[
        m_1 u^{k_1}, \dots, m_\ell u^{k_\ell}
    \]
    where $m_i \in \FilM{k_i}$ (considered as degree $k_i$ elements).

    Expressing the degree $p$ part of the Rees module using these generators:
    \[
        \FilM{p} \cdot u^p = \sum_{i=1}^{\ell} \left( \FilA{p-k_i} \cdot u^{p-k_i} \right) \cdot (m_i \cdot u^{k_i}).
    \]
    By canceling $u^p$ on both sides (or comparing coefficients), we obtain the desired formula:
    \[
        \FilM{p} = \sum_{i=1}^{\ell} \FilA{p-k_i} \cdot m_i.
    \]

\end{proof}

\begin{corollary}\label{cor:compare-good-filtrations}
    Let $M$ be a finitely generated $A$-module and let $\Omega_\bullet M$ and $F_\bullet M$
    be two compatible filtrations on $M$.
    If $F_\bullet M$ is good, then there exists an integer $a$ such that
    \[
        F_pM \subseteq \Omega_{p+a}M \qquad \text{for all } p.
    \]
    Consequently, if $\Omega_\bullet M$ is also good, then there exists an integer $a$ with
    \[
        \Omega_{p-a}M \subseteq F_pM \subseteq \Omega_{p+a}M \qquad \text{for all } p.
    \]
\end{corollary}
\begin{proof}
    By Lemma~\ref{lem:good-filtration-exists}(2), since $\FilF{\bullet}$ is a good filtration, for all $p$ we have:
    \[
        \FilF{p} = \sum_{i=1}^{\ell} \FilA{p-k_i} \cdot m_i
    \]
    for some generators $m_1, \dots, m_\ell$ and integers $k_1, \dots, k_\ell$.

    Assume that $m_1, \dots, m_\ell \in \FilM{q}$ for some sufficiently large integer $q \gg 0$.

    Then we have the following inclusion chain:
    \begin{align*}
        \FilF{p} & \subseteq \sum_{i=1}^{\ell} \FilA{p-k_i} \cdot \FilM{q}                                                                    \\
                 & \subseteq \sum_{i=1}^{\ell} \FilM{p+q-k_i} \quad (\text{since } \FilM{\bullet} \text{ is compatible with } \FilA{\bullet}) \\
                 & \subseteq \FilM{p+a}
    \end{align*}
    where we choose $a \ge \max_i(q - k_i)$.
\end{proof}


\paragraph{Grothendieck group.}
For any ring $A$, let $K_0(A)$ denote the Grothendieck group of the abelian
category $\Mod(A)$: it is generated by symbols $[M]$ for $M\in\Mod(A)$ with the
relations $[M]=[M_1]+[M_2]$ for every short exact sequence
$0\to M_1\to M\to M_2\to 0$.

\begin{remark}[Well-definedness of $K_0$]
    Strictly speaking, defining the Grothendieck group on the category of \textit{all} modules $\mathrm{Mod}(A)$ is problematic for two reasons:

    \begin{enumerate}
        \item \textbf{Set-theoretic Size Issue}: The collection of isomorphism classes in $\mathrm{Mod}(A)$ forms a \textbf{proper class}, not a set. One cannot strictly generate a free abelian group from a proper class in ZFC.

        \item \textbf{Eilenberg Swindle}: Allowing infinite direct sums makes the group trivial. For any module $M$, let $M^\infty = \bigoplus_{i=1}^\infty M$. We have the isomorphism:
              \[
                  M \oplus M^\infty \cong M^\infty.
              \]
              In the Grothendieck group, this implies $[M] + [M^\infty] = [M^\infty]$, and by cancellation, $[M] = 0$.
    \end{enumerate}

    \textbf{Correction}: To avoid these issues, $K$-theory is standardly defined on "small" subcategories of \textbf{finitely generated} modules (where infinite sums are not allowed and isomorphism classes form a set):
    \begin{itemize}
        \item $K_0(A)$ is defined on $\mathbf{Proj}(A)$ (f.g. projective modules).
        \item $G_0(A)$ is defined on $\mathbf{mod}(A)$ (f.g. modules, typically for Noetherian rings).
    \end{itemize}
\end{remark}

\chapter{Characteristic varieties$\&$cycles of D-modules}
\section{Characteristic Varieties and Cycles of $\D_X$-modules}
Let $X=\Spec R$ be smooth affine with algebraic coordinates $(x_1,\dots,x_n)$.
Define an \textbf{order function} on $\D_X$ by
\[
  \ord(f)=0\ (f\in R), \qquad \ord(v)=1\ (v\in \T_X),
\]
and for a differential operator $P=\sum_{\underline{\alpha}} f_{\underline{\alpha}}\,\partial^{\underline{\alpha}}$ set
\[
  \ord(P)=\max\{\,|\underline{\alpha}|\mid f_{\underline{\alpha}}\ne 0\,\},\qquad |\underline{\alpha}|=\alpha_1+\cdots+\alpha_n.
\]
This gives an increasing $\mathbb Z_{\ge0}$-filtration $F_p\D_X:=\{\,P\in\D_X\mid \ord(P)\le p\,\}$.

\begin{remark}\label{rmk:order-filtration-is-well-defined}
  One should verify that this filtration is well-defined, i.e., independent of the choice of algebraic coordinates.
  This can be achieved by considering an equivalent definition found in \cite[Definition 3.1.]{chen2024dmodules}, \cite[Definition 2.2.]{mustata2023dmodules}, and \cite[Exercise 1.1.4.]{hotta2008d}. 
  By showing that these definitions are equivalent, it follows that our definition of the order filtration is coordinate-independent.
\end{remark}



\begin{lemma}\label{lem:grDx}
There are natural isomorphisms of $R$-algebras
\[
  \gr \D_X \cong R[\xi_1,\dots,\xi_n] \cong \Sym\,\T_X.
\]
Consequently, $\Spec(\gr_\bullet \D_X)\simeq T^{\ast}X$.
\end{lemma}
\begin{proof}
For a coordinate-free proof, see \cite[Theorem 5.2 (PBW for $\D_X$)]{chen2024dmodules}. We prove this using a local coordinate system $x_1, \dots, x_n$ with derivations $\partial_1, \dots, \partial_n$.

Consider the homomorphism of graded $R$-algebras:
\[
    \varphi: R[\xi_1, \dots, \xi_n] \longrightarrow \gr \D_X
\]
defined by mapping variables $\xi_i \mapsto \sigma_1(\partial_i) = \partial_i \pmod{F_0 \D_X}$.

Here $\sigma_i: F_i D_X \longrightarrow \frac{F_i D_X}{F_{i-1} D_X}$ denotes the principal symbol of the order-$i$ differential operator.

\textbf{Surjectivity:} Since $\gr \D_X$ is generated as an $R$-algebra by the classes of order-1 operators $\sigma_1(\partial_1), \dots, \sigma_1(\partial_n)$, the map $\varphi$ is surjective.
(Since $D_X$ is generated as an $\mathbb{C}$-algebra by $R, \partial_1, \dots, \partial_n$).

\textbf{Injectivity:} Since $\varphi$ is a graded homomorphism, it suffices to verify injectivity for homogeneous elements.
Let $f(\xi) \in R[\xi_1, \dots, \xi_n]$ be a non-zero homogeneous polynomial of degree $m$. We can write it as:
\[
    f(\xi) = \sum_{|\alpha|=m} a_\alpha(x) \xi^\alpha, \quad \text{with coefficients } a_\alpha(x) \in R.
\]
Let $P \in F_m \D_X$ be the differential operator obtained by replacing $\xi^\alpha$ with $\partial^\alpha$:
\[
    P = \sum_{|\alpha|=m} a_\alpha(x) \partial^\alpha.
\]
By definition, $\varphi(f)$ is the principal symbol $\sigma_m(P) \in F_m \D_X / F_{m-1} \D_X$.
Since by Remark~\ref{rmk:order-filtration-is-well-defined}, the order filtration is well-defined, we have $\ord(P) = m$, so $P\not\in F_{m-1} \D_X$.
Therefore, its class $\sigma_m(P) \neq 0$ in $\gr \D_X$. This proves $\varphi$ is injective.

The second isomorphism $\gr_\bullet \D_X \cong \Sym^{\bullet} \T_X$ follows from the coordinate-free identification of the principal symbol with tangent vectors ($\sigma_1(\partial) \leftrightarrow \partial \in \Der(R)$) and the universal property of symmetric $R$-algebras. Taking spectra yields $\Spec(\gr \D_X) \simeq T^* X$.
\end{proof}

\begin{corollary}\label{cor:Dx-noeth}
$\D_X$ is Noetherian.
\end{corollary}
\begin{proof}
Use Lemma~\ref{lem:grDx} and Proposition~\ref{prop:noeth}.
\end{proof}

Let $M$ be a finitely generated $\D_X$-module.

\begin{lemma}\label{lem:gr-class-independent}
If $F_\bullet M$ and $G_\bullet M$ are two good filtrations on $M$, then
\[
  [\,\gr^{F}_\bullet M\,] = [\,\gr^{G}_\bullet M\,] \quad \text{in } K_0(\gr_\bullet\D_X).
\]
\end{lemma}
\begin{proof}
By Corollary~\ref{cor:compare-good-filtrations} there exists $a\ge0$ such that
$G_{p-a}M\subseteq F_pM\subseteq G_{p+a}M$ for all $p$. We first consider the adjacent
case: suppose $F_{p-1}M\subseteq G_{p-1}M\subseteq F_pM\subseteq G_pM$ for all
$p$. Then we have short exact sequences
\[
  0\to G_{p-1}M/F_{p-1}M\to F_pM/F_{p-1}M\to F_pM/G_{p-1}M\to 0,
\]
\[
  0\to F_pM/G_{p-1}M\to G_pM/G_{p-1}M\to G_pM/F_pM\to 0.
\]
Taking classes in $K_0(\gr_\bullet \D_X)$ and summing over $p$ yields
$[\,\gr^{F}_\bullet M]=[\,\gr^{G}M]$. 

\noindent In general, $\forall k \in \mathbb Z$ set
\[
    F^{(k)}_i M = F_i M + G_{i+k} M \quad \forall i
\]
Then
\[
    \begin{cases}
        F^{(k)}_\bullet = F_\bullet \quad k \ll 0 \\
        F^{(k)}_\bullet = G_{\bullet + k} \quad k \gg 0 \\
        F^{(k)}_\bullet \text{ and } F^{(k+1)}_\bullet \text{ are adjacent}.
    \end{cases}
\]

\vspace{1em} % 增加一点垂直间距

\noindent By induction, we have
\[
    [\operatorname{gr}^F_\bullet M] = [\operatorname{gr}^G_{\bullet+k} M] = [\operatorname{gr}^G_\bullet M].
\]
\end{proof}

\paragraph{Characteristic cycle and variety.}
Given a finitely generated $\D_X$-module $M$ with a good filtration $F_\bullet M$,
the associated graded $\gr^{F}_\bullet M$ is a finitely generated module over
$\gr_\bullet \D_X\cong R[\xi_1,\dots,\xi_n]$. Its support is a closed subset of
$T^{\ast}X\simeq\Spec(\gr_\bullet \D_X)$. \\
Let {\it $\mathfrak p$} be the minimal primes in
$\operatorname{Supp}(\gr^{F}_\bullet M)$. Define the \textbf{multiplicity} $\ell(\mathfrak p)$ to be the length of
the localization $(\gr^{F}_\bullet M)_{\mathfrak p}$ as a module over $(\gr_\bullet \D_X)_{\mathfrak p}$. The
\textbf{characteristic cycle}
and \textbf{characteristic variety}
of $M$ are
\begin{align*}
  \mathrm{CC}(M) &:= \sum_{\substack{\mathfrak p\in\operatorname{Supp}(\gr^{F}_\bullet M)\\ \text{minimal}}}
  \ell(\mathfrak p)\,\overline{\{\mathfrak p\}}, \\
  \mathrm{Ch}(M) &:= \operatorname{Supp}\,\mathrm{CC}(M)
  = \bigcup_{\mathfrak p\,\text{minimal in }\operatorname{Supp}(\gr^{F}_\bullet M)} \overline{\{\mathfrak p\}}
  \subseteq T^{\ast}X.
\end{align*}

By Lemma~\ref{lem:gr-class-independent}, $\mathrm{CC}(M)$ and $\mathrm{Ch}(M)$ do not depend on the
choice of good filtration.

\section{Early properties of characteristic varieties}
Last week we constructed characteristic varieties/cycles of $\D_X$-modules. We now discuss their
basic properties. Throughout this section:
\begin{itemize}
  \item $X$ is a smooth affine variety over $\C$ with algebraic coordinates $(x_1,\dots,x_n)$,
  \item $\D_X$ denotes the algebraic differential operators on $X$,
  \item filtrations $F_\bullet$ are increasing and indexed by $\Z$.
\end{itemize}

\begin{proposition}\label{prop:charvar-union}
Let $0\to \mathcal M_1\to \mathcal M\to \mathcal M_2\to 0$ be a short exact sequence of finitely
generated $\D_X$-modules. Then $\Ch(\mathcal M)=\Ch(\mathcal M_1)\cup \Ch(\mathcal M_2)$.
\end{proposition}
\begin{proof}
Take a good filtration $F_\bullet\mathcal M$ and define $F_\bullet\mathcal M_1:=F_\bullet\mathcal M
\cap \mathcal M_1$ and $F_\bullet\mathcal M_2:=F_\bullet\mathcal M/F_\bullet\mathcal M_1$. These
are good filtrations on $\mathcal M_1$ and $\mathcal M_2$. Passing to the associated graded gives a
short exact sequence 
\[
 0\to \gr \mathcal M_1\to \gr \mathcal M\to \gr \mathcal M_2\to 0,
\]
hence
$\Supp(\gr \mathcal M)=\Supp(\gr \mathcal M_1)\cup \Supp(\gr \mathcal M_2)$, which implies the
claim.
\end{proof}

Before moving on we recall some terminology. Let $N$ be a finitely generated $R$-module. We say that
$N$ is \textbf{locally free} if $N_{\mathfrak p}$ is free over $R_{\mathfrak p}$ for every
$\mathfrak p\in\Spec R$ (equivalently $N_{\mathfrak m}$ is free over $R_{\mathfrak m}$ for every
maximal $\mathfrak m$). Given a $\D_X$-module $\mathcal M$, we call it an \textbf{integrable
connection} if $\mathcal M$ is locally free as an $R$-module.

\begin{remark}
  One should note that $R_{\mathfrak{p}}$ is a further localization of $R_{\mathfrak{m}}$ for some maximal ideal $\mathfrak{m}$ containing $\mathfrak{p}$, i.e.
  $R_{\mathfrak{p}} = (R_{\mathfrak{m}})_{\mathfrak{p}R_{\mathfrak{m}}}$.
\end{remark}

\begin{proposition}\label{prop:integrable-connection-characterizations}
    Let $\mathcal{M}$ be a finitely generated $\D_X$-module. The following are equivalent:
    \begin{enumerate}
        \item $\mathcal{M}$ is an integrable connection.
        \item $\mathcal{M}$ is finitely generated over $R$.
        \item $\Ch(\mathcal{M}) = T_X^* X$, the zero section in $T^* X$.
    \end{enumerate}
\end{proposition}


\begin{proof}
    \textbf{(1) $\implies$ (2):} is obvious.

    \textbf{(2) $\implies$ (1):} Fix an arbitrary maximal ideal $\mathfrak{m} \subseteq R$.
    By changing algebraic coordinates, we can assume $\mathfrak{m} = (x_1, \dots, x_n)$.

    $M_{\mathfrak{m}}$ is a $\mathcal{D}_{X,\mathfrak{m}}$-module, where $\mathcal{D}_{X,\mathfrak{m}} = R_{\mathfrak{m}}\langle \partial_{x_1}, \dots, \partial_{x_n} \rangle$.

    By Nakayama's Lemma, there exist $s_1, \dots, s_{\ell} \in M_{\mathfrak{m}}$ such that
    \[
        M_{\mathfrak{m}} = \sum_{i=1}^{\ell} R_{\mathfrak{m}} \cdot s_i
    \]
    and
    \[
        M_{\mathfrak{m}} \otimes R_{\mathfrak{m}}/\mathfrak{m} = \bigoplus_{i=1}^{\ell} \overline{s_i} \cdot R_{\mathfrak{m}}/\mathfrak{m}.
    \]

    If $M_{\mathfrak{m}}$ is not free, then we have a non-trivial relation:
    \[
        \sum f_i s_i = 0, \quad f_i \in R_{\mathfrak{m}}.
    \]
    Since $\overline{s_i}$ are free generators, we must have $f_i \in \mathfrak{m}$ for all $i$.

    Define the order of vanishing:
    \[
        d(f_i) = \max \{ k \mid f_i \in \mathfrak{m}^k \}
    \]
    \[
        d(f_1, \dots, f_{\ell}) = \min_{i} \{ d(f_i) \} \ge 1.
    \]

    For some $j$, apply the derivation $\partial_{x_j}$:
    \[
        \partial_{x_j} \left( \sum f_i s_i \right) = \sum \partial_{x_j}(f_i) s_i + \sum f_i \partial_{x_j}(s_i) = \sum g_i s_i = 0.
    \]
    One can see that $d(g_1, \dots, g_{\ell}) < d(f_1, \dots, f_{\ell})$.

    If $d(g_1, \dots, g_{\ell}) \ge 1$, we take $\partial_{x_j}(\sum g_i s_i)$ again.
    Finally, by repeating this process, we get:
    \[
        \sum h_i s_i = 0 \quad \text{s.t. } d(h_1, \dots, h_{\ell}) = 0.
    \]
    So $h_i \notin \mathfrak{m}$ for some $i$.
    This implies $\sum \overline{h_i} \overline{s_i} = 0$ with some $\overline{h_i} \neq 0$, contradicting the assumption that $\overline{s_i}$ are free generators.

    Therefore, $M_{\mathfrak{m}}$ is free for any maximal ideal $\mathfrak{m}$, which implies $M$ is a locally free $R$-module (an integrable connection).
    So $(1) \iff (2)$.

    \textbf{(1) $\implies$ (3)} is known (marked as Homework). \\

    \textbf{(3) $\implies$ (1)}.
    
    Take a good filtration $F_{\bullet}M$ on $M$.
    Recall that $\operatorname{ch}(M) = \operatorname{supp}(\operatorname{gr}_\bullet M) = V(\operatorname{Ann}(\operatorname{gr}_\bullet M))$.
    Given condition (3), $\operatorname{ch}(M) = T^*_X X$ (the zero section), which implies the support is defined by the ideal generated by $\xi_1, \dots, \xi_n$.

    By Hilbert's Nullstellensatz:
    \[
        \operatorname{rad}(\operatorname{Ann}(\operatorname{gr}_\bullet M)) = \sqrt{\operatorname{Ann}(\operatorname{gr}_\bullet M)} = I := (\xi_1, \dots, \xi_n) \subseteq R[\xi_1, \dots, \xi_n].
    \]
    
    This implies that for $m \gg 0$:
    \[
        I^m \subseteq \operatorname{Ann}(\operatorname{gr}_\bullet M).
    \]
    Meanwhile, the ideal power $I^m$ is generated by monomials of degree $m$:
    \[
        I^m = \{ \xi^\alpha \mid |\alpha| = \sum \alpha_i = m \}.
    \]
    The condition $I^m \cdot \operatorname{gr}_\bullet M = 0$ means that the principal symbols of order $m$ annihilate the graded module.
    Translated to the filtration $F_{\bullet}M$, this implies:
    \[
        \text{if } |\alpha| = m, \text{ then } \partial^\alpha F_j M \subseteq F_{j+m-1} M \quad \forall j.
    \]

    Since we chose a good filtration, for $j \gg 0$, we have the property $F_{m+j} M = F_m \mathcal{D}_X \cdot F_j M$.
    We can expand the action of $F_m \mathcal{D}_X$:
    \begin{align*}
        F_{m+j} M &= F_m \mathcal{D}_X \cdot F_j M \\
        &= \sum_{|\alpha| \le m} R \cdot \partial^\alpha \cdot F_j M \\
        &\subseteq F_{m+j-1} M.
    \end{align*}

    This chain of inclusions implies that the filtration stabilizes:
    \[
        \Rightarrow F_\ell M = F_{\ell+1} M = \dots \quad \forall \ell \gg 0.
    \]
    \[
        \Rightarrow F_\ell M = M.
    \]
    Since each $F_\ell M$ is finitely generated over $R$ (by definition of good filtration), this implies $M$ is finitely generated over $R$ (i.e., $M$ is f.g. / $R$).
    
    This proves condition (2). Since we have already shown $(2) \iff (1)$, the proof is complete.
\end{proof}

\begin{remark}
  To see why $d(g_1, \dots, g_{\ell}) < d(f_1, \dots, f_{\ell})$, note that
  \[
      g_i = \partial_{x_j}(f_i) + f_i \cdot (\text{some term}).
  \]
  Since $f_i \in \mathfrak{m}^{d(f_1, \dots, f_{\ell})}$, the term $f_i \cdot (\text{some term})$ is in $\mathfrak{m}^{d(f_1, \dots, f_{\ell})}$.
  It suffices to analyze $\partial_{x_j}(f_i)$.\\
  By the definition of order of vanishing, we can write:
  \[
      f_i = \sum_{|\beta| = d(f_1, \dots, f_{\ell})} c_\beta x^\beta \quad \text{for } c_\beta \in R_{\mathfrak{m}}
  \]
  with some $c_\beta \neq 0$.
  Applying $\partial_{x_j}$ yields:
  \[
      \partial_{x_j}(f_i) = \sum_{|\beta| = d(f_1, \dots, f_{\ell})} c_\beta \beta_j x^{\beta - e_j} + \text{terms in } \mathfrak{m}^{d(f_1, \dots, f_{\ell})},
  \]
  where $e_j$ is the standard basis vector with $1$ in the $j$-th position.
  Since at least one $\beta \neq 0$, we have $\partial_{x_j}(f_i) \notin \mathfrak{m}^{d(f_1, \dots, f_{\ell})}$ for some $x_j$.
\end{remark}

\chapter{Gabber's Theorem: Involutivity of Characteristic varieties}
\section{Gabber's Involutivity Theorem}
Let $B$ be a commutative $\C$-algebra.

\begin{definition}
  \leavevmode
  \begin{enumerate}
    \item A $\C$-bilinear map $\{\,,\,\}\colon B\otimes_{\C}B\to B$ is a \textbf{Poisson bracket}
          if it is a Lie bracket and satisfies the Leibniz rule $\{a_1a_2,b\}=a_1\{a_2,b\}+a_2\{a_1,b\}$ for all $a_1,a_2,b\in B$.
    \item If $B$ is Poisson, a closed subset $V\subseteq\Spec B$ with radical ideal $I_V$ is
          \textbf{coisotropic} (or \textbf{involutive}) if $\{I_V,I_V\}\subseteq I_V$.
  \end{enumerate}
\end{definition}

\noindent $\Gamma_\bullet A : \text{a } \mathbb{Z}_{\ge 0}\text{-filtered } \mathbb{C}\text{-algebra}$
\begin{definition}
  \leavevmode
  \begin{enumerate}
    \item $\Gamma_\bullet A$ is called \textbf{almost commutative} if $\operatorname{gr}^\Gamma_\bullet A$ is a commutative $\mathbb{C}$-algebra of finite type.
    \item For $x \in \Gamma_k A$, we use $\sigma_k(x) \text{ to denote its image in } \frac{\Gamma_k A}{\Gamma_{k-1} A}$.
    \item For $x \in \Gamma_k A$, $y \in \Gamma_{\ell} A$, we define $\{ \sigma_k(x), \sigma_{\ell}(y) \} = \sigma_{k+\ell-1}(xy - yx)$
          (Since $\operatorname{gr}_\bullet A$ is commutative, $xy - yx \in \Gamma_{k+\ell-1} A$)
  \end{enumerate}
\end{definition}

\begin{lemma}[Lemma~6.3]
  $(\gr^{\Gamma}_\bullet A,\{\,\,\})$ is a Poisson algebra.
\end{lemma}
\begin{proof}
  Homework.
\end{proof}

\begin{example}
  For a smooth affine $X$, the identification $\gr \D_X\cong \Sym^{\bullet}\T_X$ transports the
  Poisson structure on $\gr \D_X$ to $T^{\ast}X$.
\end{example}

If $\mathcal M$ is a finitely generated $A$-module with a good filtration, the characteristic
variety $\Ch(\mathcal M):=\Supp(\gr^{\Gamma}_\bullet\!\mathcal M)$ is a closed subset of
$\Spec\gr^{\Gamma}_\bullet A$.

\begin{theorem}[Gabber's~Involutivity~I]\label{thm:gabber-I}
  If $\mathcal M$ is a finitely generated $A$-module, then $\Ch(\mathcal M)\subseteq\Spec\gr^{\Gamma}_\bullet A$ is involutive.
\end{theorem}

\begin{corollary}
  For a finitely generated $\mathcal{D}_X$-module $\mathcal{M}$ on a smooth affine variety $X$, the characteristic variety $\operatorname{Ch}(\mathcal{M})$ inside $T^{\ast}X$ is conic and involutive. In particular $\dim \operatorname{Ch}(\mathcal{M}) \ge n$.
\end{corollary}

To prove Theorem~\ref{thm:gabber-I}, we introduce Gabber rings.

\begin{definition}
  Let $\mathbb E:=\C[u]/(u^2)$ be the \textbf{ring of dual numbers}. An $\mathbb E$-algebra $T$ (not
  necessarily commutative) is a \textbf{Gabber ring} if
  \begin{enumerate}
    \item $T$ is flat over $\mathbb E$, and
    \item $\overline T:=T/uT$ is a commutative $\C$-algebra of finite type.
  \end{enumerate}
\end{definition}

\begin{lemma}\label{E-flatness}
  For a $T$-module $\underline{\mathcal M}$, the following are equivalent:
  \begin{enumerate}
    \item $\underline{\mathcal M}$ is $\mathbb E$-flat.
    \item The map $\underline{\mathcal M}/u\underline{\mathcal M} \to u\underline{\mathcal M},\ \bar{m} \mapsto um$ is an isomorphism.
  \end{enumerate}
\end{lemma}
\begin{proof}
  Recall that for an $R$-module $N$ over a commutative ring $R$,
  \[
    N \text{ is } R\text{-flat} \iff \forall \text{ proper ideal } I \subseteq R, \quad I \otimes_R N \to N \text{ is injective.}
  \]

  So $\underline{M}$ is $\mathbb{E}$-flat $\iff (u) \otimes_{\mathbb{E}} \underline{M} \hookrightarrow \underline{M}$
  (since $(u)$ is the only non-trivial ideal of $\mathbb{E}$).

  Note that the image of the map is $(u)\underline{M} = u \cdot \underline{M}$.
  Consider the map $\mathbb{E}/(u) \to (u)$ given by $1 \mapsto u$. This map is an $\mathbb{E}$-module isomorphism.

  Thus, we have an $\mathbb{E}$-module isomorphism:
  \[
    \underline{M}/(u)\underline{M} = \mathbb{E}/(u) \otimes_{\mathbb{E}} \underline{M} \xrightarrow{\sim} (u) \otimes_{\mathbb{E}} \underline{M}.
  \]

  Therefore, $\underline{M}$ is $\mathbb{E}$-flat $\iff$ the composite map $\underline{M}/(u)\underline{M} \longrightarrow u \cdot \underline{M} $
  is an isomorphism.
\end{proof}

Whenever $\underline{\mathcal M}$ is $\mathbb E$-flat and $\overline a,\overline b\in\overline T$, we can pick
lifts $a,b\in T$ and define $\{\overline a,\overline b\} := \overline c$ where $c$ solves $u\cdot c = ab-ba$. Since $\overline{T}$ is commutative, such a $c$ exists.

\begin{lemma}\label{lem:poisson}
  $(\overline T,\{\,\,\})$ is a Poisson algebra.
\end{lemma}
\begin{proof}
  Skew-symmetry is immediate from the definition. Take lifts $a,b,c\in T$ of
  $\overline a,\overline b,\overline c$. Direct calculation gives
  \begin{align*}
    \{\overline a,\{\overline b,\overline c\}\}
     & = \overline{a(bc-cb) - (bc-cb) a}, \\
    \{\overline b,\{\overline c,\overline a\}\}
     & = \overline{b(ca-ac) - (ca-ac) b}, \\
    \{\overline c,\{\overline a,\overline b\}\}
     & = \overline{c(ab-ba) - (ab-ba) c}.
  \end{align*}
  Adding the three lines shows Jacobi. For the Leibniz rule we compute
  \begin{align*}
    \{\overline a\,\overline a',\overline b\}
     & = \overline{aa' b - ba a'}
    = \overline{a(a' b - b a')} + \overline{(ab - ba)a'}                                        \\
     & = \overline a\,\{\overline a',\overline b\} + \{\overline a,\overline b\}\,\overline a',
  \end{align*}
  and similarly in the second argument. Hence $(\overline T,\{\,\,\})$ is Poisson.
\end{proof}

\begin{theorem}[Gabber's~Involutivity~II]\label{thm:gabber-II}
  Let $T$ be a Gabber ring and $\underline{\mathcal M}$ a finitely generated $\mathbb E$-flat $T$-module. If
  $I:=\rad\Ann_{\overline T}(\overline{\underline{\mathcal M}})$, then $\{I,I\}\subseteq I$.
\end{theorem}

\begin{proof}[Theorem~\ref{thm:gabber-II} implies Theorem~\ref{thm:gabber-I}]
  Put $\Rees^{\Gamma}_\bullet(A):=\bigoplus_{p\ge0}\Gamma_pA\,u^p\subseteq A[u]$.  This is a finitely generated
  $\C[u]$-algebra, and by construction it is torsion free over $\C[u]$; since $\C[u]$ is a PID, the
  module $\Rees^{\Gamma}_\bullet(A)$ is therefore flat over $\C[u]$.  Write
  \[
    T := \Rees^{\Gamma}_\bullet(A)\otimes_{\C[u]}\mathbb E
    = \frac{\Rees^{\Gamma}_\bullet(A)}{u^2\Rees^{\Gamma}_\bullet(A)},
    \qquad
    \overline T \cong \frac{\Rees^{\Gamma}_\bullet(A)}{u\,\Rees^{\Gamma}_\bullet(A)}
    = \gr^{\Gamma}_\bullet A.
  \]
  Thus $T$ is a Gabber ring with reduction $\overline T=\gr^{\Gamma}_\bullet A$.  If $F_\bullet\underline{\mathcal M}$ is a
  good filtration on a finitely generated $A$-module, its Rees module
  $\Rees^{F}_\bullet(\underline{\mathcal M})=\bigoplus_{p\ge0}F_p\underline{\mathcal M}\,u^p$ is again torsion free over $\C[u]$,
  hence flat.  After base change we obtain the $\mathbb E$-flat $T$-module
  \[
    \underline{\mathcal{M}} := \Rees^{F}_\bullet(\underline{\mathcal M})\otimes_{\C[u]}\mathbb E
    = \frac{\Rees^{F}_\bullet(\underline{\mathcal M})}{u^2\Rees^{F}_\bullet(\underline{\mathcal M})},
  \]
  whose reduction modulo $u$ is
  $\overline{\underline{\mathcal{M}}}=\Rees^{F}_\bullet(\underline{\mathcal M})/u\,\Rees^{F}_\bullet(\underline{\mathcal M})\cong\gr^{F}_\bullet\underline{\mathcal M}$.  Now
  Theorem~\ref{thm:gabber-II} applied to $(T,\underline{\mathcal{M}})$ says that
  $\Supp(\gr^{F}_\bullet\underline{\mathcal M})\subseteq\Spec\gr^{\Gamma}_\bullet A$ is involutive.

  \begin{lemma}\label{lem:rees}
    For $\overline T=\gr^{\Gamma}_\bullet A$ and $T=\frac{\Rees^{\Gamma}_\bullet(A)}{u^2\Rees^{\Gamma}_\bullet(A)}$
    (the $\mathbb E$-algebra that appeared above), the induced bracket coincides with the intrinsic
    bracket defined earlier on $\gr^{\Gamma}_\bullet A$.
  \end{lemma}
  \begin{proof}
    The editor of this lecture note think it is straightforward. And it is assigned as homework.
  \end{proof}

  Then we complete the deduction of Theorem~\ref{thm:gabber-I} from Theorem~\ref{thm:gabber-II}.
\end{proof}

\begin{proof}[Proof of Theorem~\ref{thm:gabber-II}]



  Write
  \[
    I=\rad\Ann_{\overline T}(\overline{\underline{\mathcal M}})
    = \bigcap_{\substack{
        \text{minimal prime ideal } \\
        \mathfrak p_i \supseteq \Ann_{\overline T}(\overline{\underline{\mathcal M}})
      }} \mathfrak p_i.
  \]
  Since $\{I,I\}\subseteq I$ iff $\{\mathfrak p_i,\mathfrak p_i\}\subseteq \mathfrak p_i$ for every minimal prime, we fix one such $\mathfrak p$ and work locally.

  \begin{lemma}\label{lem:noether}
    Let $R$ be a finite type $k$-algebra of dimension $\ell+d$ and $\mathfrak p\subset R$ a prime of codimension
    $\ell$. There exist algebraically independent elements $x_1,\dots,x_\ell,y_1,\dots,y_d\in R$ such that
    \begin{enumerate}
      \item $R$ is finite over $k[x_1,\dots,x_\ell,y_1,\dots,y_d]$.
      \item $\mathfrak p\cap k[x_1,\dots,x_\ell,y_1,\dots,y_d]=(x_1,\dots,x_\ell)$.
      \item There exists $0\ne f\in k[y_1,\dots,y_d]$ for which $(R/\mathfrak p)_f$ is free over $k[y_1,\dots,y_d]_f$.
    \end{enumerate}
    \noindent
    This is a version of the Noether Normalization Theorem (see, for instance, Eisenbud,
    \textbf{Commutative Algebra with a View Toward Algebraic Geometry}).
  \end{lemma}

  \begin{proof}[Sketch of Lemma~\ref{lem:noether}]
    By \cite[Theorem 13.3(Noether Normalization).]{eisenbud1995commutative} We proved (1) and (2).

    To prove (3), we use the following result called \textbf{Generic Freeness},
    This notion will be used again in the proof of Theorem~\ref{thm:gabber-II}:

    \begin{lemma}[{\cite[\href{https://stacks.math.columbia.edu/tag/051R}{Tag 051R}]{stacks-project} or \cite[Theorem 14.4(Grothendieck's Generic Freeness Lemma).]{eisenbud1995commutative}}]\label{algebra-lemma-generic-flatness-Noetherian}
      \mbox{}\\
      Let $R \to S$ be a ring map.
      Let $M$ be an $S$-module.
      Assume
      \begin{enumerate}
        \item $R$ is Noetherian,
        \item $R$ is a domain,
        \item $R \to S$ is of finite type, and
        \item $M$ is a finite type $S$-module.
      \end{enumerate}
      Then there exists a nonzero $f \in R$ such that
      $M_f$ is a free $R_f$-module.% Lemma content goes here
    \end{lemma}

    Using this Lemma~\ref{algebra-lemma-generic-flatness-Noetherian}, we proved (3) of Lemma~\ref{lem:noether}.
  \end{proof}



  Choose a minimal prime $\mathfrak p$ of $\overline T$ containing $\Ann_{\overline T}(\overline{\underline{\mathcal M}})$
  and write $d=\dim \mathfrak p$. Since $\overline T$ is a finite type $\C$-algebra, we have primary decomposition that is:
  \[
    \rad\Ann_{\overline T}(\overline{\underline{\mathcal M}})
    = \mathfrak p \cap \Bigl(\bigcap_{i\ne j} \mathfrak p_i\Bigr).
  \]
  and we assume $\mathfrak p$ and $\mathfrak p_i$ are all distinct minimal primes.

  Choose $g\in\bigcap_{i\ne j}\mathfrak p_i\smallsetminus \mathfrak p$. Localizing at $g$ removes all
  minimal primes except $\mathfrak p$, so $(\overline{\underline{\mathcal M}})_{g}$ is killed by some power
  $\mathfrak p_{g}^\ell$.

  Since $\overline{\underline{\mathcal M}}$ is finitely generated over $\overline T$, by repeatedly applying Lemma~\ref{algebra-lemma-generic-flatness-Noetherian}, we get:
  \[
    \exists g,\, \forall k,\, \frac{\mathfrak p^k_g \overline{\underline{\mathcal M}}}{\mathfrak p^{k+1}_g \overline{\underline{\mathcal M}}} \text{ is free over } \overline{T}_g/\mathfrak p_g.
  \]
  (One should note that $\overline{T}/ \mathfrak p$ is a noetherian domain, and $\mathfrak p^k \overline{\underline{\mathcal M}}/ \mathfrak p^{k+1} \overline{\underline{\mathcal M}}$ is finitely generated over $\overline{T}/ \mathfrak p$.)


  Apply Lemma~\ref{lem:noether} to the finite-type algebra $R=\overline T_{g}$. We obtain algebraically
  independent elements $y_1=t_1,\dots,y_d=t_d\in\overline T_{g}$ and some $h\in\C[t_1,\dots,t_d]$ such that both $\overline T_{g}$ and $(\overline T_{g}/\mathfrak p_{g})_h$ are finite and free
  over $S_h:=\C[t_1,\dots,t_d]_h$ (i.e. $S:=\C[t_1,\dots,t_d]$). Replacing $\bar f$ by $g h$ if necessary (and keeping the notation $\bar f$
  for the product).

  With this choice we have:
  $\mathfrak p_{\bar f}^\ell\overline{\underline{\mathcal M}}_{\bar f}=0$ for some $\ell>0$, and each successive
  quotient $\mathfrak p_{\bar f}^i\overline{\underline{\mathcal M}}_{\bar f} / \mathfrak p_{\bar f}^{i+1}\overline{\underline{\mathcal M}}_{\bar f}$
  is a free $S_h$-module. \\

  For later use, fix $S_h$-bases
  $\{\overline m_{i,j}\}$ such that for each $i$, $\{\overline m_{i,j}\}$(fix $i$, and change $j$) is a $S_h$-basis of
  $\mathfrak p_{\bar f}^i\overline{\underline{\mathcal M}}_{\bar f}$.
  This basis $\{\overline m_{i,j}\}$ can also be regarded as a $S_h$-basis of $\overline{\underline{\mathcal M}}_{\bar f}$
  (The $S_h$-action on $\overline{\underline{\mathcal M}}_{\bar f}$ is through $\overline{T}_{\bar f}$, not through $\overline{T}_{\bar f}/\mathfrak p_{\bar f}$).\\

  Relabelling we obtain:
  \[
    \underbrace{\overline m_{k_0},\dots,\overline m_{k_1-1}}_{\text{basis of }\overline{\underline{\mathcal M}}_{\bar f}/\mathfrak p_{\bar f}\overline{\underline{\mathcal M}}_{\bar f}},\quad
    \underbrace{\overline m_{k_1},\dots,\overline m_{k_2-1}}_{\text{basis of }\mathfrak p_{\bar f}\overline{\underline{\mathcal M}}_{\bar f}/\mathfrak p_{\bar f}^2\overline{\underline{\mathcal M}}_{\bar f}},\quad
    \ldots,\quad
    \underbrace{\overline m_{k_{\ell-1}},\dots,\overline m_{k_\ell-1}}_{\text{basis of }\mathfrak p_{\bar f}^{\ell-1}\overline{\underline{\mathcal M}}_{\bar f}}.
  \]

  We then take a lift $f\in T$ of $\bar f\in \overline T$.

  \begin{lemma}[Lemma~6.12]\label{lem:matrix}
    Let $a,b\in T$ have images $\bar a,\bar b\in \mathfrak p$. Then there exist integers
    $n_1,n_2,n_3\ge0$ and elements $e_{ij}\in T$ with $\overline{e_{ij}} \in S\subseteq \overline{T}$ such that
    \begin{enumerate}
      \item $f^{n_1}[f^{n_2}a,f^{n_3}b] m_i = u\sum_j e_{ij} m_j$ for every $i$, and
      \item $\sum_j \overline{e_{ij}} = 0$ for every $i$.
    \end{enumerate}
  \end{lemma}
  \begin{proof}
    If $k_{j_1} \le i < k_{j_1+1}$ for some $j_1$, then we have
    \[
      f^{n_2} a m_i = \sum_{j \ge k_{j_1+1}} u_{ij} m_j + u \sum_j \alpha_{ij} m_j,\quad \bar{u}_{ij}, \bar{\alpha}_{ij} \in S
    \]


    \begin{remark}
      By abuse of notation, let us clarify how to get this expression, this equality makes sense in $\underline{\mathcal{M}}$:

      \begin{enumerate}
        \item \textbf{Filtration:}
              $\overline{m_i} \in \mathfrak{p}^{k_i}_{\bar f}\overline{\underline{\mathcal{M}}}_{\bar f}/\mathfrak{p}^{k_i+1}_{\bar f}\overline{\underline{\mathcal{M}}}_{\bar f}$. Then $\overline{a}\overline{m_i} \in \mathfrak{p}^{k_i+1}_{\bar f}\overline{\underline{\mathcal{M}}}_{\bar f}$. So we have $\overline{a}\overline{m_i} = \sum_{j \ge k_{j_1+1}} \overline{u_{ij}}\overline{m_j} \in \overline{\underline{\mathcal{M}}}_{\bar f}$.

        \item \textbf{Localization:}
              So we have $\overline{f^{n_2} a m_i} = \sum_{j \ge k_{j_1+1}} \overline{u_{ij}} \overline{m_j} \in \overline{\underline{\mathcal{M}}}$, here we canceled the denominators of $\overline{u_{ij}}$, and use the same symbol to denote the new $\overline{u_{ij}}$.

        \item \textbf{Flatness:}
              By Lemma~\ref{E-flatness}, we have the isomorphism $\overline{\underline{\mathcal{M}}} \xrightarrow{\sim} u \underline{\mathcal{M}}$. So every element in $u\underline{\mathcal{M}}$ can be expressed by the form
              \[
                u \cdot \text{ elements lifted from }\overline{\underline{\mathcal{M}}}.
              \]
              So we have $f^{n_2} a m_i = \sum_{j \ge k_{j_1+1}} u_{ij} m_j + u \sum_j \alpha_{ij} m_j\quad  \in \underline{\mathcal{M}}$.
      \end{enumerate}
    \end{remark}



    Similarly,
    \[
      f^{n_3} b m_i = \sum_{j \ge k_{j_1+1}} v_{ij} m_j + u \sum_j \beta_{ij} m_j,\quad \bar{v}_{ij}, \bar{\beta}_{ij} \in S
    \]
    Then, since $u^2=0$,
    \begin{align*}
      f^{n_2} a f^{n_3} b m_i = \sum_{j \ge k_{j_1+1}} \left( v_{ij} f^{n_2} a + [f^{n_2}a, v_{ij}] \right) m_j + u \sum_j \beta_{ij} f^{n_2} a m_j
    \end{align*}

    \noindent Take $U = [u_{ij}]$, $V = [v_{ij}]$, $\alpha = [\alpha_{ij}]$, $\beta = [\beta_{ij}]$, we get:

    \begin{align*}
      f^{n_2} a f^{n_3} b m_i
      = \sum_{j \ge k_{j_1+1}} \left( (VU)_{ij} + [f^{n_2}a, v_{ij}] \right) m_j
      + u \sum_j \left( (V \cdot \alpha)_{ij} + (\beta \cdot U)_{ij} \right) m_j
    \end{align*}



    \noindent Similarly,
    \begin{align*}
      f^{n_3} b f^{n_2} a m_i
      = \sum_{j \ge k_{j_1+1}} \left( (UV)_{ij} + [f^{n_3}b, u_{ij}] \right) m_j
      + u \sum_j \left( (U \cdot \beta)_{ij} + (\alpha \cdot V)_{ij} \right) m_j
    \end{align*}

    \noindent Therefore,
    \[
      [f^{n_2}a, f^{n_3}b] m_i = \sum_{j \ge k_{j_1+1}} c_{ij} m_j + u \sum_j \gamma_{ij} m_j
    \]
    where
    \[
      c_{ij} = [V, U]_{ij} + [f^{n_2}a, v_{ij}] - [f^{n_3}b, u_{ij}],
    \]
    and
    \[
      [\gamma_{ij}] = [V, \alpha] + [\beta, U].
    \]
    \noindent On the other hand, since $\overline{T}$ is commutative
    \[
      [f^{n_2}a, f^{n_3}b] = u \cdot c \quad \text{for some } c \in T
    \]

    \noindent This implies
    \[
      [f^{n_2}a, f^{n_3}b] m_i = u \cdot c m_i = u \cdot \sum_j d'_{ij} m_j, \quad \overline{d'_{ij}} \in S_{\bar f}
    \]
    \noindent (Here we use the $E$-flatness of $\mathcal{M}$, that is the isomorphism $\overline{\underline{\mathcal{M}}} \xrightarrow{\sim} u \underline{\mathcal{M}}$.)

    \noindent Canceling the denominators of $d'_{ij}$, we set
    \[
      f^{n_1} [f^{n_2}a, f^{n_3}b] m_i = u \cdot \sum_j e_{ij} m_j, \quad \overline{e_{ij}} \in S
    \]
    So
    \[
      f^{n_1} \cdot [c_{ij}] + u f^{n_1} [\gamma_{ij}] = u [e_{ij}]
    \]
    But $[c_{ij}]$ is traceless because it is a sum of a commutator and a strict upper triangular matrix.
    $[\gamma_{ij}]$ is also traceless similarly.
    So $[\overline{e_{ij}}]$ is traceless.
  \end{proof}

  \noindent Now we go back to prove $\{\frak{p}, \frak{p}\} \subseteq \frak{p}$.\\

  \noindent We know
  $  f^{n_1} [f^{n_2}a, f^{n_3}b] = c $
  \begin{align*}
    c & = f^{n_1} f^{n_2} a f^{n_3} b - f^{n_1} f^{n_3} b f^{n_2} a                                          \\
      & = f^{n_1+n_2+n_3} ab + f^{n_1+n_2} [a, f^{n_3}] b  - f^{n_1+n_2+n_3} ba - f^{n_1+n_3} [b, f^{n_2}] a \\
      & = f^{n_1+n_2+n_3} [a, b] + f^{n_1+n_2} [a, f^{n_3}] b - f^{n_1+n_3} [b, f^{n_2}] a
  \end{align*}

  \noindent This implies $\overline{c} \in \bar{f}^{n_1+n_2+n_3} \{\bar{a}, \bar{b}\} + \frak{p}$.\\

  \noindent By Lemma~\ref{lem:matrix}, $\text{the trace of } \overline{c} \text{ on } \overline{\underline{\mathcal{M}}}_{\bar f} = \sum_i \overline{e_{ii}} = 0$.
  Since $\frak{p}^l_{\bar f} \cdot \overline{\underline{\mathcal{M}}}_{\bar f} = 0$, its trace is also 0.
  So the trace of $\{\bar{a}, \bar{b}\}$ on $\overline{\underline{\mathcal{M}}}_{\bar f}$ is also 0 (since $\bar{f} \notin \frak{p}$).\\

  \noindent Then $\forall x \in \overline{T}_{\bar f}$, $\{x\bar{a}, \bar{b}\} = x\{\bar{a}, \bar{b}\} + \{x, \bar{b}\}\bar{a} \in x\{\bar{a}, \bar{b}\} + \frak{p}_{\bar f}$.
  Similarly, we conclude
  \[
    \text{the trace of } x\{\bar{a}, \bar{b}\} \text{ on } \overline{\underline{\mathcal{M}}}_{\bar f} \text{ is also } 0.
  \]

  \noindent Meanwhile, since $\overline{\underline{\mathcal{M}}}_{\bar f}$ is free over $B = \overline{T}_{\bar f} / \frak{p}_{\bar f}$ and $S_{\bar f} \hookrightarrow B$ is a finite free extension, the trace of $y = x\{\bar{a}, \bar{b}\} \in \overline{T}_{\bar f}$ on $\overline{\underline{\mathcal{M}}}_{\bar f}$ is
  \[
    r \cdot (\text{the trace of } \bar{y} \text{ on } B),
  \]
  where $r$ is the rank of $\overline{\underline{\mathcal{M}}}_{\bar f}$ over $S_{\bar f}$.
  (For this properties of trace one can refer \cite[Theorem 7.8.5]{li2019algebraic})\\

  \noindent Since finite extension is integral and we are over $\mathbb{C}$ (thus the corresponding field extension is separable), the trace form of $S_{\bar f} \hookrightarrow B$ is non-degenerate.

  \noindent So
  \[
    \bar{y} = 0 \in B \implies \{\bar{a}, \bar{b}\} \in \frak{p}.
  \]

  This proves $\{\mathfrak p,\mathfrak p\}\subseteq \mathfrak p$, so $\mathfrak p$ (and thus $I$) is involutive.  The proof of
  Theorem~\ref{thm:gabber-II}---and hence Theorem~\ref{thm:gabber-I}---is complete.

\end{proof}
\chapter{Some symplectic geometry of affine $T^{\ast}X$}
\section{Symplectic vector spaces}

Let $V$ be a $k$-vector space equipped with a bilinear form
$\sigma\colon V\times V\to k$.  We use the following terminology:
\begin{itemize}
  \item The form $\sigma$ is \emph{skew-symmetric} if
        $\sigma(a,b)=-\sigma(b,a)$ for all $a,b\in V$.
  \item A skew-symmetric bilinear form is \emph{nondegenerate} if for every
        $0\ne a\in V$ there exists $v\in V$ such that $\sigma(a,v)\ne 0$.
\end{itemize}
A pair $(V,\sigma)$ is called a \emph{symplectic vector space} if $\sigma$
is a nondegenerate skew-symmetric bilinear form.  In this case
$\dim_k V=2n$ for some $n$.

Nondegeneracy identifies $V$ with its dual: the map
\[
  \sharp \colon V \longrightarrow V^\ast,\qquad
  v \longmapsto \sigma(-,v)
\]
is an isomorphism, called the \emph{Hamiltonian isomorphism}.

\begin{definition}[Def.~7.1]
Let $W\subseteq V$ be a linear subspace.
\begin{enumerate}
  \item The \emph{symplectic orthogonal} of $W$ is
        $W^\perp:=\{v\in V\mid \sigma(v,W)=0\}$.
  \item The subspace $W$ is called \emph{isotropic} if $W\subseteq W^\perp$,
        \emph{Lagrangian} if $W=W^\perp$, and \emph{coisotropic} if
        $W\supseteq W^\perp$.
\end{enumerate}
\end{definition}

\begin{lemma}[Lemma~7.2]
Let $W\subseteq V$ be as above.  If $W$ is isotropic then $\dim_k W\le n$.
If $W$ is Lagrangian then $\dim_k W=n$, and if $W$ is coisotropic then
$\dim_k W\ge n$.
\end{lemma}
\begin{proof}
The Hamiltonian isomorphism identifies $W^\perp$ with the kernel of the
composition $V\xrightarrow{\sharp}V^\ast\twoheadrightarrow W^\ast$.  Hence
$\dim_k W^\perp=2n-\dim_k W$.  When $W\subseteq W^\perp$ we obtain
$\dim_k W\le n$.  Equality occurs exactly when $W=W^\perp$, the defining
property of a Lagrangian subspace.  Dually, $W\supseteq W^\perp$ implies
$\dim_k W\ge n$, which is the coisotropic case.
\end{proof}

\begin{example}[Ex.~7.2]\label{ex:symplectic-standard}
Let $W$ be an $n$-dimensional $k$-vector space with dual $W^\ast$.  Then
$V:=W\oplus W^\ast$ becomes a symplectic vector space when we define
\[
  \sigma\big((x,\xi),(x',\xi')\big)
  :=\langle x',\xi\rangle-\langle x,\xi'\rangle,
  \qquad (x,\xi),(x',\xi')\in V,
\]
where $\langle\cdot,\cdot\rangle$ denotes the evaluation pairing between
$W$ and $W^\ast$.
\end{example}

\section{Symplectic structure on $T^\ast X$}

Let $X=\Spec R$ be a smooth affine algebraic variety over $\C$ (more generally
over an algebraically closed field $k$ of characteristic $0$) with algebraic
coordinates $(x_1,\dots,x_n)$.  The cotangent bundle is
\[
  T^\ast X = \Spec \gr \D_X = \Spec \Sym^\bullet T_X
           \simeq \Spec R[\xi_1,\dots,\xi_n],
\]
where $\Sym^\bullet T_X:=\bigoplus_{i\ge0}\Sym^i T_X$ carries the natural
grading.

\begin{definition}[Def.~8.1]
A subvariety $V\subseteq T^\ast X$ is called \emph{conic} if its defining ideal
$I(V)\subseteq \Sym^\bullet T_X$ is homogeneous.
\end{definition}

\begin{remark}
Under the identification $T^\ast X\simeq X\times k^n$, each fiber is an
$n$-dimensional $k$-vector space, so $T^\ast X$ carries a canonical $k^\times$
action.  A subvariety $V\subseteq T^\ast X$ is conic precisely when it is
stable under this action.
\end{remark}

The canonical $1$-form on the cotangent bundle is
\[
  \alpha_X := \sum_{i=1}^n \xi_i\,dx_i \in \Omega^1_{T^\ast X},
\]
which does not depend on the chosen coordinates.  Its exterior derivative
\[
  \sigma_X := d\alpha_X = \sum_{i=1}^n d\xi_i\wedge dx_i \in \Omega^2_{T^\ast X}
\]
is the standard symplectic form on $T^\ast X$.

For a closed point $p\in T^\ast X$ with maximal ideal $\mathfrak m_p$ we have
the tangent space
\[
  T_{T^\ast X,p}=\Der_k(\mathfrak m_p,k)
               \simeq \Omega^1_{T^\ast X}\otimes_S S/\mathfrak m_p
               \simeq \mathfrak m_p/\mathfrak m_p^2.
\]
Evaluating $\sigma_X$ at $p$ yields a bilinear form
$\sigma_p\colon T_{T^\ast X,p}\times T_{T^\ast X,p}\to k$.

\begin{lemma}[Lemma~8.1]
For every closed point $p\in T^\ast X$, the pair $(T_{T^\ast X,p},\sigma_p)$ is
a symplectic vector space.
\end{lemma}
\begin{proof}
The identification above shows that $T_{T^\ast X,p}\simeq
\Omega^1_{T^\ast X}\otimes_S k(p)$.  Example~\ref{ex:symplectic-standard}
(with $V=\Omega^1_{T^\ast X}\otimes k(p)$ and $\sigma=\sigma_p$) applies
directly, so $\sigma_p$ is a nondegenerate skew-symmetric form.
\end{proof}

The global form $\sigma_X$ yields the Hamiltonian isomorphism
\[
  \mathcal H\colon \Omega^1_{T^\ast X}\xrightarrow{\sim}T_{T^\ast X},\qquad
  \beta\longmapsto \mathcal H(\beta),
\]
characterized by $\sigma_X(-,\mathcal H(\beta))=\beta$.  For any function
$f\in \Sym^\bullet T_X$ we define its \emph{Hamiltonian vector field} to be
\[
  H_f:=\mathcal H(df)\in \Gamma(T^\ast X,T_{T^\ast X}),
\]
which satisfies $\sigma_X(v,H_f)=df(v)$ for all vector fields $v$.

\begin{lemma}[Lemma~8.2]\label{lem:hamiltonian-coordinates}
With respect to the coordinates $(x_1,\dots,x_n,\xi_1,\dots,\xi_n)$ one has
\[
  H_f = \sum_{i=1}^n
        \left(\frac{\partial f}{\partial \xi_i}\frac{\partial}{\partial x_i}
             -\frac{\partial f}{\partial x_i}\frac{\partial}{\partial \xi_i}
        \right).
\]
\end{lemma}
\begin{proof}
Write $H_f=\sum_i \big(\alpha_i\,\partial_{x_i}+\beta_i\,\partial_{\xi_i}\big)$.
Because $\sigma_X(\partial_{x_i},\partial_{x_j})=0$ and
$\sigma_X(\partial_{x_i},\partial_{\xi_j})=\delta_{ij}$, we obtain
\[
  df(\partial_{x_i})=\sigma_X(\partial_{x_i},H_f)=-\beta_i,\qquad
  df(\partial_{\xi_i})=\sigma_X(\partial_{\xi_i},H_f)=\alpha_i.
\]
Thus $\alpha_i=\partial f/\partial \xi_i$ and $\beta_i=-\partial f/\partial
x_i$, yielding the displayed formula.
\end{proof}

The Hamiltonian vector fields endow $\Sym^\bullet T_X$ with a Poisson bracket:
for $f,g\in \Sym^\bullet T_X$ set
\[
  \{f,g\}_\sigma := H_f(g)=\sigma_X(H_f,H_g).
\]
It is skew-symmetric, satisfies the Jacobi identity, and acts as a derivation
in each variable.

\begin{lemma}[Lemma~8.3]
$(\Sym^\bullet T_X,\{-,-\}_\sigma)$ is a Poisson algebra.
\end{lemma}
\begin{proof}
The identities above translate the usual properties of Hamiltonian vector
fields into the axioms of a Poisson bracket; this is standard and was assigned
as homework.
\end{proof}

\begin{lemma}[Lemma~8.4]
Under the identification $\gr \D_X\simeq \Sym^\bullet T_X$, the Poisson bracket
constructed above coincides with $\{-,-\}_\sigma$.
\end{lemma}
\begin{proof}
Because $\gr \D_X$ is a graded $R$-algebra isomorphic to $\Sym^\bullet T_X$, it
is enough to compare the brackets on elements of the form
$f=f_\alpha\xi^\alpha$, $g=g_\beta\xi^\beta$ with $f_\alpha,g_\beta\in R$ and
multi-indices $\alpha,\beta\in\mathbb Z_{\ge0}^n$.  Using
Lemma~\ref{lem:hamiltonian-coordinates} we compute
\[
  \{f,g\}_\sigma
   = \sum_{i=1}^n
       \left(
         \alpha_i f_\alpha\,\frac{\partial g_\beta}{\partial x_i}
         -\beta_i g_\beta\,\frac{\partial f_\alpha}{\partial x_i}
       \right)
       \xi^{\alpha+\beta-e_i},
\]
where $e_i$ is the $i$th standard basis vector.  This matches the explicit
formula obtained earlier for the bracket on $\gr \D_X$, so the two Poisson
structures agree.
\end{proof}

Now let $S:=\Sym^\bullet T_X$, so $T^\ast X=\Spec S$.  Let
$V\subseteq T^\ast X$ be a closed subvariety with radical ideal
$I_V\subseteq S$.  For a smooth closed point $p\in V$ write
$\mathfrak m_p\subset S$ and $\overline{\mathfrak m}_p\subset S/I_V$ for the
corresponding maximal ideals.  The exact sequence
\[
  0 \longrightarrow I_{V,p}/I_{V,p}^2
    \longrightarrow \mathfrak m_p/\mathfrak m_p^2
    \longrightarrow \overline{\mathfrak m}_p/\overline{\mathfrak m}_p^2
    \longrightarrow 0
\]
and we denote $\mathfrak m_p/\mathfrak m_p^2 \longrightarrow \overline{\mathfrak m}_p/\overline{\mathfrak m}_p^2$
in above short exact sequence by \hypertarget{triple_ast}{$\textcolor{red}{(* * *)}$},
and this short exact sequence dualizes to a natural embedding of tangent spaces
\[
  T_{V,p}\simeq (\overline{\mathfrak m}_p/\overline{\mathfrak m}_p^2)^\vee
   \hookrightarrow
  T_{T^\ast X,p}\simeq (\mathfrak m_p/\mathfrak m_p^2)^\vee.
\]
Equivalently, $T_{V,p}$ consists of those vectors $v\in T_{T^\ast X,p}$ such
that $df(v)=0$ for every $f\in I_V$.

\begin{definition}[Def.~8.4]
A closed subvariety $V\subseteq T^\ast X$ is \emph{isotropic} (resp.
\emph{Lagrangian}, \emph{coisotropic}) if for every smooth closed point
$p\in V$ the tangent space $T_{V,p}$ is isotropic (resp. Lagrangian,
coisotropic) inside the symplectic vector space $T_{T^\ast X,p}$.
\end{definition}

\begin{theorem}[Thm.~8.5]
Let $V\subseteq T^\ast X$ be a closed subvariety with ideal $I_V$.  The
following are equivalent:
\begin{enumerate}
  \item $V$ is coisotropic.
  \item $I_V$ is involutive, i.e.\ $\{I_V,I_V\}_\sigma\subseteq I_V$.
\end{enumerate}
\end{theorem}
\begin{proof}
Take a smooth closed point $p \in V$.
\begin{gather*}
    \sigma(v, H_f) = df(v) \quad \text{and} \quad \hyperlink{triple_ast}{\textcolor{red}{(* * *)}} \\
    \Downarrow \\ % 换行直接画箭头
    \begin{align*}
      (\mathcal{T}_{V,p})^\perp &= \{ v \in \mathcal{T}_{T^*X, p} \mid \sigma_p(v, w) = 0, \quad \forall w \in \mathcal{T}_{V,p} \} \\
      &= \text{the space spanned by } H_{f,p} \quad \forall f \in I_V.
    \end{align*}
\end{gather*}

% 中间的方框注释
\noindent If $h \in S$, then $h_p$ denotes the value of $h$ at the point $p$, i.e. $h_p = \bar{h} \in \frac{R}{\mathfrak{m}_p} \simeq k$.


\vspace{1em}

\noindent $V$ coisotropic
\begin{align*}
    &\implies H_{f,p} \in \mathcal{T}_{V,p} \implies \{f, g\}_p \\
    &= H_f(g)_p = dg(H_f)_p = 0, \quad \forall \text{ regular } p.
\end{align*}
\[
    \implies \exists \text{ open } U \subseteq V \quad \text{s.t.} \quad \forall p \in U \text{ is regular.}
\]
\[
    \{f, g\}|_U = 0 \xrightarrow{HW} \{f, g\}|_V = 0.
\]
So (1) $\Rightarrow$ (2).

\vspace{1em}
\noindent (2) $\Rightarrow$ (1) is similar.
\end{proof}

As an immediate corollary we obtain the standard statement for characteristic
varieties.

\begin{theorem}[Thm.~8.6]
Let $\mathcal M$ be a finitely generated $\D_X$-module.  Then its
characteristic variety $\Ch(\mathcal M)\subseteq T^\ast X$ is conic and
coisotropic.  In particular $\dim \Ch(\mathcal M)\ge n$.
\end{theorem}

\section{Conic Lagrangian subvarieties of $T^\ast X$}

Write $Y=\Spec S$ for a smooth affine $k$-variety, and let
$V\subseteq Y$ be a closed subscheme with radical ideal $I=I_V\subseteq S$.
Set $\overline S:=S/I$, so $V=\Spec \overline S$.

\begin{lemma}[Lemma~9.1]\label{lem:kahler-sequence}
There is an exact sequence of $\overline S$-modules
\[
  I/I^2 \xrightarrow{\;\delta\;} \Omega_S\otimes_S \overline S
    \longrightarrow \Omega_{\overline S}\longrightarrow 0,\qquad
  \delta(\overline b)=db\otimes 1.
\]
\end{lemma}
<...>

\chapter{Holonomicity and regular holonomicity}
\section{Holonomic $\mathcal{D}$-modules}
Let $X = \text{spec} R$ be a smooth affine scheme with algebraic coordinates $(x_1, \ldots, x_n)$ and let $M$ be a finite generated $\mathcal{D}_X$-module. We already know $\text{ch}(M) \subset T^*X$ is a conic involutive (or coisotropic) subvariety. In particular, $\dim \text{ch}(M) \geq n$. 

\begin{definition}
    $M$ is called $\textbf{holonomic}$ if $\dim \text{ch}(M) = n$.
\end{definition}

\begin{corollary}
    If $M$ is a holonomic $\mathcal{D}_X$-module, then $\text{ch}(M) \subset T^*X$ is a Lagrangian subvariety, and $$CC(M) = \sum_{E_i \subset X\atop \text{irr subvarieties}} n_i T^*_{E_i} X$$
\end{corollary}

\begin{remark}
    Recall that Lagrangian means $W=W^{\perp}\iff \mathrm{dim}_{k}W=n$, and we define $V\subseteq T^{\ast}X$ is Lagrangian if for any smooth closed point $p\in V$, we have $\mathcal{T}_{V,p}\subseteq \mathcal{T}_{T^{\ast}X,p}$ is Lagrangian. Here $T^{\ast}_{E_{i}}X$ is the conormal bundle of irreducible subvariety $E_{i}$. 
\end{remark}

\begin{example}
    \quad
    \begin{enumerate}
        \item $R$ is a $\mathcal{D}_X$-module. $CC(R) = T^*_X X\simeq X \Rightarrow R$ is holonomic.
        \item Let $M$ be an integrable connection on $X$, we obtain that $ch(M)=T^{\ast}_{X}X$ and $\mathrm{dim}_{k}V(\xi_{1},\dots,\xi_{n})=n$, we have $M$ is holonomic. 
    \end{enumerate}
\end{example}

\section{The de Rham Complexes and Analytification}

Let $M$ be a $D_X$-module, and denote $\partial_i = \partial_{x_i}$ be the vector field along $x_i$.  
Define a complex $$K(M; \partial_1) = [M \xrightarrow{\partial_1} M].$$Here $[M\xrightarrow{\partial_{1}} M]$ means a complex concentrated on $\mathrm{deg}=0$ and $\mathrm{deg}=1$, with $\partial_{1}$ as coboundary map. 

\quad

\begin{note}[The cone of a morphism of complexes]\label{note 6.2}
Let $f: A^\bullet \longrightarrow B^\bullet$ be a morphism of complexes:

\begin{center}
 \begin{tikzcd}
 \cdots \arrow[r] & A^{n-1} \arrow[r, "d_A^{n-1}"] \arrow[d, "f^{n-1}"] & A^n \arrow[r, "d_A^n"] \arrow[d, "f^n"] & A^{n+1} \arrow[r] \arrow[d, "f^{n+1}"] & \cdots \\
 \cdots \arrow[r] & B^{n-1} \arrow[r, "d_B^{n-1}"] & B^n \arrow[r, "d_B^n"] & B^{n+1} \arrow[r] & \cdots
 \end{tikzcd}
\end{center}

then $\mathrm{Cone}(f)=C^{\bullet}$ defined by $C^n = A^n \oplus B^{n-1}$ with differential $$d_C^n = \begin{pmatrix}
 d_A^n & 0 \\
 f^n & -d_B^{n-1}
 \end{pmatrix}: A^n \oplus B^{n-1} \longrightarrow A^{n+1} \oplus B^n.$$The cone construction gives a short exact sequence of complexes:
\begin{center}
\begin{tikzcd}
0 \arrow[r] & B^\bullet[-1] \arrow[r] & Cone(f)^\bullet \arrow[r] & A^\bullet \arrow[r] & 0.
\end{tikzcd}
\end{center}
\end{note}


With Note \ref{note 6.2}, for $k \geq 2$, we inductively define     $$K(M; \partial_1, \cdots, \partial_k) = \text{cone}\left(K(M; \partial_1, \cdots, \partial_{k-1}) \xrightarrow{\partial_k} K(M; \partial_1, \cdots, \partial_{k-1})\right).$$ Then the $\textbf{de Rham complex}$ of $M$ is   $$K(M; \partial_1, \cdots, \partial_n),$$ denoted as $\mathrm{DR}(M).$

\begin{remark}
    In fact, it's Koszul complex with $(\partial_{1},\dots,\partial_{n})$ a finite sequence of central elements in $M$, defined in a mapping cone way.   
\end{remark}

Meanwhile, we have a coordinate-free way to define $\mathrm{DR}(M)$: $$M \to \Omega^{1}_{X} \otimes_{R} M \to \Omega^{2}_{X} \otimes_{R} M \to \cdots \to \Omega^{n}_{X} \otimes_{R} M.$$With $m\rightarrow \sum\limits_{i}dx_{i}\otimes \partial_{i}(m)$, and $\eta \otimes m \mapsto \sum dx_i \wedge \eta \otimes \partial_i(m)$ and so on...

\begin{remark}
    Here $\sum\limits_{i} dx_i \otimes \partial_i \in \Omega^1 \otimes_R \mathcal{T}_X$ is canonical i.e. independent of coordinates.
\end{remark}

\begin{lemma}
    We have $$\mathrm{DR}(M) \cong K(M; \partial_1, \ldots, \partial_n).$$The coordinate-free de Rham complex is isomorphic to the Koszul complex on the partial derivatives.
\end{lemma}
\begin{proof}
    When $n = 2$,  $\mathrm{DR}(M):0 \to M \to \Omega^{1}_{X} \otimes M \to \Omega^{2}_{X} \otimes M \to 0.$ $$K(M; \partial_{1}, \partial_{2})= M \to M \oplus M \to M\atop m \mapsto (\partial_{1} m, \partial_{2} m)\quad(m_1, m_2) \mapsto \partial_2(m_1) - \partial_1(m_2)$$On the other hand, we have $M\rightarrow \Omega^{1}_{X}\otimes_{R} M\rightarrow \Omega_{X}^{2}\otimes_{R}M$ to be $m\mapsto dx_{1}\otimes\partial_{1} m+dx_{2}\otimes\partial_{2}m$, $dx_{1}\otimes m_{1}+dx_{2}\otimes m_{2}\mapsto dx_{1}\wedge dx_{2}\otimes(\partial_{1} m_{2}-\partial_{2} m_{1})$
	
    The general case is similar. 
\end{proof}
	
\begin{example}
\quad
\begin{enumerate}
    \item Given $\alpha \in \mathbb{C}$, define $M_{\alpha} := \mathbb{C}[x, \frac{1}{x}] \cdot x^{\alpha}$, where $x^{\alpha}$ is a formal symbol such that $\partial_x x^{\alpha} = \alpha \frac{1}{x}\cdot x^{\alpha}$, then  $M_{\alpha}$ is a $\mathcal{D}_X$-module with: $$\partial_{x} \frac{1}{x^{k}} \cdot x^{\alpha} = (\alpha - k) \cdot \frac{1}{x^{k+1}} x^{\alpha} \quad \text{for } k \gg 0$$Since we obtain $\frac{1}{x^{k}}\cdot x^{\alpha}$ as a generator of $M_{\alpha}$ as $\mathcal{D}_{X}$-module, we have $M_{\alpha} \simeq \mathcal{D}_{\alpha} /\mathrm{Ann}_{\mathcal{D}_{\alpha}}( \frac{1}{x^{k}}\cdot x^{\alpha})\simeq\mathcal{D}_X / \mathcal{D}_X \cdot (x\partial_x - (\alpha - k))$. With $\mathrm{ch}(M)=V(\sigma(P))\subseteq T^{\ast}X$ in which $\sigma$ is the principal symbol, and $\sigma(P)=x\xi$, consider $T^{\ast}X\xrightarrow{\pi} X,\quad (x,\xi)\mapsto x$, we obtain $V(x)=\pi ^{-1}(0)=T^{\ast}_{\{0\}}X$ and $V(\xi)=T^{\ast}_{X}X$ zero section.  $\mathrm{ch}(M)=V(\sigma(x\partial_{x}))=T^{\ast}_{\{0\}}X\cup T_{X}^{\ast}X$ and  $$CC(M_{\alpha}) = T_{\{0\}}^*X + T_X^*X.$$Since $\mathrm{dim}ch(M)=1$ for $\mathrm{dim}T^{\ast}_{X}X=\mathrm{dim}X=1$, and $T^{\ast}_{\{0\}}X$ is the fiber of $T^{\ast}M$ at $0$, thus the cotangent space of $X$ at $0$, $\mathrm{dim}T^{\ast}_{\{0\}}X=\mathrm{dim}X=1$ as well. So  $M_{\alpha}$ is holonomic for all $\alpha \in \mathbb{C}$.
    \item Take $N = \mathbb{C}[\delta] \simeq \mathcal{D}_X / \mathcal{D}_X \cdot x \simeq \mathbb{C}[\partial_x]$, where $\delta$ is the Dirac delta function. Then $$CC(N) = T_{\{0\}}^*X$$also be holonomic. 

\end{enumerate}

\end{example}

Now, for $\alpha \notin \mathbb{Z}$, consider the de Rham complex of $M_{\alpha}$: $$M_{\alpha} \xrightarrow{\partial_1} M_{\alpha}$$

With $\text{Ker}(\partial_1) = \{f \cdot x^{\alpha} \mid \partial_1(f \cdot x^{\alpha}) = 0\}\simeq \{f \in \mathbb{C}[x, \frac{1}{x}] \mid xf' + \alpha f= 0\} = 0$, know that $x^{-\alpha}$ should be the solution of $xf'+\alpha f$, but $\alpha\not\in \mathbb{Z}$ makes $x^{-\alpha}\not\in \mathbb{C}\left[ x, \frac{1}{x} \right]$, thus has no solution in $\mathbb{C}[x, \frac{1}{x}]$. Yet, from ODE, we know that the differential equation $xf'+\alpha f=0$ has solutions at least locally, which suggests that the current definition of de Rham complex is not enough, and we need to extend scalars. 

\quad

Let $Y = \text{Spec } S$ be an affine variety of finite type over $\mathbb{C}$ with $\dim Y = n$. Define $Y \longrightarrow \mathbb{A}^{m}_{\mathbb{C}}$ by $$ \mathbb{C}[z_1, \ldots, z_m] \longrightarrow S\atop  z_i \mapsto \text{generators of } S$$And denote $I = I(Y) = \langle f_1, \ldots, f_k \rangle, \quad f_i \in \mathbb{C}[z_1, \ldots, z_m]$. 

In particular, $f_i$ are holomorphic functions on $\mathbb{C}^m$, then $f_1 = \cdots = f_k = 0$ defines an analytic subvariety $Y^{an}$ of the complex manifold $\mathbb{C}^m$. We have a continuous map: $$ \lambda: Y^{an} \longrightarrow Y $$For all$\ y \in Y^{\text{an}}$, by Implicit Function Theorem, $Y^{\text{an}}$ is smooth around $y$ if and only if the Jacobian matrix $J(f_1, \ldots, f_k)(y)$ at $y$ has rank $m-n$. 

\begin{definition} We define

\begin{enumerate}
    \item $Y$ is $\textbf{smooth}$ at a closed point $p$ if $$\dim_{\mathbb{C}} \mathfrak{m}_p/\mathfrak{m}_p^2 = n,$$ where $\mathfrak{m}_p$ is the maximal ideal of $p$.
    \item $Y$ is $\textbf{smooth}$ if $Y$ is smooth at every closed point.
\end{enumerate}
\end{definition}
\begin{lemma}
    $Y$ is smooth $\iff$ $Y^{\text{an}}$ is a complex manifold.
\end{lemma}
\begin{proof}
    Enough to prove for closed point $p$ of $Y$ we have: $$\text{rank}(J(f_1, \ldots, f_k)(p)) = m-n \quad (*)\iff \quad \dim_{\mathbb{C}} \mathfrak{m}_{p}/\mathfrak{m}_p^2 = n.$$Know that $p$ is also a closed point of $\mathbb{A}^{m}_{\mathbb{C}}$ with maximal ideal $\mathfrak{m} = (z_1 - \alpha_1, \cdots, z_m - \alpha_m)$. Define $$\mathbb{C}[z_1, \cdots, z_m] \to V = \mathbb{C}^m\atop f \quad \mapsto \quad \left( \frac{\partial f}{\partial z_1}(p), \cdots, \frac{\partial f}{\partial z_m}(p) \right)$$By HW, this induces an iso:$$
\mathfrak{\mathfrak{m}} /\mathfrak{m}^2 \xrightarrow[\varphi]{\simeq} V\simeq \Omega_{V,p}=\Omega_{V}\otimes R/\mathfrak{m}_{p}$$Since $I\subseteq \mathfrak{m}$ is generated by $f_{1},f_{2},\dots,f_{k}$, and for $h\in \mathbb{C}[z_{1},\dots,z_{m}]$, we can write $h=c+\sum(z_{i}-\alpha_{i})+\sum(z_{i}-\alpha_{i})(z_{j}-\alpha_{j})+\dots$, then we get $[hf_{i}]=c[f]$ in $\mathfrak{m} /\mathfrak{m}^{2}$ for $\mathfrak{m}^{2}$ makes only monomial with $\mathrm{deg}\leq 1$ matter and $f=\sum(z_{i}-\alpha_{i})+(\text{higher deg})$. So we have $\varphi(f_{1}),\dots,\varphi(f_{k})$ span $\varphi(I)=I+\mathfrak{m} /\mathfrak{m}^{2}$, and further, we obtain: $$
\dim(\varphi(I)) = \operatorname{rank} \left( J(f_1, \cdots, f_k)(p) \right).
$$
On the other hand, we have short exact sequence $$0\rightarrow\varphi(I)=I+\mathfrak{m}^{2} /\mathfrak{m}^{2}\rightarrow \mathfrak{m} /\mathfrak{m}^{2}\rightarrow \mathfrak{m}/\mathfrak{m}^{2}+I\rightarrow 0.$$
Know that $\mathfrak{m}_{p}=\mathfrak{m} /I$, $\mathfrak{m}_{p} /\mathfrak{m}_{p}^{2}=\frac{\mathfrak{m}/I}{\mathfrak{m}^{2}+I /I}\simeq \mathfrak{m} /\mathfrak{m}^{2}+I$,  with $\mathrm{dim}_{\mathbb{C}} \mathfrak{m} /\mathfrak{m}^{2}=m$, we get $\operatorname{rank} \left( J(f_1, \cdots, f_k)(p) \right)+\mathrm{dim}_{\mathbb{C}}\mathfrak{m}_{p} /\mathfrak{m}_{p}^{2}=m$, done. 
\end{proof}
\quad

Let $Z$ be a complex manifold (or variety) and $\mathcal{O}_{Z}$ the sheaf of all holomorphic functions on $Z$. 

We consider $Z=Y^{an}$, have $\lambda:Y^{an}\rightarrow Y$ induces a morphism of sheaves of rings $S \longrightarrow \mathcal{O}_{Y^{an}}$. If $M$ is an $S$-module, then $$
M^{\text{an}} \coloneqq M \otimes_S \mathcal{O}_{Y^{an}}
$$Then we finally define
$$
\operatorname{DR}(M) = K(M^{\text{an}}; \partial_1, \cdots, \partial_n)
$$
And $\mathrm{DR}(M)=[0 \rightarrow M^{an} \rightarrow M^{an} \otimes_R \Omega^1_X \rightarrow \cdots \rightarrow M^{an} \otimes_R \Omega^n_X \rightarrow0]=[0 \rightarrow M^{an} \rightarrow M^{an} \otimes_{\mathcal{O}} \Omega^1_{X^{an}} \rightarrow \cdots \rightarrow M^{an} \otimes_{\mathcal{O}} \Omega^n_{X^{an}}  \rightarrow0]$. 

Now, let's calculate $\mathrm{DR}(M_{\mathfrak{\alpha}})$. 

\begin{example}
    $X = \mathbb{A}^1_{\mathbb{C}}$,  $M_{\alpha} := \mathbb{C}[x, \frac{1}{x}] \cdot x^{\alpha}$.
    \begin{enumerate} 
    \item $\alpha \not\in \mathbb{Z}$, $\mathrm{DR}(M_{\alpha}) \cong [M_{\alpha} \xrightarrow{d_{\alpha}} M_{\alpha}]$. 
$\text{Ker}\partial_{x} \mid_{x \backslash \{0\}} \\ = \mathbb{C} \cdot x^{-\alpha} \cdot x^{d}$ since $\partial_{x}(x^{-\alpha}\cdot x^{\alpha})=\partial_{x}\cdot 1=0$. This implies $\text{Ker}\partial_{x} \mid_{x \backslash \{0\}}$ is the (non-trivial) local system $L_{\lambda}$ given by the multivalued function $x^{-\alpha}$, i.e. $T \cdot x^{-\alpha} = T \cdot e^{-\alpha \log x}= e^{-\alpha(\log x + 2\pi i)} = e^{-\alpha \cdot 2\pi i} e^{-\alpha \log x}= \lambda e^{-\alpha \log x}$, here $\lambda=e^{-\alpha2\pi i}$. 
$\partial_{x} \mid_{x \backslash \{0\}}$ is surjective, and then $\mathrm{DR}(M_{\alpha}) \mid_{x \backslash \{0\}} \simeq L_{\lambda}$, and  $H^{0}\mathrm{DR}(M_{\alpha}) \mid_{\{0\}} = H^{1}DR(M_{\alpha}) \mid_{\{0\}} = 0$. 
\item $\alpha \in \mathbb{Z} \Rightarrow \lambda = 1$, $\mathrm{DR}(M_{\alpha}) \mid_{x \backslash \{0\}} = \mathbb{C}_{x \backslash \{0\}}$, $H^{0}\mathrm{DR}(M_{\alpha}) \cong H^{1}\mathrm{DR}(M_{\alpha}) \cong \mathbb{C}_{\{0\}}$
\item $N=\mathbb{C}[\delta]\simeq \mathbb{C}[\partial_{x}]$, $\mathrm{DR}(N)\simeq \mathbb{C}_{\{0\}}[-1]$. 
    \end{enumerate}
\end{example}

\section{Basic version Riemman-Hilbert correspondence}

Let $Y$ be a complex manifold, $\mathcal{O}_Y$: the sheaf of holomorphic functions on $Y$, $\mathcal{L}$: a sheaf of $\mathbb{C}$-vector spaces. 

\begin{definition}
    $\mathcal{L}$ is called a $\textbf{local system}$ if $\forall y \in Y$, $\exists$ open neighborhood $U_y$ of $y$ s.t. $\mathcal{L}|_{U_y} \simeq\mathbb{C}^r_{U_y}$, where $\mathbb{C}^{r}_{U_y}$ is the constant sheaf on $U_y$.
\end{definition}

\begin{theorem}
    Let $X = \text{Spec } R$ be a smooth affine algebraic variety over $\mathbb{C}$ and $M$ an integrable connection on $X$, then $$M^{\text{an}} \simeq \mathcal{L} \otimes \mathcal{O}_{X^{\text{an}}},$$with $\mathcal{L}$ a local system.
\end{theorem}
\begin{proof}
   \quad

    Goal: $M$ is locally free, we want to show for small enough neighborhood, we can pick a basis $s_{1},\dots,s_{r}$ makes $\partial_{x_{i}}s_{j}=0$, thus
	We have known $M$ is a locally free $R$-module, $\forall x \in X$, take an open neighborhood $U \subset X$ s.t. $M|_U$ is free, then $M^{\text{an}}|_{U^{\text{an}}}$ is a free $\mathcal{O}_{U^{\text{an}}}$-module, generated by $\underline{\mathfrak{m}} = (m_1, \cdots, m_r)$.

	Denote $(x_1, \cdots, x_n)$ coordinates of $U^{\text{an}}$, then we have $\partial_{n} \underline{\mathfrak{m}} = X \cdot\underline{ \mathfrak{m}}$, where $X$ is a matrix of holomorphic functions on $U^{an}$. Expand $X$ w.r.t. $x_n$ in a small enough neighborhood of $x\in V\subseteq U^{an}$, we have $$X = \sum_{i\ge0} x_n^i \cdot X_i,$$where $X_i$ is a matrix of holomorphic functions in $(x_1, \cdots, x_{n-1})$. 

	Take $\Psi=\sum\limits_{i\ge 0} \frac{1}{i+1} x_{n}^{i+1}X_{i}$, so $\partial_{x}(\Psi)=X$, then $\Phi=\mathrm{\exp}(-\Psi)$, which is an invertible matrix. Set $\underline{S}=(s_{1},\dots,s_{r})=\Phi(\underline{\mathfrak{m}})$, so $\partial_{x}\underline{S}=\partial_{n}(\Phi)\cdot \underline{\mathfrak{m}}+\Phi \cdot(\partial_{n} \underline{\mathfrak{m}})=(\partial_{n}(\Phi)+\Phi \cdot X)\underline{\mathfrak{m}}=0$. 

	Since $\Phi$ is invertible matrix, we obtain $s_{1},\dots,s_{r}$ also a basis of $U^{an}|_{V}$, so we have $M^{an}|_{V}=\bigoplus\limits_{i=1}^{r}\mathcal{O}_{V}s_{i}$ s.t. $\partial_{n}s_{i}=0$. 

	Want to clarify the relation of $s_{i}$ and other $\partial_{x_{j}},j<n$, then Take $M_1 = \ker[M^{\text{an}}|_V \xrightarrow{\partial_n} M^{\text{an}}|_V]=\left\{ \partial_{n}\left( \sum f_{i}s_{i} \right)=0:f_{i}\in \mathcal{O}_{V} \right\}=\left\{ \sum\limits_{i=1}^{n}\partial_{n}(f_{i})s_{i}=0 \right\}.$ So we get $\sum_{i=1}^r \partial_n(f_i) s_i = 0$ for all $i=1,\dots,r$, i.e. $f_{i}$ is independent of $x_{n}$. 

	So we have $M_{1}=\bigoplus\limits_{i=1}^{r}\mathcal{O}_{V}^{n-1}s_{i}$, $\mathcal{O}_{V}^{n-1}$ are holomorphic functions that are independent of $x_{n}$, which suggests that we do not need the var $x_{n}$. And do the same thing for $M_{1}$,we get a new basis with $\partial_{n}s_{i}=\partial_{n-1}s_{i}=0$ for all $i=1,\dots,r$, and the change of basis is via an invertible matrix of homomorphic functions in $(x_{1},\dots,x_{n-1})$. 

	Continue the process, by induction, we finally obtain:  $$ M_n = \bigoplus_{i=1}^r \mathbb{C} \cdot s_i $$with $\partial_j(s_i) = 0 \quad \forall i, j$. And since every change of basis is by an invertible matrix, we know that for a small enough neighborhood $V \subseteq U^{an}$, we get $M^{an}|_{V}=\bigoplus\limits_{i=1}^{r}\mathcal{O}_{V}s_{i}$, with $\partial_{j}(s_{i})=0$ . 
\end{proof}

\begin{corollary}
    $\mathrm{DR}(M)|_{V}\stackrel{\text{quasi-isomorphism}}{\simeq}H^{0}\mathrm{DR}(M)|_{V}\simeq \mathbb{C}^{r}_{V}$. 
\end{corollary}
\begin{example}[Algebraic vs Analytic]
    Let $M = \mathbb{C}[x, \frac{1}{x}]$, $M' = \mathbb{C}[x, \frac{1}{x}] \cdot e^{\frac{1}{x}}$. 
    
    Know that $M$ and $M'$ are different as algebraic $D$-modules, but $M^{\text{an}}|_{\mathbb{C}^*} = \mathcal{O}_{\mathbb{C}^*}$, $M'^{an}|_{\mathbb{C}^{\ast}}=\mathcal{O}_{\mathbb{C}^{\ast}}\cdot e^{1/x}=\mathcal{O}_{\mathbb{C}^{\ast}}$, where $\mathbb{C}^{\ast}$ is the punctured complex plane. So $M^{\text{an}}|_{\mathbb{C}^*} \cong M'^{\text{an}}|_{\mathbb{C}^*}$ as analytic $\mathcal{D}$-modules and $DR(M^{\text{an}}) = DR(M'^{\text{an}}) \cong \mathbb{C}_{\mathbb{C}^*}^*$ Moreover, $M^{\text{an}} = \mathcal{O}_{\mathbb{C}}\cdot \frac{1}{x},M'^{an}=\mathcal{O}_{\mathbb{C}}\cdot e^{1/x}$ so $M^{an}$ and $M'^{an}$ are different as analytic $\mathcal{D}$-module. 
\end{example}

\section{Regular holonomic $\mathcal{D}$-module}

Let $X = \text{Spec } R$ be a smooth affine algebraic variety over $\mathbb{C}$ and $M$ a holonomic $\mathcal{D}_X$-mod. 
For any point $x \in X^{\text{an}}$, we consider the embedding:$$
i: \Delta \hookrightarrow X^{\text{an}}, \quad 0 \mapsto x.
$$where $\Delta$ is a small disk in the complex plane.
Then we have a morphism of sheaves of rings:$$R \to \mathcal{O}_{\Delta}, \quad f \mapsto f|_{\Delta}.$$We define $i^*M = M \otimes_R \mathcal{O}_{\Delta}$. 
\begin{definition}
    $M$ is called $\textbf{regular holonomic}$, if $i^{\ast}M$ is generated by $M_{\alpha}^{an}$ and $N^{an}|_{\Delta}$ in $\mathbf{K}_{0}(\mathcal{D}_{\Delta})$ for every small enough disk embedding $i:\Delta \hookrightarrow X^{an}$. 
\end{definition}

\begin{example}
    $M=\mathbb{C}\left[ x, \frac{1}{x} \right]\cdot e^{1/x}$ is not regular because $M$ is simple and $M\simeq \mathcal{D}_{X} /\mathcal{D}_{X}(x^{2}\partial_{x}+1)$. 
\end{example}
\chapter{Auslander regularity and more on filtered rings$\&$modules}
\section{Auslander regular ring}
Let $A$ be a noetherian ring (not necessarily commutative), and $\mathrm{Mod}_{f}(A)$, $\mathrm{Mod}_{f}^{op}(A)$ the abelian category of left/right finitely generated $A$-module.

For $M \in \mathrm{Mod}_{f}(A)$, the $\textbf{projective dimension}$ of $M$ is
$$\mathrm{pd}(M) = \text{minimum of length of projective resolution.}$$

We say $A$ is $\textbf{of finite global dimension}$, if $\exists w \in \mathbb{Z}_{>0}$ s.t for all $M \in \mathrm{Mod}_{f}(A)$, have $\mathrm{pd}(M) \le w$.

We now assume that $A$ is finite type over $k$ and of finite global dimension. Then $\mathrm{RHom}(M,A) \in \mathbf{D}^{b}(\mathrm{Mod}^{op}_{f}(A))$, the derived category of bounded complexes in $\mathrm{Mod}_{f}^{op}(A)$.

\begin{definition}
    \quad
    \begin{enumerate}
        \item The $\textbf{grade number}$ $j(M) := \min\left\{k \middle| \mathrm{Ext}^{k}(M,A) \neq 0\right\}$, we sometimes denote it as  $j_{A}(M)$ as well. In particular, define $j(M) = +\infty$ if $M = 0$.
        \item For a noetherian and of finite global dimension ring $A$, we say $A$ is $\textbf{Auslander regular}$ (AR for short) if for all $M \in \mathrm{Mod}_{f}(A)$ and $v \in \mathbb{Z}_{\ge 0}$, given any $N \subseteq \mathrm{Ext}^{v}(M,A)$ submodule, have $j(N) \ge v$.
    \end{enumerate}
\end{definition}

Let $S$ be a commutative $k$-algebra of finite type, we say $S$ is $\textbf{regular}$ if $\Omega_{S/k}$ is locally free, see definition \cite[\href{https://stacks.math.columbia.edu/tag/00OD}{Tag 00OD}]{stacks-project}.

\begin{theorem}[Auslander]\label{theorem 7.2}
    Let $S$ be regular of dimension $m$, $M \in \mathrm{Mod}_{f}(S)$, then
    \begin{enumerate}
        \item $S$ is of finite global dimension;
        \item $j(M) + \dim(\mathrm{supp}(M)) = m$;
        \item $\dim(\mathrm{supp}(\mathrm{Ext}^{i}(M,S))) \le m - i$ for all $i$;
        \item $\dim(\mathrm{supp}(\mathrm{Ext}^{j(M)}(M,S))) = \dim(\mathrm{supp}(M))$.
    \end{enumerate}
\end{theorem}
\begin{proof}
    \mbox{}

    For (1), see \cite[Thm.~19.12.]{eisenbud1995commutative} or \cite[\href{https://stacks.math.columbia.edu/tag/00OC}{Tag 00OC}]{stacks-project} for local ring case.
    Then use \cite[\href{https://stacks.math.columbia.edu/tag/00O9}{Tag 00O9}]{stacks-project} to get global case, which is \cite[\href{https://stacks.math.columbia.edu/tag/00OE}{Tag 00OE}]{stacks-project}.

    For (2), (3) and (4), in the case of $m$ is even, see \cite[Prop.~6.4.]{schnell_dmodules_notes}.

    For general $m$, one can see \cite[Thm.~D.4.4.]{hotta2008d}, but this book refers other books for the proof of (2), (3) and (4).
\end{proof}
\begin{corollary}
    Regular commutative $k$-algebra of finite type are Auslander regular.
\end{corollary}

Let $X = \mathrm{Spec } R$ be a smooth-affine variety over $\mathbb{C}$, our next goal is to prove $\mathcal{D}_{X}$ is AR.

\section{Spectral sequences of filtered complex}

Let $F^{\bullet}$ be a bounded complex of $\mathbb{Z}$-modules:$$\cdots \to F^{i-1} \xrightarrow{\mathrm{d}} F^i \xrightarrow{\mathrm{d}} F^{i+1} \to \cdots$$with an increasing $\mathbb{Z}$-filtration: $F_j^i \subseteq F_{j+1}^i \subseteq F^i \quad (\forall j,i \in \mathbb{Z})$ satisfying:
\begin{enumerate}
    \item Exhaustive condition $\bigcup\limits_j F_j^i = F^i \quad \forall i$ ;
    \item Hausdorff condition $\bigcap_j F_j^i = 0 \quad \forall i$
    \item  Differential preserves filtration $\mathrm{d}(F_j^i) \subseteq F_j^{i+1} \, (\forall i,j)$ , which implies that $\dots\rightarrow F_{j}^{i-1}\xrightarrow{\mathrm{d}} F_{j}^{i}\xrightarrow{\mathrm{d}} F_{j}^{i+1}\rightarrow\dots$ is a complex.
\end{enumerate}
Then, to every triple $(k, i, j)$ we set:
\begin{itemize}
    \item $Z_j^i(k) = \{u \in F_j^i \mid du \in F^{i+1}_{j-k}\}.$
    \item $B^{i}_{j}(k)=F^{i}_{j}\cap \mathrm{d}(F^{i-1}_{j+k-1})=\mathrm{d}(Z^{i-1}_{j+k-1}(k-1))$
    \item $E_j^i(k) = \dfrac{Z_j^i(k) + F_{j-1}^i}{B_j^i(k) + F_{j-1}^i}$, $E_{k}^{-j,i+j}:=E_{j}^{i}(k)$, and $E_{\bullet}^i(k) = \bigoplus\limits_{j} E_j^i(k)$.
\end{itemize}

\begin{remark}
    Since we consider the increasing filtration $\dots\subseteq F_{j}\subseteq F_{j+1}\subseteq \dots \subseteq F$, while usually we consider a decreasing filtration for cochain complex, we can actually view $F_{j}$ as $F^{-j}$. That's the reason why we define $E_{k}^{-j,i+j}:=E_{j}^{i}(k)$.
\end{remark}
So $\mathrm{d}(Z_j^i(k)) \subseteq Z_{j-k}^{i+1}(k)$, which induces $$ E_k^{-j,i+j} := E_j^i(k) \xrightarrow{\mathrm{d}_{k}} E_k^{-j+k,i+j-k+1} = E_{j-k}^{i+1}(k) .$$Further, we have $$\dots\xrightarrow{\mathrm{d}_{k}}E_{j+k}^{i-1}(k)\xrightarrow{\mathrm{d}_{k}}E_{j}^{i}(k)\xrightarrow{\mathrm{d}_{k}}E_{j-k}^{i+1}(k)\xrightarrow{\mathrm{d}_{k}} \dots$$Consider $E^{i}_{\bullet}(k)=\bigoplus\limits_{j}E^{i}_{j}(k)$ and know that $\mathrm{d}_{k}:E^{i}_{\bullet}(k)\rightarrow E^{i+1}_{\bullet}(k)$, this makes $E_\bullet(k)$ a complex of graded $\mathbb{Z}$-modules.  And $E_{j}^{i}(0) = F^{i}_{j} /F_{j-1}^{i}$, i.e. $E^{i}_{\bullet}(0)=\mathrm{gr }_{\bullet}  F^{i}$ the associated graded complex of  $F^{i}$ for all $i$.

\medskip

$\{E_k^{-\ast,\bullet+\ast} = E_\ast^\bullet(k)\}_k$ is called the $\textbf{spectral sequence}$ of the filtered complex $F^{\bullet}$.

\medskip

\begin{lemma}\label{lemma 7.4}
    $E_\bullet^i(k+1) = H^i(E_{\bullet}^{\bullet}(k))$, more explicitly: $$E^{i}_{j}(k+1)\simeq \frac{\ker(E_{j}^{i}(k)\xrightarrow{\mathrm{d}_{k}}E_{j-k}^{i+1}(k) )}{\mathrm{Im}(E_{j+k}^{i-1}(k)\xrightarrow{\mathrm{d}_{k}}E_{j}^{i}(k))}.$$
\end{lemma}
\begin{proof}
    \quad

    (After verifying that $\mathrm{d}_{k}$ is well-defined) Consider $k=0$, we obtain that $E^{i}_{j}(0)= F^{i}_{j} /F^{i}_{j-1}$, so
    $$\mathrm{ker}(E^{i}_{j }(0)\rightarrow E^{i+1}_{j}(0))=\ker(F^{i}_{j} /F^{i}_{j-1}\rightarrow F^{i+1}_{j} /F^{i+1}_{j-1})=\left(\{u\in F^{i}_{j}:du\in F^{i+1}_{j-1}\}+F^{i}_{j-1}\right) /F^{i}_{j-1},$$
    and
    $$\mathrm{Im}(E^{i-1}_{j}(0)\rightarrow E^{i}_{j}(0))=\mathrm{Im}(F^{i-1}_{j} /F^{i-1}_{j-1}\rightarrow F^{i}_{j} /F^{i}_{j-1})=\mathrm{d}(F^{i-1}_{j})+F^{i}_{j-1} /F^{i}_{j-1},$$
    then we get $$E^{i}_{j}(1)\simeq \frac{\mathrm{ker}(E^{i}_{j }(0)\rightarrow E^{i+1}_{j}(0))}{\mathrm{Im}(E^{i-1}_{j}(0)\rightarrow E^{i}_{j}(0))}.$$

    For $k\ge 1$, we obtain that $F_{j-1}^{i}\subseteq F_{j+k-1}^{i}$$$E^{i}_{j}(k+1)=\frac{\{u\in F^{i}_{j}:du\in F^{i+1}_{j-k-1}\}+F^{i}_{j-1}}{F^{i}_{j}\cap \mathrm{d}(F^{i-1}_{j+k})+F^{i}_{j-1}}.$$

    For $\left(Z^{i}_{j}(k)+F^{i}_{j-1}\right) /F^{i}_{j-1}\xrightarrow{\mathrm{d}_{k}}\left(Z_{j-k}^{i+1}(k)+F^{i+1}_{j-k-1}\right) /F^{i+1}_{j-k-1}$,
    $$\ker \mathrm{d}_{k}:=\{u\in F^{i}_{j} /F_{j-1}^{i}: \exists x\in u\ s.t.\, \mathrm{d}x\in F^{i+1}_{j-k-1}\}=\left(Z^{i}_{j}(k+1)+F^{i}_{j-1}\right) /F^{i}_{j -1}$$
    And
    $$\mathrm{Im}\left(\left(Z^{i-1}_{j+k}(k)+F^{i-1}_{j+k-1}\right) /F^{i-1}_{j+k-1}\xrightarrow{\mathrm{d}_{k}}\left(Z^{i}_{j}(k)+F^{i}_{j-1}(k)\right) /F^{i}_{j-1}(k)\right)=\left(\mathrm{d}(Z^{i-1}_{j+k}(k))+F^{i}_{j-1}(k)\right) /F^{i}_{j-1}(k).$$
\end{proof}

\quad

\begin{definition}[Induced filtration on Cohomology]\label{def:induced_filtration_cohomology}
    The filtration of $F^{\bullet}$ induces a filtration of $H^i(F^{\bullet})$ by $$F_j H^i(F^\bullet) = \frac{Z_j^{i}(\infty) + \mathrm{d}F^{i-1}}{\mathrm{d}F^{i-1}}.$$where $Z_j^{i}(\infty) := \mathrm{Ker}(d) \cap F_j^i$.
\end{definition}
\begin{remark}
    For standard defintion, see \cite[Section~5.5]{Weibel1994HA}, or see Remark~\ref{rmk:induced-filtration-on-Ext} we proved two definition are equivalent there.
\end{remark}

In fact, we often view $Z^{i}_{j}(k)$ and $B^{i}_{j}(k)$ in $F^{i}_{j} /F^{i}_{j-1}$. Namely, we define
$$\eta^{i}_{j}:F^{i}_{j}\mapsto F^{i}_{j} /F^{i}_{j-1},$$
and
$$\overline{Z}^{i}_{j}(k):=\eta^{i}_{j}(Z^{i}_{j}(k)),$$
$$\overline{B}^{i}_{j}(k):=\eta^{i}_{j}(B^{i}_{j}(k)),$$
we obtain that $E^{i}_{j}(k)\simeq \frac{\overline{Z}^{i}_{j}(k)}{\overline{B}^{i}_{j}(k)}$. Moreover, we have a tower for all $i,j:$ $$
    0 \subseteq \overline{B}_j^i(0) \subseteq \overline{B}_j^i(1) \subseteq \cdots \subseteq \overline{B}_j^i(r) \subseteq \cdots \subseteq \overline{B}_j^i(\infty)
    \subseteq \overline{Z}_j^i(\infty) \subseteq \cdots \subseteq \overline{Z}_j^i(r) \subseteq \overline{Z}_j^i(r-1) \subseteq \cdots \subseteq \overline{Z}_j^i(0)
$$
We define the limit terms:
\[
    \overline{B}_j^i(\infty) := \varinjlim_{r} \overline{B}_j^i(r) \quad \text{and} \quad \overline{Z}_j^i(\infty) := \varprojlim_{r} \overline{Z}_j^i(r).
\]

\begin{lemma}
    Assuming the filtration is exhaustive and Hausdorff, we have:
    \[
        \overline{B}_j^i(\infty) = \eta_j^i(F_j^i \cap \mathrm{d}(F^{i-1}))
        \quad \text{and} \quad \overline{Z}_j^i(\infty)
        \, \textcolor{red}{\supseteq}\, \eta_j^i(Z_j^i(\infty)).
    \]
\end{lemma}
\begin{proof}
    For the first equality, recall that $\overline{B}_j^i(r) = \eta_j^i(F_j^i \cap \mathrm{d}(F_{j+r-1}^{i-1}))$.
    Since the filtration is exhaustive, we have $\bigcup_{r} F_{j+r-1}^{i-1} = F^{i-1}$.
    Thus $\bigcup_{r} \mathrm{d}(F_{j+r-1}^{i-1}) = \mathrm{d}(F^{i-1})$.
    Therefore, $\varinjlim_{r} \overline{B}_j^i(r) = \eta_j^i(F_j^i \cap \mathrm{d}(F^{i-1}))$.

    For the second formula, Hausdorff condition is NOT sufficient to ensure:
    \[
        \bigcap_{r}\eta_j^i(\{u \in F_j^i \mid \mathrm{d}u \in F_{j-r}^{i+1}\}) = \eta_j^i\left(\bigcap_{r}\{u \in F_j^i \mid \mathrm{d}u \in F_{j-r}^{i+1}\}\right)
    \]
    So we can only get an inclusion, not an equality.
\end{proof}
\begin{remark}
    See \cite[the discussion above Lemma~5.5.7.]{Weibel1994HA}.
\end{remark}

\begin{proposition}\label{prop:Boardman's_convergence}
    $\mathrm{gr}_{j}^{F}H^{i}(F)\simeq \frac{\eta_j^i(Z_j^i(\infty))}{\overline{B}^{i}_{j}(\infty)}$.
\end{proposition}

\begin{proof}
    See \cite[Lemma~5.5.7.(3)]{Weibel1994HA}, where $e^\infty_{pq}$ is just $\frac{\eta_p^q(Z_p^q(\infty))}{\overline{B}^{p}_{q}(\infty)}$ \cite[Lemma~5.5.7.(1)]{Weibel1994HA}, or see the proof below:

    Since we have defined induced filtration on cohomology, then
    \begin{align*}
        \mathrm{gr}_{j}^{F}H^{i}(F) & \simeq  F^j H^i(F^\bullet) \big/ F^{j-1} H^i(F^\bullet)                                                                                                                  \\
                                    & \simeq \left(\frac{Z_j^{i}(\infty) + \mathrm{d}F^{i-1}}{\mathrm{d}F^{i-1}} \right) \big/ \left(\frac{Z_{j-1}^{i}(\infty) + \mathrm{d}F^{i-1}}{\mathrm{d}F^{i-1}} \right) \\
                                    & \simeq \frac{Z_j^{i}(\infty) + \mathrm{d}F^{i-1}}{Z_{j-1}^{i}(\infty) + \mathrm{d}F^{i-1}}                                                                               \\
                                    & \simeq \frac{Z_j^{i}(\infty)}{Z_j^{i}(\infty) \cap (Z_{j-1}^{i}(\infty) + \mathrm{d}F^{i-1})}                                                                            \\
                                    & \simeq \frac{Z_j^{i}(\infty)}{Z_{j-1}^{i}(\infty) + (Z_j^{i}(\infty) \cap \mathrm{d}F^{i-1})}                                                                            \\
                                    & \simeq \frac{Z_j^{i}(\infty)}{Z_{j-1}^{i}(\infty) + (F_j^{i} \cap \mathrm{d}F^{i-1})}.
    \end{align*}
    The last three steps are obtained as follows:
    \begin{itemize}
        \item The third isomorphism uses the Second Isomorphism Theorem $(A+B)/B \cong A/(A \cap B)$ by setting $A=Z_j^i(\infty)$ and $B=Z_{j-1}^i(\infty) + \mathrm{d}F^{i-1}$.
        \item The fourth isomorphism uses the Modular Law $A \cap (B+C) = (A \cap B) + C$ since $C=Z_{j-1}^i(\infty) \subseteq Z_j^i(\infty)=A$.
        \item The last isomorphism holds because $Z_j^i(\infty) \cap \mathrm{d}F^{i-1} = (\Ker(\mathrm{d}) \cap F_j^i) \cap \mathrm{d}F^{i-1} = F_j^i \cap (\Ker(\mathrm{d}) \cap \mathrm{d}F^{i-1}) = F_j^i \cap \mathrm{d}F^{i-1}$ noting that $\mathrm{Im}(\mathrm{d}) \subseteq \Ker(\mathrm{d})$.
    \end{itemize}
    On the other hand,
    \begin{align*}
        \frac{\eta_j^i(Z_j^i(\infty))}{\overline{B}^{i}_{j}(\infty)} & \simeq \frac{Z_j^i(\infty) + F_{j-1}^i}{F_j^i \cap \mathrm{d}F^{i-1} + F_{j-1}^i}          \\
                                                                     & \simeq \frac{Z_j^i(\infty)}{Z_j^i(\infty) \cap (F_j^i \cap \mathrm{d}F^{i-1} + F_{j-1}^i)} \\
                                                                     & \simeq \frac{Z_j^i(\infty)}{F_j^i \cap \mathrm{d}F^{i-1} + (Z_j^i(\infty) \cap F_{j-1}^i)} \\
                                                                     & \simeq \frac{Z_j^i(\infty)}{F_j^i \cap \mathrm{d}F^{i-1} + Z_{j-1}^i(\infty)}.
    \end{align*}
    Thus $\mathrm{gr}_{j}^{F}H^{i}(F) \simeq \frac{\eta_j^i(Z_j^i(\infty))}{\overline{B}^{i}_{j}(\infty)}$.
\end{proof}

\quad

\begin{definition}\label{def:convergent_spectral_sequence}
    We say the spectral sequence $\{E_{\bullet}^{\bullet}(k)\}_k$ $\textbf{converges to}$ $H^i(F^{\bullet})$ if for all $i,j$, there exists $k = k(i,j)$ s.t. $\forall k' \ge k$, $$E_j^i(k') \cong \mathrm{gr}_j^F H^i(F^{\bullet}).$$
\end{definition}

\begin{definition}$\textbf{The regularity condition}:$
    if there exists a $\omega \in \mathbb{Z}_{\ge0}$ s.t.$$F_j^i \cap \mathrm{d}(F^{i-1}_{}) \subset \mathrm{d}(F_{j+\omega}^{i-1})$$for every pair $(i,j)$, then we say the filtered complex is $\omega$-regular.
\end{definition}

\begin{lemma}\label{lemma 7.8}
    If filtered complex $F^{\bullet}$ is $\omega$-regular and convergent in \cite[Convergence~5.2.11]{Weibel1994HA}'s sense, i.e. $\mathrm{gr}_{j}^{F}H^{i}(F) \simeq \frac{\overline{Z}_j^i(\infty)}{\overline{B}^{i}_{j}(\infty)}$. Then $$\mathrm{gr}^F_{\bullet} H^{i}(F^\bullet) \simeq E_{\bullet}^{i}(\omega+1), \quad \forall i$$In particular, the spectral sequence converges.
\end{lemma}

\begin{proof}
    First, we show that $E^{i}_{j}(\omega+k)\simeq E^{i}_{j}(\omega+1)$ for all $k\ge 1$.
    It suffices to show that the differentials $\mathrm{d}_r: E_{j+r}^{i-1}(r) \to E_j^i(r)$ are zero for all $r \ge \omega+1$.

    Any element in the image of $\mathrm{d}_r$ is represented by some $v \in F_j^i \cap \mathrm{d}(F_{j+r}^{i-1})$.
    By the regularity condition, we have
    $$ v \in F_j^i \cap \mathrm{d}(F^{i-1}) \subseteq \mathrm{d}(F_{j+\omega}^{i-1}). $$
    Since $r \ge \omega+1$, we have $j+\omega \le j+r-1$, so $F_{j+\omega}^{i-1} \subseteq F_{j+r-1}^{i-1}$ and hence $\mathrm{d}(F_{j+\omega}^{i-1}) \subseteq \mathrm{d}(F_{j+r-1}^{i-1})$.
    Thus $v \in F_j^i \cap \mathrm{d}(F_{j+r-1}^{i-1}) = B_j^i(r)$, which means the class of $v$ is zero in $E_j^i(r)$.
    Therefore $\mathrm{d}_r = 0$ for all $r \ge \omega+1$, and the spectral sequence collapses at the $(\omega+1)$-page; i.e., $E^{i}_{j}(\omega+k)\simeq E^{i}_{j}(\omega+1)$ for all $k\ge 1$.

    By Proposition~\ref{prop:Boardman's_convergence}, we obtain that $\frac{\overline{Z}^{i}_{j}(\infty)}{\overline{B}^{i}_{j}(\infty)}\simeq \mathrm{gr}_{\bullet}^{F}H^{i}(F)$, take $r\gg 0$, we consider the following map:

    \begin{center}
        \begin{tikzcd}
            {\frac{\overline{Z}_j^i(r)}{\overline{B}_j^i(r)}} & {\frac{\overline{Z}_j^i(\infty)}{\overline{B}_j^i(r)}} & {\frac{\overline{Z}_j^i(\infty)}{\overline{B}_j^i(\infty)}} \\
            & {\frac{\overline{Z}_j^{i+1}(r+k)}{\overline{B}_j^{i+1}(r+k)}}
            \arrow["\simeq", from=1-1, to=2-2, sloped, dash]
            \arrow["\Phi"', hook', from=1-2, to=1-1]
            \arrow["\Psi", two heads, from=1-2, to=1-3]
            \arrow[from=1-2, to=2-2]
        \end{tikzcd}
    \end{center}

    1. Since $\mathrm{d}_k = 0$ for all $k \ge r$, we have $\frac{\overline{Z}_j^i(k)}{\overline{B}_j^i(k)} = \frac{\overline{Z}_j^i(k+1)}{\overline{B}_j^i(k+1)}$ for $k \ge r$.
    Since we have tower
    \[
        0 \subseteq \overline{B}_j^i(0) \subseteq \overline{B}_j^i(1) \subseteq \cdots \subseteq \overline{B}_j^i(r) \subseteq \cdots \subseteq \overline{B}_j^i(\infty)
        \subseteq \overline{Z}_j^i(\infty) \subseteq \cdots \subseteq \overline{Z}_j^i(r) \subseteq \overline{Z}_j^i(r-1) \subseteq \cdots \subseteq \overline{Z}_j^i(0)
    \]
    This implies the filtration on $Z$ and $B$ stabilize:
    $$ \overline{Z}_j^i(r) = \overline{Z}_j^i(r+1) = \cdots = \overline{Z}_j^i(\infty) $$
    and
    $$ \overline{B}_j^i(r) = \overline{B}_j^i(r+1) = \cdots = \overline{B}_j^i(\infty). $$
    Thus the map $\Phi$ (induced by inclusion $\overline{Z}_j^i(\infty) \subseteq \overline{Z}_j^i(r)$) is an isomorphism because the inclusion is actually an equality.

    2. Similarly, the map $\Psi$ (induced by projection modulo $\overline{B}_j^i(\infty)$ vs $\overline{B}_j^i(r)$) is an isomorphism because the denominators are equal.

    Therefore, $E_j^i(k)\simeq\frac{\overline{Z}^{i}_{j}(\infty)}{\overline{B}^{i}_{j}(\infty)} \simeq \mathrm{gr}_j H^i(F) $ for all $j$.
\end{proof}

\begin{remark}\label{rmk:classical_convergence_theorem}
    We introduce the \textbf{Classical Convergence Theorem} \cite[Classical~Convergence~Theorem~5.5.1]{Weibel1994HA} for a complex with \textbf{bounded below filtration} in the sense of convergence introduced in \cite[Convergence~5.2.11]{Weibel1994HA}.

    We will use this theorem in accompany by Lemma~\ref{lemma 7.8}.
\end{remark}

\begin{theorem}\label{tehorem 7.9}
    If $F^\bullet$ is $\omega$-regular and convergent in \cite[Convergence~5.2.11]{Weibel1994HA}'s sense, then $\mathrm{gr}^F_{\bullet} H^{i}(F^{\bullet})$ is a subquotient of $H^{i} \mathrm{gr}^{F}_{\bullet }(F^{\bullet})$.
\end{theorem}
\begin{proof}
    By Lemma \ref{lemma 7.4} + Lemma \ref{lemma 7.8}.
\end{proof}

\section{More on filtered rings }\label{sect:more_on_filtered_rings}

Recall that we have defined a $\mathbb{Z}$-filtration of a ring in Definition \ref{def:z_filtration}.

Let $A$ be a ring and $\Gamma_{\bullet} = \Gamma_{\bullet}A$ a $\mathbb{Z}$-filtration on $A$, recall that $\Gamma_{\bullet}A$ is called a $\textbf{noetherian } \mathbb{Z}\textbf{-filtration}$ if additionally $\mathrm{gr}^{\Gamma}_\bullet A$ is also noetherian.

From now on, we assume that $A$ is a $\C$-algebra and has a noetherian $\mathbb{Z}$-filtration $\Gamma_{\bullet}$, so it's rees algebra $R_{\bullet}^{\Gamma}A$ is a $\C[u]$ algebra.

Let $M$ be a finitely-generated $A$-module, and $\Omega_{\bullet}M$ a good filtration on $M$ s.t. $\mathrm{gr}^{\Omega} M \neq 0$.
Since $\mathrm{gr}^{\Gamma} A$ is noetherian, $R^{\Omega} M$ has a graded free resolution:
$$ \dots\rightarrow P_{\bullet}^{-1}\rightarrow P^{0}_{\bullet}\rightarrow R^{\Omega} _{\bullet}M\rightarrow0$$
s.t. each $P^{i}_{\bullet}$ is a free graded $R_{\bullet}^{\Gamma}A$-module of finite rank.
So in derived category, $\mathrm{RHom}_{R_{\bullet}A}(R_{\bullet}M, R_{\bullet}A)$ is the complex \cite[\href{https://stacks.math.columbia.edu/tag/0A8H}{Tag 0A8H}, \href{https://stacks.math.columbia.edu/tag/0A66}{Tag 0A66}]{stacks-project} (well-defined up to quasi-isomorphism in $\mathbf{D}(R_{\bullet}A)$):
$$\mathrm{Hom}_{R_{\bullet}A}(P^{\bullet}_{\bullet}, R_{\bullet}A).$$
Tensoring with $\mathbb{C}[u]/(u-1)$ over $\mathbb{C}[u]$, we get a free resolution:
$$ \cdots \to F^{-1} \to F^0 \to M \to 0$$
with $F^i = P^i_{\bullet} \otimes_{\mathbb{C}[u]} \mathbb{C}[u]/(u-1)$.
Define the filtration on $F^i_{\bullet}$ by:$$F_j^i := \mathrm{Im}(P_j^i \longrightarrow F^i)$$
This filtered complex:
$$\cdots \longrightarrow F_{\bullet}^{-1} \longrightarrow F_{\bullet}^0 \longrightarrow 0$$
is called a $\textbf{filtered free resolution}$ of $M$, which is $\C[u]/(u-1)\otimes^{L}_{\C[u]} R^{\Omega}M = A\otimes^{L}_{R_{\bullet}^{\Gamma}A} R^{\Omega}M$ in derived category, \cite[\href{https://stacks.math.columbia.edu/tag/0FNB}{Tag 0FNB}, \href{https://stacks.math.columbia.edu/tag/064M}{Tag 064M}]{stacks-project}
so it is well-defined up to quasi-isomorphism in $\mathbf{D}(A)$.

\bigskip

Since $\mathrm{Hom}({-},A)$ functor is contravariant, we define a filtered complex $$ \check{F}^{\bullet} = \operatorname{Hom}_{A}(F^{\bullet}, A) \quad \mathrm{s.t.}\quad \check{F}^{i} = \operatorname{Hom}_{A}(F^{-i}, A) $$
We consider $\mathrm{Hom}_{R_{\bullet}A}(P^{\bullet}_{\bullet},R_{\bullet}A)\otimes \mathbb{C}[u] /(u-1)$,
since $P^\bullet_{\bullet}$ is a free $R_{\bullet}A$-module,
we then have
$$\mathrm{Hom}_{R_{\bullet}A}(P^\bullet_{\bullet},R_{\bullet}A)\otimes \mathbb{C}[u]/(u-1)\simeq \mathrm{Hom}_{A}(F^{\bullet},A) \quad (\text{quasi-isomorphism in } \mathbf{D}(A)),$$
and filtration on $\check{F}^{\bullet}$ is defined as follows.

Recall that for two graded modules $M_{\bullet}, N_{\bullet}$ over a graded ring $R_{\bullet}$, the group of graded morphisms of degree $k$ is
(see \cite[\href{https://stacks.math.columbia.edu/tag/09MN}{Tag 09MN}]{stacks-project})
$$\mathrm{Hom}_{R_{\bullet}}(M_{\bullet}, N_{\bullet})_k := \{ f: M_{\bullet} \to N_{\bullet} \mid f(M_j) \subseteq N_{j+k}, \forall j \}.$$
Let $\mathcal{H}^i_{\bullet} := \mathrm{Hom}_{R_{\bullet}A}(P^{-i}_{\bullet}, R_{\bullet}A)$, which is a graded $R_{\bullet}A$-module. The isomorphism above induces a map from each graded piece:
$$ \Phi_k: \mathcal{H}^i_k \longrightarrow \mathrm{Hom}_A(F^{-i}, A). $$
We define the filtration on $\check{F}^i$ by the image of this map:
$$ \check{F}^i_k := \mathrm{Im}(\Phi_k). $$

\begin{proposition}\label{prop:filtration-on-Hom}
    The filtration defined above coincides with the standard induced filtration on $\mathrm{Hom}_A(F^{-i}, A)$, i.e.,
    $$ \check{F}^i_k = \{ f \in \mathrm{Hom}_A(F^{-i}, A) \mid f(F_j^{-i}) \subseteq F_{j+k}A, \forall j \}. $$
    (cf. \cite[Appendix~D, after Lemma~D.2.2]{hotta2008d})
\end{proposition}
\begin{proof}
    It suffices to check this for $P^{-i} = R_{\bullet}A(p)$, a shifted free module.
    Then $F^{-i} \cong A$ with filtration shifted by $p$, i.e., $F_j^{-i} = F_{j+p}A$.

    LHS: $\mathcal{H}^i_k = \mathrm{Hom}_{R_{\bullet}A}(R_{\bullet}A(p), R_{\bullet}A)_k \cong R_{\bullet}A_{k-p}$.
    Its image in $\mathrm{Hom}_A(A, A) \cong A$ is exactly $F_{k-p}A$.

    RHS: $\{ f: A \to A \mid f(F_j^{-i}) \subseteq F_{j+k}A \} = \{ a \in A \mid a \cdot F_{j+p}A \subseteq F_{j+k}A, \forall j \}$.
    Since $A$ has a filtration, this condition is equivalent to $a \in F_{k-p}A$ (by definition of filtration on ring).

    Thus LHS = RHS. Since $P^{-i}$ is a direct sum of such terms, the result holds.
\end{proof}

\begin{lemma}\label{lem:checkF-regular}
    $\check{F}^{\bullet}$ is  $\omega$-regular for some $\omega \gg 0.$
\end{lemma}
\begin{proof}
    Take $N = \mathrm{d}(\check{F}^{k}) \subseteq \check{F}^{k+1}$, then $N_j = \mathrm{d}(\check{F}_{j}^{k})$ gives a good filtration on $N$.
    On the other hand, $N \cap \check{F}_{j}^{k+1}$ gives another good filtration. By Corollary \ref{cor:compare-good-filtrations}, there exists a uniform $\omega$ s.t. $$ \check{F}_{j}^{k+1} \cap \mathrm{d}(\check{F}^{k}) = N \cap \check{F}_{j}^{k+1} \subseteq N_{j+\omega} = \mathrm{d}(\check{F}_{j+\omega}^{k}), \quad \forall j. $$
\end{proof}

\begin{corollary}\label{corollary 7.11}
    $\mathrm{gr}_{\bullet}\mathrm{Ext}_A^j (M, A)$ is a subquotient of $\mathrm{Ext}_{\mathrm{gr}_{\bullet} A}^j (\mathrm{gr} _{\bullet}M, \mathrm{gr}_{\bullet} A)$ for all $j$.
\end{corollary}
\begin{proof}
    Applying Theorem \ref{tehorem 7.9} to the filtered complex $\check{F}^{\bullet}$.
    \begin{itemize}
        \item By Lemma \ref{lem:checkF-regular}, $\check{F}^{\bullet}$ is $\omega$-regular.
        \item Since the filtration on $M$ is a good filtration, it is bounded from below (see \cite[Lemma~D.2.3]{hotta2008d}). This implies the filtration on $\check{F}^\bullet$ is also bounded from below, so the spectral sequence converges (see Remark~\ref{rmk:classical_convergence_theorem}).
    \end{itemize}
    Thus $\mathrm{gr}_{\bullet}H^j(\check{F}^{\bullet})$ is a subquotient of $H^j(\mathrm{gr}_{\bullet}\check{F}^{\bullet})$.

    \medskip

    Now we identify the terms:
    \begin{enumerate}
        \item By definition, $H^j(\check{F}^{\bullet}) = \mathrm{Ext}_A^j(M, A)$. Thus $\mathrm{gr}_{\bullet}H^j(\check{F}^{\bullet}) = \mathrm{gr}_{\bullet}\mathrm{Ext}_A^j(M, A)$.
        \item Since $F^\bullet$ comes from a graded free resolution $P^\bullet$, $\mathrm{gr}_{\bullet} F^{\bullet} \cong P^\bullet / u P^\bullet$ is a free resolution of $\mathrm{gr}_{\bullet} M \cong R^\Omega M / u R^\Omega M$ over $\mathrm{gr}_{\bullet} A \cong R^\Gamma A / u R^\Gamma A$.
        \item Since $F^k$ are free modules, taking the associated graded commutes with $\mathrm{Hom}$, i.e.,
              \[
                  \mathrm{gr}_{\bullet}\check{F}^{\bullet} = \mathrm{gr}_{\bullet}\mathrm{Hom}_A(F^{\bullet}, A) \text{ and } \mathcal{G}^{\bullet} = \mathrm{Hom}_{\mathrm{gr}_{\bullet}A}(\mathrm{gr}_{\bullet}F^{\bullet}, \mathrm{gr}_{\bullet}A).
              \]
              We have a canonical map of graded modules:
              \[ \Psi: \mathrm{gr}_{\bullet}\mathrm{Hom}_A(F^{\bullet}, A) \longrightarrow \mathrm{Hom}_{\mathrm{gr}_{\bullet}A}(\mathrm{gr}_{\bullet}F^{\bullet}, \mathrm{gr}_{\bullet}A). \]
              It is defined by the principal symbol: for $f \in \check{F}^k_p = F_p\mathrm{Hom}_A(F^k, A)$, its image $\Psi(f \bmod \check{F}^k_{p-1})$ is the graded morphism:
              \[
                  \mathrm{gr}_{\bullet}F^k \ni (x \bmod F_{j-1}F^k) \longmapsto (f(x) \bmod F_{j+p-1}A) \in \mathrm{gr}_{\bullet}A.
              \]
              Since $F^k$ is a filtered free module (direct sum of shifted $A$), and for $A(s)$, the map $\Psi$ corresponds to the identity on $(\mathrm{gr}_{\bullet}A)(-s)$. Thus $\Psi$ is an isomorphism (or see \cite[Lemma~D.2.3]{hotta2008d}).
        \item Therefore, $H^j(\mathrm{gr}_{\bullet}\check{F}^{\bullet}) \cong \mathrm{Ext}_{\mathrm{gr}_{\bullet}A}^j(\mathrm{gr}_{\bullet}M, \mathrm{gr}_{\bullet}A).$
    \end{enumerate}
    The result follows.
\end{proof}

\begin{remark}[{Induced filtration on $\Ext$ (cf. \cite[Appendix~D, in the proof of Lemma~D.2.4]{hotta2008d})}]\label{rmk:induced-filtration-on-Ext}
    \mbox{}

    To make sense of the spectral sequence convergence, we define the filtration on $\mathrm{Ext}$ group.

    Using the filtration on $\check{F}^\bullet$ defined above, the \textbf{induced filtration} on
    \[
        \Ext_A^i(M,N)\cong H^i(\check{F}^\bullet)
    \]
    is defined by
    \[
        F_p\Ext_A^i(M,N)
        := F_pH^i(\check{F}^\bullet)
        := \mathrm{Im}\!\left(H^i(F_p\check{F}^\bullet)\longrightarrow H^i(\check{F}^\bullet)\right),
    \]
    where $F_p\check{F}^\bullet$ denotes the subcomplex with $(F_p\check{F}^\bullet)^i=F_p\check{F}^i$.
    Note that this definition coincides with Definition \ref{def:induced_filtration_cohomology}.

    Indeed, the map $H^i(F_pK^\bullet) \to H^i(K^\bullet)$ sends the class $[z]_p$ of a cycle $z \in Z_p^i(\infty) = \ker(d) \cap F_p K^i$ to its class $[z]$ in $H^i(K^\bullet)$.
    Thus the image is
    \[
        \mathrm{Im}\!\left(H^i(F_pK^\bullet)\longrightarrow H^i(K^\bullet)\right) = \left\{ z + \mathrm{d}K^{i-1} \in H^i(K^\bullet) \mid z \in Z_p^i(\infty) \right\} = \frac{Z_p^i(\infty) + \mathrm{d}K^{i-1}}{\mathrm{d}K^{i-1}},
    \]
    which matches the definition above.
\end{remark}


\section{Duality for $\mathcal{D}_{X}$-mod}

Let $X = \mathrm{Spec}\, R$ a smooth affine variety with algebraic coordinates $(x_1, \ldots, x_n)$ and $M$ a nonzero f.g. $\mathcal{D}_X$-module.

\begin{theorem}\label{theorem 7.12}
    We have:
    \begin{enumerate}
        \item $\dim \operatorname{Ch}(\operatorname{Ext}_{\mathcal{D}_X}^i(M, \mathcal{D}_X)) \leq 2n - i$.
        \item $\operatorname{Ext}_{\mathcal{D}_X}^i(M, \mathcal{D}_X) = 0$ for $i < 2n - \dim(\operatorname{Ch}(M))$ or $i > n$.
        \item $j(M) + \dim \operatorname{Ch}(M) = 2n$.
    \end{enumerate}
\end{theorem}
\begin{proof}
    \textbf{(1) and (2)}
    By Corollary~\ref{corollary 7.11}, and the fact that the support of a subquotient is contained in the support of the original module, we have:
    $$ \operatorname{Ch}(\operatorname{Ext}_{\mathcal{D}_X}^i(M, \mathcal{D}_X)) = \operatorname{Supp}(\operatorname{gr}_{\bullet}(\operatorname{Ext}_{\mathcal{D}_X}^i(M, \mathcal{D}_X))) \subseteq \Ch(\mathrm{Ext}^{j}_{\mathrm{gr}_{\bullet}\mathcal{D}_{X}}(\mathrm{gr}_{\bullet}M,\mathrm{gr}_{\bullet}\mathcal{D}_{X})). $$
    Then apply Theorem \ref{theorem 7.2} to $\Ch(\mathrm{Ext}^{j}_{\mathrm{gr}_{\bullet}\mathcal{D}_{X}}(\mathrm{gr}_{\bullet}M,\mathrm{gr}_{\bullet}\mathcal{D}_{X}))$, since $\gr_{\bullet}\mathcal{D}_{X}$ is commutative.

    \medskip

    \textbf{(3)}  We know that $j(\mathrm{gr}_{\bullet}M)+\dim(\operatorname{Ch}(M)) = \dim R[\xi_{1},\dots,\xi_{n}] = 2n$ by Theorem \ref{theorem 7.2}(2). Thus, it suffices to show that $j(M)=j(\mathrm{gr}_{\bullet}M)$.

    Since by Theorem \ref{tehorem 7.9}, $\mathrm{gr}_{\bullet}\mathrm{Ext}_{\mathcal{D}_{X}}^{j}(M,\mathcal{D}_{X})$ is a subquotient of $\mathrm{Ext}^{j}_{\mathrm{gr_{\bullet}\mathcal{D}_{X}}}(\mathrm{gr}_{\bullet}M,\mathrm{gr}_{\bullet}\mathcal{D}_{X})$.
    Thus we obtain the inequality:
    \[
        j(M) \ge j(\mathrm{gr}_{\bullet}M) = 2n - \dim(\operatorname{Ch}(M)).
    \]

    Now, assume $j(\mathrm{gr}_{\bullet}M)=j_{0}$. This means $\mathrm{Ext}^{j_{0}}_{\mathrm{gr_{\bullet}\mathcal{D}_{X}}}(\mathrm{gr}_{\bullet}M,\mathrm{gr}_{\bullet}\mathcal{D}_{X})\neq 0$.
    Consider the spectral sequence (of $\check{F}^\bullet$) converging to $\mathrm{Ext}_{\mathcal{D}_{X}}^{\bullet}(M,\mathcal{D}_{X})$ (resulting from the $\omega$-regularity of $\check{F}^\bullet$,
    see Lemma \ref{lemma 7.8}, Remark~\ref{rmk:classical_convergence_theorem} and Lemma \ref{lem:checkF-regular}).
    We have $E^{j}_{\bullet}(1) \simeq \mathrm{Ext}^{j}_{\mathrm{gr}_{\bullet}\mathcal{D}_{X}}(\mathrm{gr}_{\bullet}M,\mathrm{gr}_{\bullet}\mathcal{D}_{X})$ (the filtration on $\mathrm{Ext}$ mentioned in Rmk.~\ref{rmk:induced-filtration-on-Ext} is the same as the one in Def.~\ref{def:induced_filtration_cohomology}).
    \begin{itemize}
        \item For $j < j_0$, $E^{j}_{\bullet}(1) = 0$.
        \item Thus $E^{j}_{\bullet}(k+1) \simeq \ker(E^{i}(k)\rightarrow E^{i+1}(k))$ for $k \ge 1$.
    \end{itemize}

    The order filtration on $\mathrm{gr}_{\bullet}\mathcal{D}_{X}$ is a $\mathbb{Z}_{\ge 0}$-filtration. Since $\Omega_{\bullet}M$ is a good filtration, it is bounded below (i.e., $\Omega_{p}M=0$ for $p \ll 0$). Consequently, the induced filtration on the resolution $\check{F}^{\bullet}$ is also bounded below.

    Let $p_{0}$ be the minimal integer such that $\check{F}^{j_{0}}_{p_{0}} \neq 0$.
    Then:
    \[
        E^{j_{0}}_{p_{0}}(0) = \check{F}^{j_{0}}_{p_{0}} \neq 0.
    \]
    For $k \ge 1$, the differential $\mathrm{d}_{k}: E^{j_{0}}_{p_{0}}(k) \rightarrow E^{j_{0}+1}_{p_{0}-k}(k)$ must be the zero map (since $p_0-k < p_0$, the target term is 0).
    Therefore,
    \[
        E^{j_{0}}_{p_{0}}(\infty) = E^{j_{0}}_{p_{0}}(0) \neq 0.
    \]
    This implies $\mathrm{gr}_{p_{0}}\mathrm{Ext}_{\mathcal{D}_{X}}^{j_{0}}(M,\mathcal{D}_{X}) \neq 0$, so $\mathrm{Ext}_{\mathcal{D}_{X}}^{j_{0}}(M,\mathcal{D}_{X}) \neq 0$.

    Hence $j(M) \le j_{0} = j(\mathrm{gr}_{\bullet}M)$. Combining with the reverse inequality, we get $j(M) = j(\mathrm{gr}_{\bullet}M)$.
\end{proof}

\begin{proposition}\label{prop:DX-has-finite-global-dimension}
    The global dimension of $\mathrm{d}_X$ is also finite. Furthermore, the global dimension of $\mathcal{D}_X$ $=n$.
\end{proposition}
\begin{proof}
    First, we show $\operatorname{gl.dim}\mathcal{D}_X \ge n$. Take $M=\mathcal{O}_{X}$. Since $M$ is holonomic, $\dim(\operatorname{Ch}(M))=n$. Using Theorem \ref{theorem 7.12}(3), we have $j(M)=2n-n=n$, which implies $\operatorname{Ext}^n(M, \mathcal{D}_X) \neq 0$. Thus $\operatorname{pd}(M) \ge n$, so $\operatorname{gl.dim}\mathcal{D}_X \ge n$.

    Now we show $\operatorname{gl.dim}\mathcal{D}_X \le n$. It suffices to show that for any non-zero finitely generated $\mathcal{D}_X$-module $M$, $\operatorname{Ext}^i(M, \mathcal{D}_X) = 0$ for all $i > n$.
    Suppose $\operatorname{Ext}^i(M, \mathcal{D}_X) \neq 0$ for some $i > n$. By Theorem \ref{theorem 7.12}(1), we have:
    $$\dim \operatorname{Ch}(\operatorname{Ext}^i(M, \mathcal{D}_X)) \le 2n - i < 2n - n = n.$$
    However, by Bernstein's Inequality (Theorem \ref{thm:Bernstein-inequality}, which implies $\dim \operatorname{Ch}(M) \ge n$), the dimension of the characteristic variety must be at least $n$. This is a contradiction.
    Therefore, $\operatorname{Ext}^i(M, \mathcal{D}_X) = 0$ for all $i > n$, which means $\operatorname{pd}(M) \le n$. Consequently, $\operatorname{gl.dim}\mathcal{D}_X \le n$.
\end{proof}

\begin{theorem}
    $\mathcal{D}_{X}$ is Auslander regular.
\end{theorem}
\begin{proof}
    By Proposition \ref{prop:DX-has-finite-global-dimension}, $\operatorname{gl.dim}\mathcal{D}_X = n$, a finite number.

    If $N \subset \operatorname{Ext}^{i}(M, \mathcal{D}_X)$ is a nonzero submodule, then $\operatorname{Ch}(N)\subset\operatorname{Ch}(\operatorname{Ext}^{i}(M, \mathcal{D}_X))$.

    So $\dim(\operatorname{Ch}(N)) \leq \mathrm{dim}(\operatorname{Ch}(\operatorname{Ext}^{i}(M, \mathcal{D}_X)))\leq 2n - i$ by Theorem \ref{theorem 7.12}(1).
    Then $j(N) \geq i$ by Theorem \ref{theorem 7.12}(3).
\end{proof}

\begin{corollary}
    \quad
    \begin{enumerate}
        \item $M$ is holonomic iff $\operatorname{Ext}^i(M, \mathcal{D}_X) = 0$, $i \neq n$.
        \item If $M$ is holonomic, then $\operatorname{Ext}^n(M, \mathcal{D}_X)$ is also holonomic.
    \end{enumerate}
\end{corollary}

\chapter{Algebraic normal deformation and V-filtrations}
\section{$V$-filtration and normal deformation}
Let $X = \text{Spec } R$ be a smooth affine variety over $\mathbb{C}$ with algebraic coordinates $(t, x_2, \ldots, x_n)$. 

Take $I^j =\begin{cases}\left< t \right>^{j},& j > 0\\R,&j\leq0\end{cases}$, we define a filtration of $\mathcal{D}_{X}$: $$V_{k} \mathcal{D}_X = \{P \in \mathcal{D}_X \mid P I^{j} \subseteq I^{j-k}\text{ for all }j\in \mathbb{Z}\}.$$Then we obtain that $V_0 \mathcal{D}_X = R \langle t\partial_t, \partial_{x_2}, \ldots, \partial_{x_n} \rangle$ and for all $i>0$: $$V_{-i} \mathcal{D}_{X} = t^{i} V_{0 }\mathcal{D}_{X }= V_{0} \mathcal{D}_{X} t^{i}\atop V_i \mathcal{D}_X = \partial_t V_{i-1} \mathcal{D}_X + V_{i-1} \mathcal{D}_X.$$Next, we will prove $V_{\bullet} \mathcal{D}_X$ is a noetherian $\mathbb{Z}$-filtration s.t. $\mathrm{gr}^V \mathcal{D}_X$ is Auslander regular. 

We obtain that $V_{j} \mathcal{O}_X = I^{-j} \in \mathcal{O}_X = R$ for $j \in \mathbb{Z}$ defines a $\mathbb{Z}$-filtration on $\mathcal{O}_X$.
We have 
$$\mathcal{R}_{\bullet}(\mathcal{O}_X) = \mathcal{R}^V_{\bullet}(\mathcal{O}_X) = \bigoplus_{j \in \mathbb{Z}} I^{-j} \otimes u^j \in R[u, u^{-1}]$$
and $\mathbb{C}[u] \hookrightarrow R_{\bullet}(\mathcal{O}_X)=R[u]\bigoplus\limits_{j\ge 1}I^{j}\otimes u^{-j}$ induces a family $\widetilde{X}= \operatorname{Spec} R_{\bullet}(\mathcal{O}_X)\xrightarrow{\varphi}\mathbb{A}^{1}_{\mathbb{C}}$ s.t. 
\begin{enumerate}
    \item $\varphi$ is smooth affine; 
    \item  $\varphi^{-1}(a) \simeq X, \quad \forall \text{ nonzero } a \in \mathbb{A}^1_{\mathbb{C}}$, since $\varphi ^{-1}(a)=X\times_{\mathbb{A}^{1}_{C}}\mathrm{Spec}(\mathbb{C}[u] /(u-a))=\mathrm{Spec}(R\otimes_{\mathbb{C}[u]}\mathbb{C}[u] /(u-a))=\mathcal{R}_{\bullet}(\mathcal{O}_{X}) /(u-a)\simeq R$. 
    \item $\varphi^{-1}(0) = \operatorname{Spec} \bigoplus\limits_{j \in \mathbb{Z}_{\geq 0}} I^j/I^{j+1} = T_Y X$ the normal bundle of $Y$ in $X$, where $Y = \left<  t = 0\right> \subseteq X$. 

\end{enumerate}

We say $\varphi$ is the $\textbf{normal deformation}$ along $Y$. 

Recall the following definition: 

\begin{definition}
    Let $A$, $B$ be two commutative rings with units with $\phi: A \to B$: a ring homomorphism. Then $\Omega_{\phi} = \Omega_{B/A}$ is the $B$-module generated by $df$, $f \in B$ s.t.
    \begin{enumerate}
        \item  $d(fg) = fdg + gdf$
        \item $d(af + bg) = adf + bdg$, $\quad \forall f,g \in B$, $a,b \in A$
    \end{enumerate}
 called the $\textbf{module of K?hler differentials}$ of $B$ over $A$. 

 Furthermore, we have the following definition: 
 \begin{enumerate}
    \item  $\phi$ is called $\textbf{smooth}$ if $\Omega_{\phi}$ is a free $B$-module of finite rank.
    \item If $\phi$ is smooth, we define $$\mathcal{T}_{\phi} := \mathrm{Hom}_B(\Omega_{\phi}, B),$$which is called the $\textbf{module of derivatives}$ of $B$ over $A$.
    \item $\mathcal{D}_{\phi} \subseteq \mathrm{End}_A(B)$ is the $A$-subalgebra generated by $\mathcal{T}_{\phi}$ and $B$, the $\textbf{ring of relative differential operators}$ of $B$ over $A$. 
 \end{enumerate}

\end{definition}

\begin{remark}
    if $A'$ is another $A$-algebra, then $\Omega_{B'/A'} \simeq \Omega_{B/A} \otimes_B B',$ where $B' = B \otimes_A A'$.
\end{remark}

\begin{lemma}
    
    Consider the ring homomorphism $\mathbb{C}[u]\hookrightarrow \mathcal{R}_\bullet(\mathcal{O}_X)$, we obtain that:

    \begin{enumerate}
        \item $\Omega_\varphi = \bigoplus_{j \in \mathbb{Z}} I^{-j} (dt \otimes u^{j-1}  \bigoplus\limits_{i=2}^{n} dx_{i} \otimes u^{j}) = \mathcal{R}_{\bullet}(\mathcal{O}_X) dt \otimes u^{-1}  \bigoplus\limits_{i=2}^{n} \mathcal{R}_{\bullet}(\mathcal{O}_X) dx_{i}.$
        \item $\mathcal{T}_\varphi = \bigoplus\limits_{j \in \mathbb{Z}} I^{-j} (\partial_t \otimes u^{j+1} \bigoplus_{i=2}^n \partial_i \cdot u^j)$
    \end{enumerate}
\end{lemma}
\begin{proof}
     Since $d(t^k \cdot u^{-k}) = t^{k-1} dt \cdot u^{-k}$, 1 follows. And taking graded dual, 2 follows.
\end{proof}

\begin{corollary}
    \quad
    \begin{enumerate}
        \item $\mathcal{R}_{\bullet}^V \mathcal{D}_X = \mathcal{D}_\varphi$
        \item  $\mathrm{gr}_{\bullet}^V \mathcal{D}_X \cong \mathcal{D}_{T_Y X}$
    \end{enumerate}
    Then $\mathcal{R}^V_{\bullet} \mathcal{D}_X$ is noetherian and  $\mathrm{gr}^{V}_{\bullet} \mathcal{D}_X$ is Auslander regular. Meanwhile, we obtain that $\bigcap\limits_{i \in \mathbb{Z}} V_i \mathcal{D}_X = 0$.  

\end{corollary}

\chapter{Bj\"ork's approach to existance of b-functions}
\section{Bj\"ork's approach to existence of b-functions}

\begin{theorem}[Bj\"ork]\label{Theorem 8.4}
    Let $M$ be a holonomic $\mathcal{D}_X$-module with a good filtration $\Omega_{\bullet} M$ over $V_{\bullet} \mathcal{D}_X$. There exists a non zero $b(s) \in \mathbb{C}[s]$ such that
    \[
        b(t\partial_t + k) \text{ kills } \frac{\Omega_k M}{\Omega_{k-1} M}, \quad \forall k \in \mathbb{Z}.
    \]
    The monic minimal degree $b(s)$ is called the $\textbf{$b$-function}$ of $M$ along $Y$.
\end{theorem}
\begin{remark}
    In the handwritten notes provided by the lecturer, says $b(t\partial_t \textcolor{red}{ - } k) \text{ kills } \frac{\Omega_k M}{\Omega_{k-1} M}$,
    but this is not compatible with the increasing filtration setting in Rmk.~\ref{rmk:increasing_V_filtration}.
\end{remark}


\begin{lemma}\label{lemma 8.5}
    Let $Y$ be a smooth variety over $\mathbb{C}$, $N$ a holonomic $\mathcal{D}_Y$-module, and $\theta\in \mathrm{End}_{\mathcal{D}_Y}(N)$. Then there exists $b(s)\in\mathbb{C}[s]$ such that $b(\theta)=0$.
\end{lemma}
\begin{proof}
    Since $N$ is holonomic, we obtain that
    \[
        \mathrm{Ch}(N) = \bigcup_{i} T^*_{Y_i} Y,
    \]
    with $Y_{i}$ locally closed.

    Take a $Y_i$ of maximal dimension and take an (affine) open subset $U \subset Y$
    such that $Y_{i}^{\circ} \hookrightarrow U$ is a closed embedding
    with $Y_{i}^{\circ}$ smooth, $\overline{Y_i^{\circ}} = Y_{i}$, $U \cap Y_i = Y_i^\circ$
    and $U \cap (\bigcup_{j\neq i}Y_j) = \varnothing$. Then
    \[
        \mathrm{Ch}(N|_U) = T^*_{Y_i^\circ} U.
    \]

    Since $Y^\circ$ is smooth, locally we can split algebraic coordinates to use Kashiwara's equivalence \ref{thm:kashiwara_equiv}, by shrinking $U$:
    \[
        N|_U = i_{+}M = M[\partial_1, \ldots, \partial_k].
    \]

    Consider $U^{\mathrm{an}} \to U$, the analytification functor, since $Y^\circ \hookrightarrow U$ is a closed embedding,
    by \cite[Proposition~4.7.2.(ii)]{hotta2008d} we have $i_{+}(M)^{\mathrm{an}} \to i^{\mathrm{an}}_{+}(M^{\mathrm{an}})$ is an isomorphism.

    We use Riemann-Hilbert correspondence on $(N|_U)^{\mathrm{an}}$:
    \begin{align*}
        (N|_U)^{\mathrm{an}} & \cong (i_{+}M)^{\mathrm{an}}                                                                                                       &  & \text{(By \ref{thm:kashiwara_equiv})}                                                                                   \\
                             & \cong i^{\mathrm{an}}_{+}(M^{\mathrm{an}})                                                                                                                                                                                                                      \\
                             & \cong M^{\mathrm{an}}[\partial_1, \ldots, \partial_k]                                                                              &  & \text{\begin{tabular}[t]{@{}l@{}}($i_+$ of a closed embedding behaves\\ the same on variety and manifold)\end{tabular}} \\
                             & \cong \mathcal{O}^{\mathrm{an}}_{Y_i^\circ} \otimes_{\mathbb{C}} L \otimes_{\mathbb{C}} \mathbb{C}[\partial_1, \ldots, \partial_k] &  & \text{(By \ref{cor:kashiwara_equiv_conormal}, $M$ is an integrable connection, use \ref{thm:basic-RH-corr})}
    \end{align*}
    Then we get
    \[
        (N|_U)^{\mathrm{an}} = \mathcal{O}^{\mathrm{an}}_{Y_i^\circ} \otimes L[\partial_1, \ldots, \partial_k]
    \]
    with $L$ a local system on $Y_i^\circ$.

    For the name \emph{Riemann-Hilbert correspondence}, one may think it is a \textbf{categorical equivalence}, that means there is not only a correspondence of objects but also a correspondence of \textbf{morphisms}.

    That is \cite[Theorem~4.2.4.]{hotta2008d}'s $H^{-\mathrm{d}_X}(\mathrm{DR}(-))$ ($X$ a complex manifold),
    which means for a morphism between two holomorphic $\mathcal{D}_X$-modules $f: M_1 \to M_2$ locally free over $\mathcal{O}_X$ (holomorphic),
    with respective flat sections
    \[
        L_1 := \Ker(\nabla_{M_1}) \quad \text{and}\quad L_2 := \Ker(\nabla_{M_2}) \quad \text{(see \ref{rmk:basic-RH-corr})},
    \]
    we have a morphism between their corresponding local systems
    \[
        H^{-\mathrm{d}_X}(\mathrm{DR}(f)): L_1 \to L_2
    \]
    which is just a restriction of $f$.

    Recall \ref{prop:classification-of-local-system} says if $X$ is connected, we have an another categorical equivalence
    to identify local system and finite diemsional representation of $\pi_1(X,x)$ (take an arbitrary base point $x$), morphisms of representations
    is called \textbf{intertwining operator} with respect to the action of $\pi_1(X,x)$.

    They compose a third categorical equivalence, since categorical equivalence is fully faithful, we have
    \begin{align*}
        \mathrm{End}_{\mathcal{D}_U}(N|_U) & = \mathrm{End}_{\mathcal{D}_{Y^\circ}}(M)                                     &  & \text{(Kashiwara equivalence is also a categorical equivalence)}          \\
                                           & \subseteq \mathrm{End}_{\mathcal{D}^{\mathrm{an}}_{Y^\circ}}(M^{\mathrm{an}})                                                                                \\
                                           & = \mathrm{End}_{\mathrm{LocSys}_{X}}(L)                                       &  & \text{(By \cite[Theorem~4.2.4.]{hotta2008d})}                             \\
                                           & = \mathrm{End}_{\pi_1(X,x)}(L_x)                                              &  & \text{(By \ref{prop:classification-of-local-system}, and $L_x$ is fiber)} \\
                                           & \subseteq \mathrm{End}_{\C}(L_x)                                              &  & \text{(Definition of intertwining operator)}
    \end{align*}
    Therefore, by Cayley-Hamilton theorem, there exists $b_i(s) \in \mathbb{C}[s]$ such that $b_i(\theta|_{U}) = 0$.

    Before we do induction, recall a fact that for a smooth variety $Y$, a irreducible closed subset $X$, we have
    \[
        T^*_X Y = \overline{T^*_{X\cap U} U} \subseteq T^*Y \quad \text{(for affine open $U \subset Y$ such that $X\cap U \subseteq X^{\mathrm{sm}}$ dense in $X$)}
    \]
    is also irreducible, by the fact that
    \[
        T^*_{X \cap U} U \subseteq T^*U.
    \]
    is closed irreducible in $T^*U$.

    Now, take  $N_1 = b_1(\theta)N$.
    Since $b_1(\theta|_U) = 0$, we have $N_1|_U = 0$, which implies $\mathrm{Ch}(N_1) \cap T^*U = \varnothing$.

    By Bernstein's inequality \ref{thm:Bernstein-inequality},
    $N_1$, as a submodule of a holomonic $\D_Y$-module is also a holomonic $\D_Y$-module.
    So irreducible component of $\mathrm{Ch}(N_1)$ has dimension $n = \dim Y$,
    and is contained in some irreducible component of $\mathrm{Ch}(N) = \bigcup_k T^*_{Y_k} Y$.

    Since all irreducible component of a holomonic $\D_Y$-module has dimension $n = \dim Y$ (see Cor.~\ref{cor:holonomic-char-cycle}),
    and we have a fact that a proper closed subset of an irreducible variety has strictly smaller dimension,
    it follows that each irreducible component of $\mathrm{Ch}(N_1)$ must be equal to some $T^*_{Y_k} Y$.

    But it can't be $T^*_{Y_i} Y$, because $\mathrm{Ch}(N_1) \cap T^*U = \varnothing$.
    So we strictly reduce the number of components in the characteristic variety,
    so we are done by induction.


\end{proof}




\begin{theorem}\label{Theorem 8.6}
    Let $M$ be a finitely generated $\mathcal{D}_{X}$-module, and $\Omega_{\bullet}M$ a good filtration over $V_{\bullet}\mathcal{D}_{X}$. Then $j(\mathrm{gr}_{\bullet}^\Omega M)\ge j(M)$.
\end{theorem}
\begin{remark}
    In Thm.~\ref{theorem 7.12} we also proved $j(\gr_\bullet M)\ge j(M)$, but this time we have a different filtration, the $V$-filtration.

    We need the fact that $\gr_\bullet^{V}\D_X$ is Auslander regular, see Cor.~\ref{cor:Rees-ring-of-V-filtration-D-module}. So it can't replace the first proof in Thm.~\ref{theorem 7.12}.
\end{remark}
\begin{proof}
    If $\mathrm{gr}^\Omega_{\bullet}M=0$, there is nothing to prove, so we assume $\mathrm{gr}^\Omega_{\bullet}M\neq 0$ and $\ell=j(\mathrm{gr}_{\bullet}^{\Omega}M)<j(M)$.

    Then the $E_1$-page of the spectral sequence of $\check{F}^\bullet$
    introduced in Section~\ref{sect:more_on_filtered_rings},
    that converges to $\mathrm{Ext}_{\mathcal{D}_{X}}^{j}(M,\mathcal{D}_{X})$,
    satisfies:
    \[
        0 \rightarrow \operatorname{Ext}^{\ell}(\operatorname{gr}_{\bullet}^{\Omega} M, \operatorname{gr}^{V}_{\bullet}\mathcal{D}_X)
        \xrightarrow{\mathrm{d}_1}
        \operatorname{Ext}^{\ell+1}(\operatorname{gr}_{\bullet}^{\Omega} M, \operatorname{gr}^{V}_{\bullet}\mathcal{D}_X).
    \]
    Where we identify $\gr_\bullet \Hom_{\D_X}(F^\bullet, \D_X)$ with $\Hom_{\gr_\bullet^{V} \D_X}(\gr_\bullet F^\bullet, \gr_\bullet^{V} \D_X)$
    as in the proof of Cor.~\ref{corollary 7.11}.

    Since $\ell = j(\gr_\bullet^\Omega M)$, so we obtain that $E^{\ell}_{\bullet}(2)=\ker \mathrm{d}_1\subseteq\operatorname{Ext}^{\ell}(\operatorname{gr}_{\bullet}^{\Omega} M, \operatorname{gr}^{V}_{\bullet}\mathcal{D}_X)$ since no images map to $E^{\ell}_{\bullet}(1)$.

    Then we have the following exact sequence with the quotient $N=E^{\ell}_{\bullet}(1) /\ker \mathrm{d}_{1}\simeq \mathrm{d}_{1}(E_{\bullet}^{\ell}(1))\subseteq E^{\ell+1}_{\bullet}(1)$:
    \[
        0 \rightarrow E_{\bullet}^{\ell}(2) \rightarrow \operatorname{Ext}^{\ell}(\operatorname{gr}_{\bullet}^{\Omega} M, \operatorname{gr}^{V}_{\bullet}\mathcal{D}_X)
        \rightarrow N \rightarrow 0.
    \]
    In general, by induction, $\forall k\ge 1$, we also have such exact sequence:
    \[
        0 \rightarrow E_{\bullet}^{\ell}(k+1) \rightarrow E_{\bullet}^{\ell}(k) \rightarrow N_k \rightarrow 0.
    \]

    Since $j(M)>\ell$, and $\check{F}^\bullet$ is convergent (see Definition~\ref{def:convergent_spectral_sequence}),
    we have $E_{\bullet}^{\ell}(r) = \operatorname{gr}_{\bullet} \operatorname{Ext}^{\ell}(M, \mathcal{D}_X) = 0$ for some $r\gg 0$,
    and $E^{\ell}_{\bullet}(r)=\Ker \mathrm{d}_{r}$.

    By triangle inequality of grade number (one can get this inequality by a long exact sequence corresponds to the short exact sequence), we have
    \[
        j(E_{\bullet}^{\ell}(k)) \ge \min(j(E_{\bullet}^{\ell}(k+1)), j(N_k)).
    \]
    Note that $N_k$ is a subquotient of $E_{\bullet}^{\ell+1}(1) = \operatorname{Ext}^{\ell+1}(\operatorname{gr}_{\bullet}^{\Omega} M, \operatorname{gr}^{V}_{\bullet}\mathcal{D}_X)$.
    Since $\gr_\bullet^{V}\D_X$ is Auslander regular, we have $j(N_k) \ge \ell+1$.

    Since $E_{\bullet}^{\ell}(r) = 0$ for $r \gg 0$, we have $j(E_{\bullet}^{\ell}(r)) = \infty$.
    By induction (downwards on $k$), we obtain
    \[
        j(E_{\bullet}^{\ell}(1)) \ge \ell+1.
    \]
    However, $E_{\bullet}^{\ell}(1) = \operatorname{Ext}^{\ell}(\operatorname{gr}_{\bullet}^{\Omega} M, \operatorname{gr}^{V}_{\bullet}\mathcal{D}_X)$.
    Since $\ell = j(\operatorname{gr}_{\bullet}^{\Omega} M)$, we have $j(E_{\bullet}^{\ell}(1)) = \ell$ by Theorem~\ref{theorem 7.2}'s (2) and (4).
    This contradicts $j(E_{\bullet}^{\ell}(1)) \ge \ell+1$.
\end{proof}

\begin{proof}[proof of (\ref{Theorem 8.4})]
    \quad

    By Theorem~\ref{Theorem 8.6},
    $j(\mathrm{gr}_{\bullet}^\Omega M)\ge j(M)=n$.
    If $\mathrm{gr}_{\bullet}^\Omega M=0$, take $b(s)=1$.

    Since by Cor.~\ref{cor:Rees-ring-of-V-filtration-D-module} and filtration is good, we view $\mathrm{gr}_{\bullet}^\Omega M$ as a finitely generated$\gr_\bullet^{V}\D_X = \D_{T_Y X}$-module.
    And in the proof of Cor.~\ref{cor:Rees-ring-of-V-filtration-D-module}, we verified that $T_{Y}X$ is a smooth variety of dimension $n$ over $\mathbb{C}$.
    So if $\gr_\bullet^{\Omega}M \neq 0$, then $\mathrm{dim}(\mathrm{Ch}(\mathrm{gr}^\Omega_{\bullet}M))\ge n$ by Thm.~\ref{thm:Bernstein-inequality}.

    Since $\mathrm{gr}_{\bullet}^V\mathcal{D}_{X}$ is a Auslander regular ring,
    by Thm.~\ref{theorem 7.12},
    we obtain that $j(\mathrm{gr}_{\bullet}^\Omega M)+\mathrm{dim}(\mathrm{Ch}(\mathrm{gr}_{\bullet}^\Omega M))=2n=2\,\mathrm{dim}T_{Y}X$, then $\mathrm{gr}^\Omega_{\bullet} M$ is holonomic over $T_Y X$.

    Now, take
    \[\theta=\bigoplus \limits_{i}(t\partial_{t}+i) : \gr_\bullet^\Omega M \to \gr_\bullet^\Omega M\quad \text{ for } t\partial_t + i: \gr_i^\Omega M \to \gr_i^\Omega M\]
    since one can verify that for $P\in\D_X$
    \[P \in V_{i}\D_X \Leftrightarrow P = p(x,\partial_x) t^{\alpha}\partial_t^\beta\quad \text{ for } p(-,-) \text{ a polynomial, and }\beta - \alpha \le i\]
    which implies
    \[[t\partial_t, P] = -i P \text{ for } P\in\gr_i^{V}\mathcal{D}_X,\]
    so we have $\theta\in\mathrm{End}_{\mathcal{D}_{T_{Y}X}}(\mathrm{gr}_{\bullet}^\Omega M)$.
    Indeed, for $u \in \gr_j^\Omega M$ and $P \in \gr_k^{V} \D_X$, we calculate
    \[
        \theta(Pu) = (t\partial_t + j + k) Pu = (P t\partial_t - k P + j + k) u = P(t\partial_t + j) u = P \theta(u).
    \]

    By Lemma \ref{lemma 8.5}, there exists $b(s)\in\mathbb{C}[s]$ s.t. $b(\theta)=0$,
    In particular, $b(t\partial_{t}+k)$ kills $\frac{\Omega_{k}M}{\Omega_{k-1}M}$ for all $k$.
\end{proof}

\begin{remark}\label{rmk:b_function_of_C[x,1/x]e^1/x}
    It is possible that $\mathrm{gr}_{\bullet}^\Omega M=0$. In fact, take
    $M := \mathcal{D}_{\Spec{\C[t]}} \big/ \mathcal{D}_{\Spec{\C[t]}}\langle t^2\partial_t + 1 \rangle = \mathcal{D}_{X}\cdot e^{\frac{1}{t}}$,
    for $e^{\frac{1}{t}}$ represents the canonical section $[1]\in\mathcal{D}_{\Spec{\C[t]}} \big/ \langle t^2\partial_t + 1 \rangle$.

    We note that $t$ also acts bijectively on $M$ (just as Example~\ref{ex:first-examples-of-d-modules}),
    so $M=M[t^{-1}]  = \mathbb{C}[t, \frac{1}{t}]e^{\frac{1}{t}}$,
    take $\Omega_{k}M=V_{k}\mathcal{D}_{X}\cdot e^{\frac{1}{t}}$,
    which is a good filtration.
    Since $t\partial_{t}e^{\frac{1}{t}}= -\frac{1}{t} e^{\frac{1}{t}}$, we have $\Omega_{k}M=M$ and $b(s)=1$.
\end{remark}

\begin{theorem}[Kashiwara-Malgrange]\label{Theorem 8.7}
    Let $M$ be a holonomic $\mathcal{D}_{X}$-module, there exists unique good filtration $V_{\bullet}M$ and $b(s)\in\mathbb{C}[s]$ s.t.
    \begin{enumerate}
        \item $\mathrm{Re}(\text{roots of }b(s))\subseteq[0,1)$;
        \item $b(t\partial_{t}+k)$ kills $\frac{V_{k}M}{V_{k-1}M}$.
    \end{enumerate}
    And the filtration is called $V$-filtration.
\end{theorem}
\begin{proof}
    This is Kashiwara's trick:
    First, set $\Gamma_{n}=\Omega_{n+k}$. Note that $b(t\partial_{t}+n+k)$ kills $\Gamma_{n} /\Gamma_{n-1}$ for all $n$,
    which means the roots of the $b$-function are shifted by $k$.
    Thus by taking $k\gg 0$, after replacing $b(s)$ by $s\mapsto b(s+k)$, we can assume $\mathrm{Re}(\lambda)\le 0$ for all roots $\lambda$ of $b(s)$.

    Next, suppose that $b(s)$ has a root $\alpha$ whose real part is $< 0$.
    Set $\Gamma_{n}'=\Gamma_{n-1}+(t\partial_{t}+n-\alpha)\Gamma_{n}$.
    If $\beta(s) = b(s)/(s-\alpha)$, then we verify that $\beta(s+n)(s+n-\alpha-1)$ annihilates $\Gamma'_{n} / \Gamma'_{n-1}$ for all $n$.

    Denote $s = t\partial_t$. Let $u' \in \Gamma'_{n}$. We can write $u' = v + (s+n-\alpha)w$ with $v \in \Gamma_{n-1}$ and $w \in \Gamma_{n}$.
    We calculate the action of $Q(s) = \beta(s+n)(s+n-\alpha-1)$:
    \[
        Q(s)u' = \beta(s+n)(s+n-\alpha-1)v + \beta(s+n)(s+n-\alpha-1)(s+n-\alpha)w.
    \]
    For the first term involving $v \in \Gamma_{n-1}$: note that $(s+n-\alpha-1)v \in (s+(n-1)-\alpha)\Gamma_{n-1} \subseteq \Gamma'_{n-1}$.
    For the second term involving $w \in \Gamma_{n}$: we use the fact that polynomials in $s$ commute, so
    \[
        \beta(s+n)(s+n-\alpha-1)(s+n-\alpha)w = (s+n-\alpha-1) \beta(s+n)(s+n-\alpha) w = (s+n-\alpha-1) b(s+n) w.
    \]
    By the hypothesis on $\Gamma_\bullet$, we have $b(s+n)w \in \Gamma_{n-1}$. Let $z = b(s+n)w$.
    Then $(s+n-\alpha-1)z \in (s+(n-1)-\alpha)\Gamma_{n-1} \subseteq \Gamma'_{n-1}$.
    Thus, $Q(s)u' \in \Gamma'_{n-1}$.

    Continue doing this if $\alpha$ has multiplicity, we will finally move a root $\alpha$ of $b(s)$ to $\alpha+1$.
\end{proof}

\begin{definition}
    $M$ is a holomonic module, then:

    \begin{enumerate}
        \item  $M$ is called $Q\textbf{-specializable}$ along $Y$ if the $b$-function has all its roots in $Q$, where $Q \subseteq \mathbb{C}$ is a subfield.

              Since any subfield of $\C$ must contain its \textbf{prime subfield} $\mathbb{Q}$, and $b(t\partial_t + k)$ kills $V_kM/V_{k-1}M$,
              the eigenvalues $\lambda$ of $t\partial_t$ on $V_kM/V_{k-1}M$ kills its eigenvectors by multiplying $b(\lambda + k)$.

              So $\lambda$ differs from the roots of $b(s)$ by integers.
              Thus, if all roots of $b(s)$ are in $Q$, then all eigenvalues of $t\partial_t$ are in $Q$ (as $\mathbb{Z} \subseteq \mathbb{Q} \subseteq Q$).

        \item If $M$ is $\mathbb{Q}$-specializable,
              then for all $\beta \in \mathbb{Q}$,
              take $k = \lceil \beta \rceil$,
              then define
              $$V_{\beta} M = \pi^{-1} \left( \bigoplus_{\alpha \leq \beta} \left( \frac{V_k M}{V_{k-1} M} \right)^{\alpha} \right),$$
              and
              $$V_{<\beta} M = \pi^{-1} \left( \bigoplus_{\alpha < \beta} \left( \frac{V_k M}{V_{k-1} M} \right)^{\alpha} \right),$$
              where $\left( \frac{V_k M}{V_{k-1} M} \right)^{\alpha}$ denotes the $\alpha$-generalized eigenspace of $\frac{V_k M}{V_{k-1} M}$
              with respect to the $t\partial_t$-action, and $\pi: V_k M \rightarrow \frac{V_k M}{V_{k-1} M}$;
              Then we get a $\mathbb{Q}$-indexed $V$-filtration s.t.
              $$\mathrm{gr}_{\beta}^{V} M = \frac{V^{\beta} M}{V^{<\beta} M} = \left( \frac{V_k M}{V_{k-1} M} \right)^{\beta}$$
              with $t\partial_t - \beta$ acts on $\mathrm{gr}_{\beta}^{V} M$ nilpotently.
    \end{enumerate}
\end{definition}
\begin{example}
    Let $X=\mathbb{A}^{1}_{\mathbb{C}}$ and $M = \mathbb{C}\left[ t, \frac{1}{t} \right]e^{\frac{1}{t}}$ an irregular holonomic $\mathcal{D}$-mod.

    Consider $V$-filtration (of $\D_X$) along $t = 0$. The $V$-filtration on $M$ exists by Theorem \ref{Theorem 8.7}.

    Recall Rmk.~\ref{rmk:b_function_of_C[x,1/x]e^1/x}. The filtration $\Omega_\bullet M$ is defined by $\Omega_i M := V_i \mathcal{D}_{X} \cdot e^{\frac{1}{t}}$ for $i \in \mathbb{Z}$.
    We have $\Omega_i M = M$ for all $i\in \mathbb{Z}$.

    This filtration gives a $b$-function $b(s) = 1$. Which has all its roots in $[0,1)$, so $\Omega_\bullet M$ is exactly the $V$-filtration on $M$.
\end{example}

\begin{theorem}
    If M is a regular holonomic $\mathcal{D}_X$-mod, then its V-filtration is not trivial. In particular, its b-function is not $1$.
\end{theorem}

\begin{proof}[Sketch of proof]
    First, each $V_i M$ is a $V_0 \mathcal{D}_X$-module and $V_0 \mathcal{D}_X = \mathcal{D}_X^{\log}$ w.r.t. $(t = 0)$. Since $V_i M$ has no $t$-torsion for all $i < 0$, $\mathrm{Ch}(V_i M)$ has no component over $(t = 0)$, for all $i < 0$.

    Further, we obtain that $\mathrm{Ch}^{\log}\left(\frac{V_i M}{V_{i-1} M}\right) = \mathrm{Ch}^{\log}\left(\frac{V_i M}{t V_i M}\right)= \mathrm{Ch}^{\log}(V_i M)\big|_{(t=0)}$, thus $\frac{V_i M}{V_{i-1} M} \neq 0.$

\end{proof}

\begin{definition}[$\D$-module theoretic direct image of a graph embedding]
    Firstly we introduce the definition of graph embedding: 1

\end{definition}

\begin{example}\label{Example 8.11}

    Let $X = \mathbb{A}_{\mathbb{C}}^n$ and $Y = \mathbb{A}_{\mathbb{C}}^n \times \mathbb{A}_{\mathbb{C}}^1$, here $\mathbb{A}^{1}_{\mathbb{C}}$ is with coordinate $t$.
    For $f \in \mathbb{C}[x_1, \cdots, x_n]$, Consider $M = \mathbb{C}[x_1, \cdots, x_n]_f[s] f^s$.

    Define $t M = f M$ and then $t^{-1} M = f^{-1} M$.  Take $s = -t\partial_t$, then $\partial_t = -t^{-1} s$, the $\partial_{t}$-action is defined, and further, $M$ is a $\mathcal{D}_Y = \mathcal{D}_X \langle t, \partial_t \rangle$-module.

    Then $N := \mathbb{C}[x_1, \cdots, x_n][\partial_t] f^s = \mathbb{C}[x_1, \cdots, x_n][\partial_t] f^s= \mathbb{C}[x_1, \cdots, x_n][t^{-1} s] f^s$, and $\partial_{x_i} f^s = \partial_{x_i}(f) \cdot \frac{1}{f} s f^s= \partial_{x_i}(f) \cdot (-\partial_t) f^s$, thus $N$ is a $\mathcal{D}_Y$-submodule of $M$. Furthermore, we can check that $\mathrm{Ch}(N)$ = the conormal bundle of the graph of $X$, which indicates that $N$ is a holonomic $\mathcal{D}_{Y}$-module.

    We consider $V_{i}\mathcal{D}_{Y}\cdot f^{s}$ forming a good filtration of $N$: $$V_0 \mathcal{D}_Y \cdot f^s = \mathcal{D}_X \langle t, s \rangle f^s = \mathcal{D}_X[s] f^s$$Similarly, $V_{-1} \mathcal{D}_Y \cdot f^s = \mathcal{D}_X[s] f^{s+1}$.

    By Theorem \ref{Theorem 8.4}, there exists $b(s) \in \mathbb{C}[s]$ s.t.$$
        b(s) \cdot \frac{V_0 \mathcal{D}_Y \cdot f^s}{V_{-1} \mathcal{D}_Y \cdot f^s} = b(s) \cdot \frac{\mathcal{D}_X[s] f^s}{\mathcal{D}_X[s] f^{s+1}} = 0.$$which is to say there exists $P(s)\in \mathcal{D}_{X}[s]$ such that $b(s) f^s = P(s) \cdot f^{s+1}$

    In fact, $N = \iota_{f,+} (\mathbb{C}[x_1, \ldots, x_n])$, the $\mathcal{D}$-module pushforward, where $$\iota_f: X \to Y = X \times \mathbb{A}^{1}_{\mathbb{C}}\atop x \mapsto (x, f(x))$$the graph embedding. So $N$ is regular holonomic, and $b(s) \neq 1$.
\end{example}

\begin{remark}
    One can replace $\mathbb{C}[x_1, \ldots, x_n]$ by any regular holonomic $\mathcal{D}_X$-module, to get similar results in general.
\end{remark}

\chapter{Application: Nearby cycles and comparison}
\section{Two definitions of nearby cycles}

On $\mathbb{A}_{\mathbb{C}}^1$ with coordinate $t$, given any $\alpha \in \mathbb{C}$, we have a left indecomposable $\D_{{A}^1_\C}$ module
\[
    K^\alpha_m :=  \frac{\D_{\mathbb{A}^1_\C}}{\D_{\mathbb{A}^1_\C}\cdot(t\partial_t - \alpha)^m}
\]
First note that $t\in K^\alpha_m$ has \textbf{both} a left and a right inverse, i.e., action by multiplying $t$ induce a bijection.
So $K^\alpha_m$ not only a (left, see 2. in Lem.~\ref{lem:quasi-coherent-sheaf-and-localization-of-D-modules}) $\C[t]\subseteq \D_{\mathbb{A}^1_\C}$ module, but also a (left) $\C[t,t^{-1}]$-module that extending the $\C[t]$-module structure.

We take sections $e^\alpha_i := \overline{(t\partial_t - \alpha)^{m-1-i}}$ (equivalence class) then we have a presentation:
\[
    K^\alpha_m = \bigoplus_{i=0}^{m-1} \mathbb{C}[t, t^{-1}] e^\alpha_i
\]
with $\partial_t e^\alpha_i = \frac{1}{t}(\alpha e^\alpha_i + e^\alpha_{i-1})$ for all $i$ (note we have claimed that $t$ is invertible), where $e^\alpha_{-1} = 0$,
and the direct sum notation should be treated as \textbf{internal} direct sum.

Then when restricted to $\mathbb{A}_{\mathbb{C}}^1 \setminus \{0\}$, $K^\alpha_m\mid_{\mathbb{A}_{\mathbb{C}}^1 \setminus \{0\}}$
becomes a \textbf{finitely generated}, (locally) free $\mathcal{O}_{\mathbb{A}_{\mathbb{C}}^1 \setminus \{0\}}$-module, i.e., an \textbf{integrable connection} (see \hyperref[def:integrable-connection]{definition}).

Then the matrix (note $K^\alpha_m$ is not a vector space since $\C[t,t^{-1}]$ is not a field) of $t\partial_t$-action on $K_m^{\alpha}$ is $J_{\alpha,m}$, the $m\times m$ Jordan block with eigenvalue $\alpha$:
\[
    J_{\alpha,m} = \begin{bmatrix}
        \alpha & 1      &        &        \\
               & \alpha & \ddots &        \\
               &        & \ddots & 1      \\
               &        &        & \alpha
    \end{bmatrix}_{m \times m}
\]

\begin{remark}
    Here $e_\ell^{\alpha}$ can be understood as the formal symbol
    of the multivalued function
    \[
        t^{\alpha} \cdot \frac{(\log t)^\ell}{\ell!}.
    \]
    And $K^\alpha_{m}$ indicates about the section not flat,
    a \textbf{flat section} means the section is in the kernel of connection
    \[
        \nabla:=t\partial_t \in\End_\C(K^\alpha_m).
    \]

    Indeed in \cite[Ch.~5]{hotta2008d},
    they develop arguments or definitions such as ``moderate growth'' by multivalued functions,
    so it is not only as a heuristic symbol to introduce multivalued function.
\end{remark}

Then locally around a simply connected neighborhood (Euclidean topology) of $(\mathbb{A}_{\mathbb{C}}^1 \setminus \{0\})^{\mathrm{an}}$,
with simply connectedness, we obtain that $\mathrm{\log}\, t$ is a well-defined analytic function (so we should consider the analytification of the module $K^\alpha_m$).
Consider
$$
    v_{\ell} := e^{-J_{\alpha,m} \log t} \cdot e_{\ell}^{\alpha} \in (K^\alpha_m)^{\mathrm{an}}
    = \frac{\D^{\mathrm{an}}_{(\mathbb{A}^1_\C)^\mathrm{an}}}{\D^{\mathrm{an}}_{(\mathbb{A}^1_\C)^\mathrm{an}}\cdot(t\partial_t - \alpha)^m}
    = \bigoplus_{i=0}^{m-1} \mathcal{O}^{\mathrm{an}}_{(\mathbb{A}^1_\C)^\mathrm{an}}\left[\frac{1}{t}\right] e^\alpha_i,
$$
for $\ell = 0, \cdots, m-1$,
we then have, by the Leibniz rule,
\begin{align*}
    t\partial_t \circ \left(e^{-J_{\alpha,m} \log t} \cdot e_{\ell}^{\alpha}\right)
     & := t\partial_t \cdot e^{-J_{\alpha,m} \log t} \cdot e_{\ell}^{\alpha}                                                                                                                                                                                             \\
     & = \left[t\partial_t, e^{-J_{\alpha,m} \log t}\right] \cdot e_{\ell}^{\alpha} + e^{-J_{\alpha,m} \log t} \cdot t\partial_t \cdot e_{\ell}^{\alpha}                                                                                                                 \\
     & = t\partial_t\left(e^{-J_{\alpha,m} \log t}\right) \cdot e_{\ell}^{\alpha} + e^{-J_{\alpha,m} \log t} \cdot t\partial_t \cdot e_{\ell}^{\alpha}                              &  & \text{(see \cite[Prop~2.25(iii)]{mustata2023dmodules})}                         \\
     & = t \left( -J_{\alpha,m} \frac{1}{t} e^{-J_{\alpha,m} \log t} \right) \cdot e_{\ell}^{\alpha} + e^{-J_{\alpha,m} \log t} \cdot \left( J_{\alpha,m} e_{\ell}^{\alpha} \right)                                                                                      \\
     & = -J_{\alpha,m} \left( e^{-J_{\alpha,m} \log t} \cdot e_{\ell}^{\alpha} \right) + J_{\alpha,m} \left( e^{-J_{\alpha,m} \log t} \cdot e_{\ell}^{\alpha} \right)               &  & \text{($\C$ is in the \textbf{center} of $\D^{\an}_{(\mathbb{A}^1_\C)^{\an}}$)} \\
     & = 0
\end{align*}
we use $\cdot$ to denote the multiplication in $\D^{\an}_{(\mathbb{A}^1_\C)^{\an}}$, and we get $v_{\ell}$ are $t\partial_t$-flat.

\bigskip

We have a local system, recall the equivalence given in Prop.~\ref{prop:classification-of-local-system}:
$$
    L_m^{\lambda}=\left(\bigoplus \limits_{i=0}^{m-1}\mathbb{C}\cdot v_{i}, e^{-2\pi iJ_{\alpha,m}}\right)
$$
There are several things to be point out:
\begin{enumerate}
    \item Taking fiber (not stalk) at any given point $p\in(\mathbb{A}_{\mathbb{C}}^1 \setminus \{0\})^{\mathrm{an}}$, $\Ker(\nabla)$ becomes a finite ($m$) dimensional $\C$ vector space.
    \item On a simply connected Euclidean open set $p \in U \subset (\mathbb{A}_{\mathbb{C}}^1 \setminus \{0\})^{\mathrm{an}}$, we can give $\log\, t(p)$ a concreate $\C$-value.
    \item View $e^\alpha_i(p)$ as the canonical basis vector (the $i+1$-th entry is $1$, others are $0$)
          \[
              \begin{pmatrix}
                  0 \\ \vdots \\ 1 \\ \vdots \\ 0
              \end{pmatrix}\in\C^m,
          \]
          where the $1$ is at the $(i+1)$-th row. Then we identify
          \[
              v^\alpha_i(p) = e^{-J_{\alpha,m} \log t}(p) \cdot e^\alpha_i(p)
          \]
          with
          \[
              e^{-J_{\alpha,m} \log t}(p) \cdot \begin{pmatrix}
                  0 \\ \vdots \\ 1 \\ \vdots \\ 0
              \end{pmatrix} \in \C^m.
          \]
    \item So the monodromy action (of a generator of $\pi_1((\mathbb{A}_{\mathbb{C}}^1 \setminus \{0\})^{\mathrm{an}}, p)$) acts on $\C^m$ by left multiplication with $e^{-2\pi i J_{\alpha,m}}$.
\end{enumerate}
We denote $\lambda = e^{-2\pi i \alpha}\in \C$.

\bigskip

Since $K^\alpha_m\mid_{\mathbb{A}_{\mathbb{C}}^1 \setminus \{0\}}$ is an integrable connection on $\mathbb{A}_{\mathbb{C}}^1 \setminus \{0\}$ (note that $\C[t, t^{-1}]$ is also smooth finite type over $\C$), recall by Riemann-Hilbert correspondence Thm.~\ref{thm:basic-RH-corr}, we have \cite[Thm.~7.2.1]{hotta2008d} that in $\D^b(\underline{\C}_{X^\mathrm{an}})$,
$$
    \text{DR}\left((K_m^{\alpha}\mid_{\mathbb{A}_{\mathbb{C}}^1 \setminus \{0\}})^{\mathrm{an}}\right) \xrightarrow{\substack{\text{quasi-isomorphic}\\ \sim}}  L_m^{\lambda}.
$$
Since the direct system (in category of coherent (finitely generated) $\D_{\mathbb{A}_{\mathbb{C}}^1}$-modules)
\begin{equation*}
    \begin{array}{ccccccc}
        K_{m-1}^{\alpha} & \lhook\joinrel\xrightarrow{\quad} & K_m^{\alpha}                              & \lhook\joinrel\xrightarrow{\quad} & K_{m+1}^{\alpha}                            & \lhook\joinrel\xrightarrow{\quad} & \cdots \\[1.5ex]
        \overline P      & \xmapsto{\quad}                   & \overline{P \cdot (t\partial_t - \alpha)} & \xmapsto{\quad}                   & \overline{P \cdot (t\partial_t - \alpha)^2} & \xmapsto{\quad}                   & \cdots
    \end{array}
\end{equation*}
induces a direct system of local systems
(categorical equivalence is given by functors, so it also transports morphisms),
arrows should be viewed as $\pi_1((\mathbb{A}_{\mathbb{C}}^1 \setminus \{0\})^\mathrm{an}, p)$-\textbf{equivariant} $\C$-linear maps by Prop.~\ref{prop:classification-of-local-system}:
\begin{equation*}
    \begin{array}{ccccccc}
        L_{m-1}^{\lambda}=\bigoplus \limits_{i=0}^{m-2}\C\cdot v_i & \lhook\joinrel\xrightarrow{\quad} & L_m^{\lambda}=\bigoplus \limits_{i=0}^{m-1}\C\cdot v_i & \lhook\joinrel\xrightarrow{\quad} & L_{m+1}^{\lambda}=\bigoplus \limits_{i=0}^{m}\C\cdot v_i & \lhook\joinrel\xrightarrow{\quad} & \cdots \\[1.5ex]
        v_i                                                        & \xmapsto{\quad}                   & v_i                                                    & \xmapsto{\quad}                   & v_i                                                      & \xmapsto{\quad}                   & \cdots
    \end{array}
\end{equation*}
To see why the morphism are $\pi_1((\mathbb{A}_{\mathbb{C}}^1 \setminus \{0\})^\mathrm{an}, p)$-equivariant,
since $J_{\alpha, m}$ is an upper-triangular matrix, then $e^{-2\pi iJ_{\alpha,m}}$ is also upper-triangular, so each $L^\lambda_{i}\subseteq L^\lambda_{i+1}$ is a sub $\pi_1((\mathbb{A}_{\mathbb{C}}^1 \setminus \{0\})^\mathrm{an}, p)$-representation.

\bigskip

Let $X = \text{Spec } R$ be a smooth affine algebraic variety over $\mathbb{C}$, and $f \in R$,
a regular function, we have the following commutative diagram:
\begin{center}
    \begin{tikzcd}
        {(f=0)\atop \text{divisor}} & X && {\mathbb{A}^1_{\mathbb{C}}} \\
        \\
        & {U=X\backslash(f=0)} && {\mathbb{A}^1_{\mathbb{C}}\backslash \{0\}}
        \arrow["i", hook, from=1-1, to=1-2]
        \arrow["f", from=1-2, to=1-4]
        \arrow["j", hook', from=3-2, to=1-2]
        \arrow["f", from=3-2, to=3-4]
        \arrow[hook', from=3-4, to=1-4]
    \end{tikzcd}
\end{center}
where the morphism $f$ is given by the ring map:
\begin{align*}
    \C[t] & \xrightarrow{\quad} R                       \\
    \C    & \xrightarrow{\mathrm{id}_\C} \C \subseteq R \\
    t     & \xmapsto{\quad} f
\end{align*}
And also note that the locus of $f$ (as an \textbf{analytic set} in complex manifold $X^\mathrm{an}$) can be \textbf{non-smooth}.

We define
$$
    \Psi_{f,\lambda} = \lim_{\substack{\longrightarrow \\ m \to \infty}} i^{-1} Rj_\ast L_m^{\lambda}
$$
which is called the $\lambda\textbf{-nearby cycles}$ of $f$.
And $T$-action ($\pi_1((\mathbb{A}_{\mathbb{C}}^1 \setminus \{0\})^\mathrm{an})$-action) on $L_m^{\lambda}$ induces $T$-action on $\Psi_{f,\lambda}$.



\begin{remark}
    Firstly, note that since $\{f=0\}$ may not smooth, so $i^{-1}$ is the sheaf-theoretic pull-back between topological (ringed) spaces.
    Which is somewhat different from scheme case.

    Secondly, $j_{\ast}$ is the direct image functor,
    and it's a left exact functor,
    we write $Rj_\ast$ the right derived functor of $j_{\ast}$,
    which can be viewed as the complex of sheaves obtained by taking the \textbf{injective resolution} of $L_{m}^{\lambda}$ and applying $j_{\ast}$.
    Since $j_\ast$ maybe have nontrivial cohomology even in a non-$\D$-module-theoretic context, we choose to use a derived functor.
\end{remark}

The above definition is a topological definition, which is a bit differed from Deligne's. And total cycles is defined by combining all information from nearby cycles.


On the other hand, recall Example~\ref{Example 8.11}, we have
$$
    \mathcal{D}_{X}[s]f^{s} \subseteq R\left[\frac{1}{f}\right][s] \cdot f^{s}.
$$

For closed point $\alpha \in \mathbb{A}_{\mathbb{C}}^{1} \simeq \operatorname{Spec} \mathbb{C}[s]$, we use $\mathfrak{m}_{\alpha}$ to denote its maximal ideal. We have proved in Example \ref{Example 8.11}:

\begin{remark}
    We now justify the symbol $R_f[s]f^s$,
    since in Rmk.~\ref{remark:second_treatment_of_example_8.11} we only get the module $R_f[\partial_t]$,
    and $f^s$ represent the canonical section $1\otimes 1$.

    One should note that the notation $R_f[\partial_t]$ (or more precisely $\C[\partial_t]\otimes_\C R_f$) itself doesn't imply any module structure.
    If we want to make it an $R_f$ (left or right) module, we should induce the action from its left $\D_Y$-module structure.

    The Step 5 Block A of Rmk.~\ref{remark:second_treatment_of_example_8.11} gives us an answer.
    We have a projection
    \[
        p: Y = X\times_{\Spec \C}\mathbb{A}^1_{\C} \to X
    \]
    which makes $\Ox$ a subring of $\Oy$.
    So we get a left $R$ module structure by
    \[
        R=\Ox\xhookrightarrow{p^\sharp} \Oy \subseteq \D_Y.
    \]

    Since before we pushforward left $\D_X$-module $R_f$ along $\Gamma_f$,
    $f\in\D_X$ already acts bijectively (left and right) on $R_f$,
    which makes $t\in\D_Y$ act bijectively on $\Gamma_{f,+}(R_f) = \C[\partial_t] \otimes_{\C} R_f$ \cite[Lemma~2.1.5]{PopaDMBG2021}.
    Then by Step 5 Block A+D of Rmk.~\ref{remark:second_treatment_of_example_8.11},
    we discover $f\in R$ also act bijectively on $\Gamma_{f,+}(R_f)$.
    So we finally give a left $R_f$ module structure on $\C[\partial_t] \otimes_{\C} R_f$.

    Then we discover $s = -t\partial_t$ and $R$ commute in $\D_Y$, this is because we have
    \[
        [\partial_t, g] = \partial_t(g) = 0 \quad (\text{see  \cite[Prop~2.25(iii)]{mustata2023dmodules}})
    \]
    By universal property of polynomial ring, we have $R_f[s]$ embedded in $\D_Y$ (as a $\C$-subalgebra),
    if we map formal symbol $s$ to $-t\partial_t \in \D_Y$.

    \medskip

    Then we can act $R_f[s]$ onto $1\otimes 1 \in \Gamma_{f,+}(R_f)$, then we claim:
    \begin{enumerate}
        \item this action is faithful, so we get a free rank 1 $R_f[s]$ submodule in $\Gamma_{f,+}(R_f)$.
        \item This submodule is the whole of $\Gamma_{f,+}(R_f)$.
    \end{enumerate}
    We denote it as $R_f[s]f^s$.
    Then we can give the abstractly constructed $R_f[s]$ a left $\D_Y$ module structure by pulling back the one on $\Gamma_{f,+}(R_f)$.
\end{remark}

\bigskip

\begin{remark}\label{rmk:base_change_of_R_f[s]_f^s}
    We then try to define $R_f(s)f^s$ and $R_f[s]_{\frak{m}_\alpha}f^s$, the construction should satisfy:
    \begin{enumerate}
        \item From (the proof in) the corollary~\ref{cor:bernstein_beilinson} below, which suggests that this construction should be functorial.
        \item Be consistent with \cite[The $\D_{U_s}$ module $\mathcal{M}$ defined in Proof of Theorem 2.2.1.]{PopaDMBG2021}.
    \end{enumerate}

    This is equivalent to say that we have a ring ($\C$-algebra) map
    \[
        \D_Y \to \End_\C(R_f[s])
    \]
    And we have a composition
    \[
        \D_X\otimes_\C \C[s] \to \D_Y \to \End_\C(R_f[s])
    \]
    by assigning $s$ to $-t\partial_t\in\D_Y$

    \medskip

    We define $R_f(s) := R_f[s]\otimes_{\C[s]}\C(s)$ and $R_f[s]_{\frak{m}_\alpha} := R_f[s]\otimes_{\C[s]}\C[s]_{\frak{m}_\alpha}$.
    Before giving it a left $\D$ module structure, we discover that this construction is functorial (in the category of $\C[s]$-modules).

    \medskip

    Then we give it a left $\D$ module structure, this structure should also be functorial.
    \begin{enumerate}
        \item $R_f(s)$ should have a left $\D_{R(s)/\C(s)} = \D_{R/\C}\otimes_\C \C[s]\otimes_{\C[s]}\C(s)$ module structure, see Lem.~\ref{lem:base_change_of_D_modules}.
        \item $R_f[s]_{\frak{m}_\alpha}$ should have a left $\D_{R[s]_{\frak{m}_\alpha}/\C[s]_{\frak{m}_\alpha}} = \D_{R/\C}[s]\otimes_{\C[s]} \C[s]_{\frak{m}_\alpha}$ module structure, but one should note $\C[s]_{\frak{m}_\alpha}$ is not a field.
    \end{enumerate}
    Thanks to the construction in the proof in Lem.~\ref{lem:base_change_of_D_modules}, the base change of the \textbf{ring of \underline{relative} differential operator} is functorial.

    \medskip

    To give a (left) module structure is equivalent to give the ring map
    \begin{align*}
        \D_{R/\C}[s]\otimes_{\C[s]}\C(s) & \to \End_\C(R_f[s]\otimes_{\C[s]}\C(s))                    \\
        P \otimes a                      & \mapsto \left( u \otimes b \mapsto P(u) \otimes ab \right)
    \end{align*}
    and
    \begin{align*}
        \D_{R/\C}[s]\otimes_{\C[s]}\C[s]_{\frak{m}_\alpha} & \to \End_\C(R_f[s]\otimes_{\C[s]}\C[s]_{\frak{m}_\alpha})  \\
        P \otimes a                                        & \mapsto \left( u \otimes b \mapsto P(u) \otimes ab \right)
    \end{align*}
    From the two formulas above, one see this definition of ring action (on module) is functorial. Where $u\mapsto P(u)$ is just
    \[
        \D_{R/\C} \otimes_\C \C[s] = \D_X\otimes_\C \C[s] \to \D_Y \to \End_\C(R_f[s])
    \]
    mentioned above.
\end{remark}

\begin{theorem}
    There exists $b(s) \in \mathbb{C}[s]$ s.t.
    \[
        b(s) \quad \mathrm{ kills } \quad \frac{\mathcal{D}_{X}[s]f^{s}}{\mathcal{D}_{X}[s]f^{s+1}}
    \]
    see Thm.~\ref{Theorem 8.4} or Ex.~\ref{Example 8.11} or Rmk.~\ref{remark:second_treatment_of_example_8.11}.
\end{theorem}



\begin{corollary}[Bernstein-Beilinson]\label{cor:bernstein_beilinson}
    Let $\mathbb{C}(s)$ be the fractional field of $\mathbb{C}[s]$, and $\mathcal{D}_{X}(s) = \mathcal{D}_{X} \otimes_{\mathbb{C}} \mathbb{C}(s)$; for all $\alpha \in \mathbb{A}_{\mathbb{C}}^{1}$ closed point, we obtain that:
    \begin{enumerate}
        \item $\mathcal{D}_{X}(s) f^{s+k} \simeq \mathcal{D}_{X}(s) f^{s}, $ for all $  k \in \mathbb{Z}$.
        \item There exists $k_{0} \geqslant 0$ s.t.
              \begin{itemize}
                  \item $\mathcal{D}_{X}[s]_{\mathfrak{m}_{\alpha}} f^{s-k} = \mathcal{D}_{X}[s]_{\mathfrak{m}_{\alpha}} f^{s-k_{0}}, $ for all $ k > k_{0}$
                  \item $\mathcal{D}_{X}[s]_{\mathfrak{m}_{\alpha}} f^{s+k} = \mathcal{D}_{X}[s]_{\mathfrak{m}_{\alpha}} f^{s+k_{0}},$ for all $  k > k_{0}$
              \end{itemize}
    \end{enumerate}

\end{corollary}

\begin{remark}\label{rmk:D_X(s)f^s=R_f(s)f^s_and_D_X[s]_m_alpha_f^s=R_f[s]_m_alpha_f^s}
    Noted that here  $\mathcal{D}_{X}(s) f^{s} = R[\frac{1}{f}](s) f^{s},\ \mathcal{D}_{X}[s]_{\mathfrak{m}_{\alpha}} f^{s} = R[\frac{1}{f}][s]_{m_{\alpha}} f^{s}$.
\end{remark}

\begin{proof}
    \mbox{}

    $1)$ Since $\mathbb{C}(s)$ can be view as the localization at $(0)$ of $\mathbb{C}[s]$,
    we have $\mathbb{C}(s)$ is flat over $\mathbb{C}[s]$,
    thus
    \[
        \frac{\mathcal{D}_{X}[s]f^{s}}{\mathcal{D}_{X}[s]f^{s+1}}\otimes \mathbb{C}(s)=\frac{\mathcal{D}_{X}(s)f^{s}}{\mathcal{D}_{X}(s)f^{s+1}}
    \]
    Note that
    \begin{enumerate}
        \item To commute tensor product with quotient, the ring map $\mathbb{C}[s]\to \mathbb{C}(s)$ is not needed to be flat, all tensor product is right exact.
        \item To make the incuced $\D_X(s)$-structure the same on LHS and RHS, the $\D_X(s)$-structure should be functorial, which is required in Rmk.~\ref{rmk:base_change_of_R_f[s]_f^s}.
    \end{enumerate}

    From the $\D_X(s)$-structure mentioned in Rmk.~\ref{rmk:base_change_of_R_f[s]_f^s}, we have
    \[
        b(s+k) \quad\mathrm{kills} \quad\frac{\mathcal{D}_{X}(s)f^{s}}{\mathcal{D}_{X}(s)f^{s+1}} \implies \D_{R/\C}[s]\otimes_{\C[s]} \C(s) \ni b(s+k)\otimes 1  \quad\mathrm{kills } \quad\frac{\mathcal{D}_{X}[s]f^{s}}{\mathcal{D}_{X}[s]f^{s+1}} \otimes_{\C[s]} \C(s)
    \]
    And $b(s+k) \neq 0$ is invertible in $\mathbb{C}(s)$, so
    \[
        b(s+k) \quad\mathrm{kills} \quad\frac{\mathcal{D}_{X}(s)f^{s}}{\mathcal{D}_{X}(s)f^{s+1}} \implies \frac{\mathcal{D}_X(s) f^{s+k}}{\mathcal{D}_X(s) f^{s+k-1}} = 0.
    \]

    $2)$  Note that Bernstein-Sato polynomial $b(s)\in\D_X[s]$ has coefficients in $\C$, see Thm.~\ref{Theorem 8.4}.

    Since the number of roots of $b(s)$ is finite,
    there exists $k_0 \ge 0$, $b(\alpha \pm k) \neq 0 \in\mathbb{C}$ for all $k > k_0$,
    then $b(s \pm k) \notin \mathfrak{m}_\alpha$ for all $k > k_0$.

    So $b(s \pm k)$ is invertible in $\mathbb{C}[s]_{\mathfrak{m}_\alpha}$ for all $k > k_0$,
    similarly, we get
    \[
        \frac{\mathcal{D}_X[s]_{\mathfrak{m}_\alpha} f^{s \pm k}}{\mathcal{D}_X[s]_{\mathfrak{m}_\alpha} f^{s \pm k + 1}} = 0.
    \]
\end{proof}

\begin{definition}
    For all $\alpha \in \mathbb{C}$:
    \begin{enumerate}
        \item $\mathcal{D}_X[s]_{\mathfrak{m}_\alpha} f^{s+k},  k \gg 0$ is stabled, called the $\textbf{minimal extension}$ of $R[\frac{1}{f}][s]_{\mathfrak{m}_\alpha} f^s$.
        \item $\Psi_{f,\alpha} := \frac{\mathcal{D}_X[s]_{\mathfrak{m}_\alpha} f^{s-k}}{\mathcal{D}_X[s]_{\mathfrak{m}_\alpha} f^{s+k}}$ is called the $\alpha\textbf{-nearby cycles}$ of $R$ along $f$.

              \cite[Definition~3.3.2.]{PopaDMBG2021} defines that $\alpha$-nearby cycles is a $\mathcal{D}_X$-module, NOT a $\mathcal{D}_X[s]$-module.
        \item $\Lambda := \mathbb{Z} + \{\text{roots of } b_{f}(s)\}$.
    \end{enumerate}
\end{definition}

\begin{remark}
    For one wonder why here we \emph{define} the minimal extension,
    see \cite[Theorem~2.10]{Wu2020ComparisonNearbyCyclesBfunctions}, which use
    \cite[Def~3.4.1]{hotta2008d} to define the minimal extension and claim the equivalence of the two.

    In $\mathcal{D}$-module theory,
    an extension is usually about given a module over an open subset $U$,
    and extending it to the whole space $X$ to get a $\mathcal{D}_{X}$-module,
    and we shall pay attention on how it extends around $f=0$.
    Take $U=\mathrm{Spec}\, R_{f}$,
    if the module over $U=X\backslash(f=0)$ can be viewed as $R$-module,
    then it's an extension,
    and we call $R_{f}[s]_{\mathfrak{m}_{\alpha}}f^{s}$  $\textbf{maximal extension}$.
\end{remark}

\begin{proposition}
    \quad

    \begin{enumerate}
        \item $s - \alpha$ acts on $\Psi_{f,\alpha}$ nilpotently
        \item $\Psi_{f,\alpha} \cong \Psi_{f,\alpha+1}$ (Note: a $\D_X$-module homomorphism, not a $\D_X[s]$-module homomorphism)
        \item $\Psi_{f,\alpha} \neq 0$ if and only if $\alpha \in \Lambda$
        \item $\Psi_{f,\alpha}$ is a regular holonomic $D_X$-module.
    \end{enumerate}
\end{proposition}
\begin{proof}
    \mbox{}

    \begin{enumerate}
        \item The descending chain of submodules
              \[
                  \mathcal{D}_{X}[s]_{\mathfrak{m}_{\alpha}}f^{s+k_{0}}\subseteq\mathcal{D}_{X}[s]_{\mathfrak{m}_{\alpha}}f^{s+k_{0}-1}\subseteq\dots\subseteq \mathcal{D}_{X}[s]_{\mathfrak{m}_{\alpha}}f^{s-k_{0}}
              \]
              induces a finite filtration on the quotient module
              \[
                  \frac{\mathcal{D}_{X}[s]_{\mathfrak{m}_{\alpha}}f^{s-k_{0}}}{\mathcal{D}_{X}[s]_{\mathfrak{m}_{\alpha}}f^{s+k_{0}}},
              \]
              whose successive quotients are of the form
              \[
                  \frac{\mathcal{D}_{X}[s]_{\mathfrak{m}_{\alpha}}f^{s+i}}{\mathcal{D}_{X}[s]_{\mathfrak{m}_{\alpha}}f^{s+i+1}}.
              \]
              For each such quotient, there exists a polynomial $b_{i}(s)$ that annihilates it (before taking localization). Because $s-\alpha$ is the only non-invertible irreducible polynomial in the local ring $\mathbb{C}[s]_{\mathfrak{m}_{\alpha}}$, each $b_{i}(s)$ takes the form $u_{i}(s-\alpha)^{n_{i}}$ for some unit $u_{i} \in \mathbb{C}[s]_{\mathfrak{m}_{\alpha}}$.
              By taking $n = \sum n_{i}$, the element $(s-\alpha)^{n}$ annihilates the entire module $\frac{\mathcal{D}_{X}[s]_{\mathfrak{m}_{\alpha}}f^{s-k_{0}}}{\mathcal{D}_{X}[s]_{\mathfrak{m}_{\alpha}}f^{s+k_{0}}}$. Hence, $s-\alpha$ acts nilpotently on it.

        \item Recall that $t$ acts on $R_f[s]f^s$ bijectively. Consider the $\mathcal{D}_{X}$-linear automorphism of $\mathcal{D}_{X}(s)f^{s}$ defined by the shift
              \[
                  T_{-1}\big(P(s)f^{s}\big) = P(s-1)f^{s-1}.
              \]
              For any element $m$, we have
              \[
                  T_{-1}(s \cdot m) = (s-1) \cdot T_{-1}(m), \quad\text{and generally}\quad T_{-1}(P(s) \cdot m) = P(s-1) \cdot T_{-1}(m).
              \]
              This means $T_{-1}$ is not $\D_X[s]$-linear.

              When localizing at $\mathfrak{m}_{\alpha}$,
              a polynomial $g(s)$ becomes invertible if $g(\alpha) \neq 0$.
              Under the shift, the action of $g(s)$ becomes the action of $g(s-1)$, evaluated at $s = \alpha+1$ is also non-zero, making it invertible in $\mathbb{C}[s]_{\mathfrak{m}_{\alpha+1}}$.
              So we can safely define the map below:
              \begin{align*}
                  T_{-1}: \mathcal{D}_{X}[s]_{\mathfrak{m}_{\alpha}}f^{s-k} & \stackrel{\sim}{\longrightarrow} \mathcal{D}_{X}[s]_{\mathfrak{m}_{\alpha+1}}f^{s-k-1} \\
                  \frac{m}{g(s)}                                            & \longmapsto \frac{T_{-1}(m)}{g(s-1)}
              \end{align*}
              for any integer $k$. And one should verify that it is indeed an isomorphism.

              For $k \gg 0$, the sequence of modules stabilizes. Passing to the respective quotients, this induces the desired $\mathcal{D}_X$-module isomorphism $\Psi_{f,\alpha} \xrightarrow{\sim} \Psi_{f,\alpha+1}$.

        \item Since $\alpha \not\in\Lambda$, $\alpha$ is not a root of $b(s+k)$, then $b(s+k)$ is a unit for all $k$, thus $\Psi_{f,\alpha}=0$.

              If $\Psi_{f,\alpha}\neq 0$,
              there exists $i$ such that $\frac{\mathcal{D}_{X}[s]_{\mathfrak{m}_{\alpha}}f^{s+i}}{\mathcal{D}_{X}[s]_{\mathfrak{m}_{\alpha}}f^{s+i+1}}\neq 0$,
              and then $\alpha$ is a root of $b(s+i)$ such that $b(s+i)$ is not invertible in $\mathbb{C}[s]_{\mathfrak{m}_{\alpha}}$. Thus $\alpha\in\Lambda$.

        \item $M_{\alpha} = \mathbb{C}[t, \frac{1}{t}] \cdot t^{\alpha}$ (see Ex.~\ref{ex:first-examples-of-d-modules}) is regular
              (at least analytically at the stalk of $0$, see Rmk.~\ref{rmk:indecomposable_objects_in_RH_D} and Prop.~\ref{prop:classification-unipotent-indecomposable-regular-holonomic},
              to be verified algebraically it should be embedded into $\mathbb{P}^1_\C$, see Rmk.~\ref{rmk:algebraic_KK_regularity_holomonic})
              holonomic $\mathcal{D}_{\mathbb{A}^{1}_{\mathbb{C}}}$-module.

              Consider
              \begin{align*}
                  f: X   & \longrightarrow \mathbb{A}^{1}_{\mathbb{C}} \\
                  \alpha & \longmapsto t=f(\alpha)
              \end{align*}
              which is induced by
              \begin{align*}
                  \mathbb{C}[t] & \longrightarrow R \\
                  t             & \longmapsto f
              \end{align*}
              We admit the fact that the $\D$-module theoretical (derived) \textbf{inverse image} functor \cite[Sect.~1.5, It is the derived one that comptible with Thm.~6.1.5]{hotta2008d} preserves (algebraic we use here) regular holonomicity \cite[Thm.~6.1.5]{hotta2008d}, then
              \[
                  f^{*} M_{\alpha}:=M_{\alpha}\otimes_{\mathbb{C}[t]} R\cong R_{f}\cdot f^\alpha= R[\frac{1}{f}] f^{\alpha}
              \]
              makes $R_{f}\cdot f^\alpha$ a regular holonomic $\mathcal{D}_{X}$-module.

              With $(2)$, we can assume $\alpha\neq 0$, otherwise, we just shifts $\alpha$ to $\alpha+1$.
              We have a short exact sequence:
              \[
                  0 \longrightarrow \frac{R_f[s]_{\mathfrak{m}_{\alpha}} f^{s}}{(s-\alpha) R_f[s]_{\mathfrak{m}_{\alpha}} f^{s}} \xrightarrow{\;(s-\alpha)\cdot\;} \frac{R_f[s]_{\mathfrak{m}_{\alpha}} f^{s}}{(s-\alpha)^2 R_f[s]_{\mathfrak{m}_{\alpha}} f^{s}} \xrightarrow{\;\text{proj}\;} \frac{R_f[s]_{\mathfrak{m}_{\alpha}} f^{s}}{(s-\alpha) R_f[s]_{\mathfrak{m}_{\alpha}} f^{s}} \longrightarrow 0
              \]
              where the injective mapping is given by \textbf{left} multiplication by $(s-\alpha)$ (i.e., sending the residue class $\overline{m}$ to $\overline{(s-\alpha)m}$), and the surjective mapping is the natural canonical projection.

              Here, the isomorphism
              \[
                  \frac{R_f[s]_{\mathfrak{m}_{\alpha}} f^{s}}{(s-\alpha) R_f[s]_{\mathfrak{m}_{\alpha}} f^{s}} \cong R_{f} \cdot f^{\alpha}
              \]
              arises because quotienting out the $(s-\alpha)$-multiples algebraically corresponds to taking the evaluation map $s \mapsto \alpha$.
              \begin{enumerate}
                  \item The ring $\mathbb{C}[s]_{\mathfrak{m}_{\alpha}}$ is a local ring with (its unique) maximal ideal $(s-\alpha)$ reduces canonically to $\mathbb{C}$, collapsing the coefficients from $R_f[s]_{\mathfrak{m}_{\alpha}}$ down to $R_{f}$.
                  \item And the formal symbol $f^{s}$ becomes $f^{\alpha}$, this step didn't actually do evaulation, since they are \textbf{both} the canonical section $1\otimes 1$ in their respective presentations. So the real thing goes $1\otimes 1 \mapsto 1\otimes 1$.
              \end{enumerate}
              Thus, this quotient module is precisely the regular holonomic $\mathcal{D}_X$-module $R_{f} \cdot f^{\alpha}$.

              Since regular holonomic $\mathcal{D}_{X}$-modules form an abelian category \cite[Combine Prop.~3.1.2 and Def.~6.1.1]{hotta2008d}, we have $\frac{R_{f}[s]_{\mathfrak{m}_{\alpha}}f^{s}}{(s-\alpha)^{2}R_{f}[s]_{\mathfrak{m}_{\alpha}}f^{s}}$ is regular holonomic.

              And further, we obtain the generalized short exact sequence:
              \[
                  0 \longrightarrow \frac{R_f[s]_{\mathfrak{m}_{\alpha}} f^{s}}{(s-\alpha) R_f[s]_{\mathfrak{m}_{\alpha}} f^{s}} \xrightarrow{\;(s-\alpha)^N\cdot\;} \frac{R_f[s]_{\mathfrak{m}_{\alpha}} f^{s}}{(s-\alpha)^{N+1} R_f[s]_{\mathfrak{m}_{\alpha}} f^{s}} \xrightarrow{\;\text{proj}\;} \frac{R_f[s]_{\mathfrak{m}_{\alpha}} f^{s}}{(s-\alpha)^N R_f[s]_{\mathfrak{m}_{\alpha}} f^{s}} \longrightarrow 0
              \]
              By induction, $\frac{R_f[s]_{\mathfrak{m}_{\alpha}} f^{s}}{(s-\alpha)^{N} R_f[s]_{\mathfrak{m}_{\alpha}} f^{s}}$ is regular holonomic for all $N$.

              By $(1)$, the action of $(s-\alpha)$ on $\Psi_{f,\alpha} = \frac{\mathcal{D}_{X}[s]_{\mathfrak{m}_{\alpha}}f^{s-k}}{\mathcal{D}_{X}[s]_{\mathfrak{m}_{\alpha}}f^{s+k}}$ is nilpotent, meaning there exists some large enough integer $N \gg 0$ such that $(s-\alpha)^N$ entirely annihilates $\Psi_{f,\alpha}$.
              Because of this annihilation, the natural surjective $\mathcal{D}_X[s]_{\mathfrak{m}_{\alpha}}$-module morphism
              \[
                  \mathcal{D}_{X}[s]_{\mathfrak{m}_{\alpha}}f^{s-k} \xrightarrow{\;\text{proj}\;} \frac{\mathcal{D}_{X}[s]_{\mathfrak{m}_{\alpha}}f^{s-k}}{\mathcal{D}_{X}[s]_{\mathfrak{m}_{\alpha}}f^{s+k}} = \Psi_{f,\alpha}
              \]
              factors through the quotient by the submodule generated by $(s-\alpha)^N$. This explicitly means that the map descends to a well-defined surjective homomorphism:
              \[
                  \frac{\mathcal{D}_X[s]_{\mathfrak{m}_{\alpha}} f^{s-k}}{(s-\alpha)^N \mathcal{D}_X[s]_{\mathfrak{m}_{\alpha}} f^{s-k}} \twoheadrightarrow \Psi_{f,\alpha}.
              \]
              Since one can verify
              \begin{enumerate}
                  \item $\mathcal{D}_{X}[s]_{\mathfrak{m}_{\alpha}}f^{s-k} \cong \mathcal{D}_{X}[s]_{\mathfrak{m}_{\alpha}}f^{s}$,
                  \item Rmk.~\ref{rmk:D_X(s)f^s=R_f(s)f^s_and_D_X[s]_m_alpha_f^s=R_f[s]_m_alpha_f^s},
                        i.e.,
                        $\mathcal{D}_X[s]_{\mathfrak{m}_{\alpha}} f^{s}$ as a submodule $R_f[s]_{\mathfrak{m}_{\alpha}} f^{s}$ actually equals it.
              \end{enumerate}
              Thus, we can view $\Psi_{f,\alpha}$ as a quotient of the $\mathcal{D}_{X}$-module $\frac{R_f[s]_{\mathfrak{m}_{\alpha}} f^{s}}{(s-\alpha)^N R_f[s]_{\mathfrak{m}_{\alpha}} f^{s}}$.
              So $\Psi_{f,\alpha}$ itself must also be a regular holonomic $\mathcal{D}_{X}$-module.
    \end{enumerate}

\end{proof}


\section{Comparison theorem}
Our next goal is to prove the following theorem:
\begin{theorem}[Comparison] \label{Theorem 9.5}
    For all $\alpha \in \mathbb{C}$, $\lambda = e^{2\pi i\alpha}$, $$\mathrm{DR}(\Psi_{f,-\alpha}) \simeq i_* \Psi_{f,\lambda}[-1]$$such that $e^{2\pi is} = \lambda^{-1}\cdot e^{2\pi i(s+\alpha)}$ is the corresponding $T$-action.
\end{theorem}
\begin{remark}
    with $(s-\alpha)$ acts on $\Psi_{f,\alpha}$ nilpotent, $e^{2\pi i(s+\alpha)}$ will be a well-defined operator for $\Psi_{f,-\alpha}$.
\end{remark}


Define $N_m^{\alpha,k} = \bigoplus_{\ell=1}^m \mathcal{D}_X[s] f^{s-k} \cdot e_\ell^{\alpha}$, know that $s=-t\partial_{t}$,  $s \cdot (\eta \cdot e_\ell^{\alpha}) = (s\eta) \cdot e_\ell^{\alpha}+\eta(-\alpha e^{\alpha}_{\ell}-e^{\alpha}_{\ell-1})=(s-\alpha)\eta \cdot e^{\alpha}_{\ell}-\eta \cdot e^{\alpha}_{\ell-1}$, so $N_m^{\alpha,k}$ is also an $\mathcal{D}_X[s]$-mod.

We define $b_f(s)$ to be the monic polynomial of least degree among $b(s)$ such that $$b(s) \cdot \frac{\mathcal{D}_X[s] f^s}{\mathcal{D}_X[s] f^{s+1}} = 0$$called the $\textbf{b-function} \text{ or } \textbf{Bernstein-Sato polynomial of }f$.

\begin{lemma}\label{lemma 9.6}
    \quad

    \begin{enumerate}
        \item $\mathrm{DR}\left(R_f[s] f^s \otimes^L_{\mathbb{C}(s)} \mathbb{C}_\alpha\right) \simeq \mathrm{DR}\left(R_f[s] f^{s+\alpha} \otimes^L_{\mathbb{C}(s)} \mathbb{C}_\alpha\right)\simeq R j_* (f^{-1} L_1^\lambda)$
        \item $\mathrm{DR}\left(\mathcal{D}_X[s] f^{s+k} \otimes^{L}_{\mathbb{C}(s)} \mathbb{C}_\alpha\right) \simeq \mathrm{DR}\left(\mathcal{D}_X[s] f^{s+k+\alpha} \otimes^{L}_{\mathbb{C}(s)} \mathbb{C}_\alpha\right)\simeq j_! (f^{-1} L_1^\lambda), \quad k \gg |\alpha|$
    \end{enumerate}
    where $\mathbb{C}_\alpha$ is the residue field of $\alpha \in \mathbb{A}_{\mathbb{C}}^1$, and $\lambda = e^{2\pi i \alpha}$.
\end{lemma}
\begin{proof}
    Take $0\rightarrow \mathbb{C}[s]\xrightarrow{s-\alpha} \mathbb{C}[\alpha]\rightarrow \frac{\mathbb{C}[s]}{s-\alpha}=\mathbb{C}_{\alpha}$ free resolution, then we obtain that $R_f[s]f^s \otimes^L_{\mathbb{C}(s)} \mathbb{C}_\alpha$ $$\stackrel{\text{q.i}}{\simeq}[ R_f[s]f^s \xrightarrow{s-\alpha} R_f[s]f^s ]\stackrel{\text{q.i}}{\simeq}[R_{f}[s]f^{s+\alpha}\xrightarrow{s} R_f[s]f^{s+\alpha}]\simeq R_f f^\alpha.$$with $U\xrightarrow{f}\mathbb{A}^{1}_{\mathbb{C}}\backslash\{0\}$, $\mathrm{DR}(R_f \cdot f^\alpha)|_U \simeq f^{-1} L_1^\lambda$, using log resolution of $(X,f=0)$, one can prove $\mathrm{DR}(R_f \cdot f^\alpha) \simeq R j_*(f^{-1} L_1^\lambda)$.
    Using resolution of singularities and duality, one can prove 2.

\end{proof}
\begin{lemma} \label{lemma 9.7}
    $$(b_f(s-k-\alpha))^m \text{ kills } N_m^{\alpha,k}/N_m^{\alpha,k-1}$$
\end{lemma}
\begin{proof}
    By construction, we have a SES: $$0\longrightarrow N_{m-1}^{\alpha,k} \longrightarrow N_m^{\alpha,k} \longrightarrow Q \longrightarrow 0$$with $Q=\mathcal{D}_{X}[s]f^{s-k}\bar{e}_{m-1}^{\alpha}$, with $s\cdot \bar{e}^\alpha_{m-1}=(s-\alpha)\overline{e}^\alpha_{m-1}$, we choose $s-\alpha$ to replace $s$, then  $Q\simeq \mathcal{D}_{X}[s]f^{s-k-\alpha}$.

    Since $b_{f}(s-k-\alpha)$ kills $\mathcal{D}_{X}[s]f^{s-k-\alpha}$, and by induction, we have $(b_f(s-k-\alpha))^{m -1}$ kill $N_{m-1}^{\alpha,k}$, done.
\end{proof}
\begin{proposition}
    There exists $k_0 > 0$ such that for all $k \geq k_0$, we obtain that:
    \begin{enumerate}
        \item $\mathrm{DR}(N_m^{-\alpha,k} \xrightarrow{s} N_{m}^{-\alpha,k}) \simeq R j_*\left(f^{-1} L_m^{\lambda}\right)$;
        \item $\mathrm{DR}(N_m^{-\alpha,-k} \xrightarrow{s} N_{m}^{-\alpha,-k})\simeq j_!\left(f^{-1} L_m^{\lambda}\right)$.
    \end{enumerate}
\end{proposition}
\begin{proof}
    Let $\underline{\mathfrak{m}}_0$ be the maximal ideal of $0 \in \mathbb{A}_{\mathbb{C}}^1$, then $$[N_{m}^{-\alpha,k} \xrightarrow{s} N_{m}^{-\alpha,k}]\stackrel{\text{q.i}}{\simeq}N^{-\alpha,k} \otimes^{L}_{\mathbb{C}[s]} \mathbb{C}_{0} \stackrel{\text{q.i}}{\simeq} [N^{-\alpha,k}_{m,\underline{\mathfrak{m}}_{0}} \xrightarrow{s}N^{-\alpha,k}_{m,\underline{\mathfrak{m}}_{0}}]\text{ localization}\atop \stackrel{\text{q.i}}{\simeq}[\bigoplus \limits_{\ell=1}^{m-1} R_{f}[s]f^{s} \cdot e_{\ell}^{-\alpha} \xrightarrow{s} \bigoplus \limits_{\ell=1}^{m-1}R_{f}[s]f^{s} \cdot e_{\ell}^{-\alpha}]\text{ by Lemma \ref{lemma 9.7}}.$$Using Lemma \ref{lemma 9.6} (1) for $\mathrm{DR}(R_{f}[s]f^{s+\alpha}\otimes_{\mathbb{C}[s]}\mathbb{C}_{0})$, by induction, Part (1) is proved. And part (2) is similar.
\end{proof}

\begin{lemma}
    Let $W$ be a $\mathbb{C}$-vector space, and let $\varphi$ be a $\mathbb{C}$-linear operator on $W$ admitting a minimal polynomial.
    Set $W_{\infty} = \bigoplus_{k=0}^{\infty} W \otimes e_k$ and define $$\varphi_{\infty}(w \otimes e_i) = (\varphi - \alpha) \cdot w \otimes e_i - w \otimes e_{i-1} \quad \text{for all }  w \in W, \; (e_{-1} = 0).$$Then $\varphi_{\infty}$ is surjective and $\ker(\varphi_{\infty}) = W_{\alpha}$ is the generalized $\alpha$-eigenspace.
\end{lemma}
\begin{proof}
    Define a map
    \begin{equation}
        W_{\alpha} \longrightarrow \ker(\varphi_{\infty})\atop w \longmapsto \sum\limits_{i \geqslant 0} (\varphi - \alpha)^i w \otimes e_{i}.
        \tag{$\ast$}
        \label{eq}
    \end{equation}
    This map is an isomorphism.

    Then we prove surjectivity: If $w\in W_{\alpha}$, then $\varphi_{\infty}( -\sum\limits_{i>j} (\varphi-\alpha)^{i-j-1}w\otimes e_{i})=w\otimes e_{j}$; If $w\in W_{\alpha}^{\perp}$, then $\varphi_{\infty}( \sum\limits_{i=1}^{j}(\varphi-\alpha)^{-i}w\otimes e_{j-i+1} )=w\otimes e_{j}$.
\end{proof}
\begin{corollary}
    For $k \gg 0$, $$
        \Psi_{f,\alpha}[1] \stackrel{\text{q.i}}{\simeq}\varinjlim_m \left( \frac{N_m^{\alpha,k}}{N_m^{\alpha,-k}} \xrightarrow{s} \frac{N_m^{\alpha,k}}{N_m^{\alpha,-k}} \right).
    $$
\end{corollary}

\begin{proof}[Proof of  (\ref{Theorem 9.5})]
    $$
        \mathrm{DR}(\Psi_{f,-\alpha}[1]) \cong \mathrm{DR}\left( \varinjlim_m  \left( \frac{N_m^{-\alpha,k}}{N_m^{-\alpha,-k}} \xrightarrow{s} \frac{N_m^{-\alpha,k}}{N_m^{-\alpha,-k}} \right) \right)$$
    $$ \cong \varinjlim_m  \mathrm{DR}\left( \frac{N_m^{-\alpha,k}}{N_m^{-\alpha,-k}} \xrightarrow{s} \frac{N_m^{-\alpha,k}}{N_m^{-\alpha,-k}} \right)$$
    $$\simeq \varinjlim_m\ \text{cone}\left( DR\left(N_{m}^{-\alpha,-k} \xrightarrow{s} N_m^{-\alpha,-k}\right) \rightarrow DR\left(N_m^{-\alpha,k} \xrightarrow{s} N_m^{-\alpha,k}\right)\right)$$
    $$\simeq\varinjlim_m\ \text{cone}\left(j_! f^{-1} L_m^{\lambda} \rightarrow R j_* f^{-1} L_m^{\lambda}\right)\simeq \Psi_{f,\lambda}$$

    By $\eqref{eq}$, we consider $$   \Psi_{f,-\alpha} \longrightarrow \ker(-s_{\infty})\atop \eta \longmapsto \sum\limits_{i \geqslant 0} (-s-\alpha)^i \eta\otimes  e_{i}.$$

    Then $s+\alpha = -J_{0,\infty}$, and $e^{2\pi i(s+\alpha)} = e^{-2\pi i J_{0,\infty}}$. Furthermore, we obtain that $$e^{-2\pi i\alpha} \cdot e^{2\pi i(s+\alpha)} = e^{-2\pi i J_{\alpha,\infty}} = T$$
\end{proof}


\chapter{Homework}
\section{Homework: week 1 (September 19, 2025)}
\documentclass[12pt]{article}
\usepackage[letterpaper,margin=1.01in]{geometry}
\usepackage{amsmath,amssymb,amsfonts}
\usepackage{mathrsfs}
\usepackage{microtype}

\title{Homework: week 1}
\date{September 19, 2025}

\begin{document}
\maketitle

$R$: a commutative ring with a unit.

\bigskip
\noindent
\textbf{Problem 1.} Let $M$ be a $R$-module. Let $T^n(M) = M \otimes_R M \otimes_R \cdots \otimes_R M$ be the $n$-th tensor product. For $n = 0$, we put $T^0(M) = R$. Then
\[
T(M) = \bigoplus_{n\ge 0} T^n(M)
\]
is a (noncommutative) $R$-algebra, called the \emph{tensor algebra} of $M$. We define the \emph{symmetric
algebra}
\[
\operatorname{Sym}^\bullet(M) = \bigoplus_{n\ge 0} \operatorname{Sym}^n(M)
\]
to be the quotient quotient of the two-sided ideal generated by all expressions $x \otimes y - y \otimes x$,
for all $x,y \in M$; $\operatorname{Sym}^n(M)$ is called the $n$-th \emph{symmetric power}.

We also define the \emph{exterior algebra}
\[
\bigwedge(M) = \bigoplus_{n\ge 0} \bigwedge^n(M)
\]
to be the quotient of $T(M)$ by the two-sided ideal generated by all expressions $x \otimes x$ for
$x \in M$; $\bigwedge^n(M)$ is called the $n$-th \emph{Wedge power} of $M$. Now we assume that $F$ is a free
$R$-module of rank $r$. Then prove:

\begin{enumerate}
  \item $\bigwedge(M)$ is a skew commutative graded $R$-algebra, which means that if $u \in \bigwedge^k(M)$ and
  $v \in \bigwedge^l(M)$, then $u \wedge v = (-1)^{kl} v \wedge u$.
  \item $\operatorname{Sym}^k(F)$ and $\bigwedge^k(F)$ are also free. Find their ranks as $R$-modules.
  \item $\operatorname{Sym}^\bullet(F) \simeq R[\xi_1,\ldots,\xi_r]$.
  \item For every $k$, the multiplication map $\bigwedge^k(F) \otimes \bigwedge^{r-k}(F) \to \bigwedge^r(F)$ induces an isomorphism
  \[
  \bigwedge^k(F) \simeq (\bigwedge^{r-k}(F))^* \otimes \bigwedge^r(F),
  \]
  where $(\bigwedge^{r-k}(F))^* = \operatorname{Hom}_R(\bigwedge^{r-k}(F),R)$ is its dual module.
\end{enumerate}

\begin{enumerate}
  \setcounter{enumi}{4}
  \item Let $0 \to F' \to F \to F'' \to 0$ be a short exact sequences of free $R$-modules. Then for
  every $k$, there exists a decreasing finite filtration
  \[
  \operatorname{Sym}^k(F) = G^0 \supseteq G^1 \supseteq \cdots \supseteq G^k \supseteq G^{k+1} = 0
  \]
  such that $G^p/G^{p+1} \simeq \operatorname{Sym}^p(F') \otimes_R \operatorname{Sym}^{k-p}(F'')$.
  \item For every $k$, there exists a decreasing finite filtration
  \[
  \bigwedge^n(F) = G^0 \supseteq G^1 \supseteq \cdots \supseteq G^k \supseteq G^{k+1} = 0
  \]
  such that $G^p/G^{p+1} \simeq \bigwedge^p(F') \otimes_R \bigwedge^{k-p}(F'')$. In particular,
  \[
  \bigwedge^r(F) = \bigwedge^{r'}(F') \otimes_R \bigwedge^{r''}(F''),
  \]
  where $r'$ and $r''$ are ranks of $F'$ and $F''$ respectively.
\end{enumerate}

$A$: a ring with unit (possibly not commutative).

\bigskip
\noindent
\textbf{Problem 2.} Let $\Gamma_\bullet A$ be a $\mathbb{Z}_{\ge 0}$-filtration of $A$. Prove that if $\operatorname{gr}_\bullet A$
is noetherian, then so is $A$ itself.

\bigskip
Let $\Omega_\bullet A$ be a $\mathbb{Z}$-filtration of $A$. We define a topology on $A$ by declaring a subset to be
open precisely if it is a union of cosets $a + \Omega_i A$ for $a \in A$ and $i \in \mathbb{Z}$, called the filtration
topology of $A$.

\bigskip
\noindent
\textbf{Problem 3.} Prove that the filtration topology of $A$ is Hausdorff.

Hint: use the condition $\bigcap_i \Omega_i M = 0$.

\bigskip
A \emph{Cauchy sequence} in $A$ is a sequence $(m_n)_{n\ge 0}$ such that for every open set $U$ containing
$0$, there exists $c=c(U) \in \mathbb{N}$ such that if $i,j>c$ then $m_i-m_j \in U$. We say two Cauchy
sequences $(m_n)$ and $(m_n')$ are equivalent if for every open set $U$ containing $0$, there exists
$c=c(U) \in \mathbb{N}$ such that if $n>c$ then $m_n-m_n' \in U$.

Let $\widehat{A}$ be the set of all all equivalence classes of Cauchy sequences.

\bigskip
\noindent
\textbf{Problem 4.} Prove that $\widehat{A}$ is a ring and that we have a ring homomorphism
\[
A \to \widehat{A}, \quad a \mapsto [(a)],
\]
where $(a) = (a,a,\ldots,a,\ldots)$ is the constant sequence.

\bigskip
$A$ is called complete with respect to $\Omega_\bullet A$ if $A \simeq \widehat{A}$.

\bigskip
\noindent
\textbf{Problem 5.} Prove that if $A$ is complete and if $\operatorname{gr}_\bullet A$ is noetherian, then $A$ is also noetherian.

\bigskip
\noindent
\textbf{Problem 6.} Let $X$ be a smooth affine algebraic variety over $\mathbb{C}$. Prove
\[
\operatorname{gr}^F_\bullet D_X \simeq \operatorname{Sym}^\bullet T_X
\]
where $F_\bullet D_X$ is the order filtration and $T_X$ is the module of $\mathbb{C}$-derivations on $X$.
\end{document}

\newpage
\section{Homework: week 2 (October 7, 2025)}
\documentclass[12pt]{article}
\usepackage[letterpaper,margin=1.01in]{geometry}
\usepackage{amsmath,amssymb,amsfonts}
\usepackage{mathrsfs}
\usepackage{microtype}

\title{Homework: week 2}
\date{October 7, 2025}

\begin{document}
\maketitle

\bigskip
$A$: a ring with unit (possibly not commutative) with a $\mathbb{Z}_{\ge 0}$-filtration $\Gamma_\bullet A$.

$X=\operatorname{Spec}R$: a smooth affine algebraic variety over $\mathbb{C}$ of finite type

\bigskip
\noindent
\textbf{Problem 1.} Let $M$ be a finitely generated $D_X$-module. Prove that if $F_\bullet M$ is a good
filtration, then $F_p M$ is a finitely generated $R$-module for every $p$.

\bigskip
\noindent
Hint: prove $F_p D_X$ finitely generated over $R$ first.

\bigskip
\noindent
\textbf{Problem 2.} Let $M$ be a $D_X$-module with a compatible filtration $\Omega_\bullet M$ such that $\Omega_p M$ is
finite generated over $R$ for every $p$. Then $\Omega_\bullet M$ is called a ``good'' filtration if there exists
$n_0 \gg 0$ such that
\[
F_j D_X\cdot \Omega_i M = \Omega_{i+j} M \quad \forall j \ge 0,\, i \ge n_0.
\]
Prove that $\Omega_\bullet M$ is ``good'' iff $\operatorname{gr}^\Omega_\bullet M$ is finitely generated over $\operatorname{gr}^F_\bullet D_X$,
i.e. ``good'' filtrations $\Longleftrightarrow$ good filtrations we defined previously.

\bigskip
\noindent
\textbf{Problem 3.} For $x \in X$ a closed point, we use $m_x$ to denote its maximal ideal. Prove that
we have canonical isomorphisms
\[
(m_x/m_x^2)^* := \operatorname{Hom}(m_x/m_x^2,\mathbb{C}) \simeq T_{X,x} := T_X \otimes_R R/m.
\]
and
\[
 m_x/m_x^2 \simeq \Omega_{X,x} := \Omega_X \otimes_R R/m.
\]

\bigskip
\noindent
Hint: use the natural map $d:R\to \Omega_X$ to induce a map $(m_x/m_x^2)^* \to T_{X,x} = \operatorname{Hom}_R(\Omega_X,R/m_x);$
then prove both of them are surjective.

\noindent
$(m_x/m_x^2)^*$ is called the \emph{zariski tangent space} of $X$ at $x$.

\bigskip
\noindent
\textbf{Problem 4.} Use problem 3 to prove that if $m \subset R$ is a maximal ideal, then we can find
$x_1,\ldots,x_n \in m$ such that $m = (x_1,\ldots,x_n)$ and $(x_1,\ldots,x_n)$ is an algebraic coordinate system
on $X$.

\bigskip
\noindent
Hint: use Nakayama's Lemma.

Assume that $\Gamma_\bullet A$ is a almost commutative ring, i.e. $\operatorname{gr}^\Gamma_\bullet A$ is a commutative ring. For
$x \in \Gamma_k A$, $\sigma_k(x)$ is the image $x$ in $\Gamma_k A/\Gamma_{k-1}A$. For $x \in \Gamma_k A$, $y \in \Gamma_l A$, since $\operatorname{gr}_\bullet A$ is commutative,
we define
\[
\{\sigma_k(x),\sigma_l(y)\} = \sigma_{k+l-1}(xy - yx).
\]

\bigskip
\noindent
\textbf{Problem 5.} Prove that $(\operatorname{gr}_\bullet A,\{\cdot,\cdot\})$ is a well-defined Poisson algebra.

\bigskip
Meanwhile, we have its Rees ring $R_\bullet(A) = \bigoplus_i \Gamma_i A\cdot u^i$. Define $T := \frac{R_\bullet(A)}{u^2\cdot R_\bullet(A)}$.

\bigskip
\noindent
\textbf{Problem 6.} Prove that $T/u\cdot T \xrightarrow{u} u\cdot T$ is an isomorphism.

\bigskip
\noindent
We have another Poisson bracket on $\overline{T} = T/u\cdot T \simeq \operatorname{gr}^\Gamma_\bullet A$ defined by for $a,b\in T$
\[
\{\overline{a},\overline{b}\} = \overline{c}
\]
such that $\overline{c} \in \overline{T}$ is given by $ab - ba \in u\cdot T$ by using the isomorphism in Problem 6.

\bigskip
\noindent
\textbf{Problem 7.} Prove that the above two Poisson structure on $\operatorname{gr}_\bullet A$ are the same.

\bigskip
\noindent
\textbf{Problem 8.} If $M$ is an integrable connection on $X$, then prove that $Ch(M)=T_X^*X$.

\bigskip
\noindent
Hint: First, prove that there exists a finite open covering $X=\bigcup_i U_i$ with $U_i=\operatorname{Spec}R_{f_i}$ for
some $f_i \in R$ such that $M_{f_i}$ is free over $R_{f_i}$.
\end{document}

\newpage
\section{Homework: week 3 (September 30, 2025)}
\input{HwWeek3.tex}
\newpage
\section{Homework: week 4 and 5 (October 16, 2025)}
\documentclass[12pt]{article}
\usepackage[letterpaper,margin=1.01in]{geometry}
\usepackage{amsmath,amssymb,amsfonts}
\usepackage{mathrsfs}
\usepackage{microtype}

\title{Homework: week 4 and 5}
\date{October 16, 2025}

\begin{document}
\maketitle

$X=\operatorname{Spec}R$: a smooth affine algebraic variety over $\mathbb{C}$ of finite type with algebraic coor-\\
dinates $(x_1,\ldots,x_n)$.

\bigskip
\noindent
\textbf{Problem 1.} Let $M$ be an integrable connection on $X$. Then prove
\[
\mathrm{DR}(M) \stackrel{q.i.}{\simeq} H^0(\mathrm{DR}(M))
\]
and that $H^0(\mathrm{DR}(M))$ is a local system on $X^{\mathrm{an}}$, where $q.i.$ denotes quasi-isomorphism.

\bigskip
Let $i:X\hookrightarrow Y$ be a closed embedding such that $Y$ is also a smooth affine algebraic
variety over $\mathbb{C}$ of finite type with algebraic coordinates $(x_1,\ldots,x_n,y_1,\ldots,y_l)$ and that $X$ is
given by the ideal $(y_1,\ldots,y_l)$.

Let $N$ be a finitely generated $D_Y$-module. Then we define
\[
 i^\sharp N = H^0(K(N;y_1,\ldots,y_l)),
\]
where $K(N;y_1,\ldots,y_l)$ is the Koszul complex of $N$ with respect to the operations $N \xrightarrow{y_i} N$
for all $i=1,\ldots,l$.

\bigskip
\noindent
\textbf{Problem 2.} If $N$ is supported on $N$, then prove that $H^i(K(N;y_1,\ldots,y_l))=0$ for every
$i\ne 0$.

\bigskip
\noindent
\textbf{Problem 3.} If $N$ is supported on $N$, then prove that $\mathrm{Ch}(N)=T_X^*Y$ iff $i^\sharp N$ is an integrable
connection on $X$.
\end{document}

\newpage
\section{Homework: week 6 (October 23, 2025)}
\documentclass[12pt]{article}
\usepackage[letterpaper,margin=1.01in]{geometry}
\usepackage{amsmath,amssymb,amsfonts}
\usepackage{mathrsfs}
\usepackage{microtype}

\title{Homework: week 6}
\date{October 23, 2025}

\begin{document}
\maketitle

Let $F^\bullet$ be a complex with a $\mathbb{Z}$-indexed filtration as we defined in our lecture and let\\
$\{E^\bullet_\bullet(k)\}_k$ be the associated spectral sequence.

\bigskip
\noindent
\textbf{Problem 1.} Prove that
\[
E^i_\bullet(k+1)=H^i(E^\bullet_\bullet(k))
\]
for every $k\ge 0$.

\bigskip
$A$: a $\mathbb{C}$-algebra (not necessarily commutative)\\
$\Gamma_\bullet A$: a $\mathbb{Z}$-filtration on $A$\\
$M$: a finitely generated $A$-module

\bigskip
\noindent
\textbf{Problem 2.} Let $\Omega_\bullet M$ be a good filtration over $\Gamma_\bullet A$. Then prove that
\[
\mathrm{Rhom}_{R_\bullet^{\Gamma}A}\left(R_\bullet^{\Omega}M,R_\bullet^{\Gamma}A\right)\otimes \mathbb{C}[u]/(u\cdot\mathbb{C}[u])
\]
is quasi-isomorphic to
\[
\mathrm{Rhom}_{\operatorname{gr}_\bullet^{\Gamma}A}(\operatorname{gr}_\bullet^{\Omega}M,\operatorname{gr}_\bullet^{\Gamma}A).
\]
$X=\operatorname{Spec}R$: a smooth affine variety with coordinate $(x_1,\ldots,x_n)$

\bigskip
\noindent
\textbf{Problem 3.} Let $M$ be a finitely generated $D_X$-module. Prove that
\[
 j(M)+\dim(\mathrm{Ch}(M))=2n,
\]
where $j(M)$ is the grade number of $M$ over $\mathcal{D}_X$.
\end{document}

\newpage
\section{Homework: week 8 (November 6, 2025)}
\documentclass[12pt]{article}
\usepackage[letterpaper,margin=1.01in]{geometry}
\usepackage{amsmath,amssymb,amsfonts}
\usepackage{mathrsfs}
\usepackage{microtype}

\title{Homework: week 6}
\date{November 6, 2025}

\begin{document}
\maketitle

$X=\operatorname{Spec}R$: a smooth affine variety with coordinate $(x_1,\ldots,x_n)$

\bigskip
\noindent
\textbf{Problem 1.} Prove that the global dimension of $\mathcal{D}_X$ is $2n$.

\bigskip
\noindent
\textbf{Problem 2.} Calculate the $V$-filtrations of $\mathcal{M}_\alpha=\mathbb{C}[t,1/t]\cdot t^\alpha$ and
$\operatorname{Ext}^1_{\mathcal{D}_X}(\mathcal{M}_{-1},\mathcal{D}_X)$ along $(t=0)$ for $\alpha\in \mathbb{C}$.

\bigskip
Let $Y\subseteq X$ be the subvariety defined by $x_1=0$.

\bigskip
\noindent
\textbf{Problem 3.} Calculate the $V$-filtration of the $\mathcal{D}_X$-module $N=\overline{R}[\partial_{x_1}]$ along $Y$, where
$\overline{R}=R/(x_1)$.

\bigskip
\noindent
Hint: use Kashiwara's equivalence.

\bigskip
\noindent
\textbf{Problem 4.} Let $W$ be a $\mathbb{C}$-vector space and let $\varphi$ be a $\mathbb{C}$-linear operator on $W$ admitting
a minimal polynomial. Set
\[
W_\infty=\bigoplus_{k=0}^{\infty} W\otimes e_k
\]
and for some $\alpha\in \mathbb{C}$ define
\[
\varphi_\infty(w\otimes e_i)=(\varphi-\alpha)\cdot w\otimes e_i-w\otimes e_{i-1},\quad \forall w\in W
\]
with $e_{-1}=0$. Prove that $W_\alpha$ is isomorphic to $\ker(\varphi_\infty)$, where $W_\alpha$ is the generalized
$\alpha$-eigenspace with respect to $\varphi$-operation.
\end{document}


\cleardoublepage
\phantomsection
\addcontentsline{toc}{chapter}{Bibliography}
\bibliographystyle{alphaurl}
\bibliography{ref}

\end{document}
